\documentclass[11pt]{article}
% What the occupation functional of a finite Markov chain determines.
% arXiv version. It inputs the section files of the journal version problem131_reconstruction.tex
% and keeps the reversible design (sec_reversible.tex) as its final section.
% Build: pdflatex problem131_reconstruction_arxiv.tex, three times.

\usepackage[margin=1in]{geometry}
\usepackage{amsmath,amssymb,amsthm}
\usepackage{microtype}
\usepackage{xcolor}
\usepackage[shortlabels]{enumitem}
\usepackage[colorlinks=true,linkcolor=blue!45!black,citecolor=blue!45!black,urlcolor=blue!45!black]{hyperref}
\hypersetup{hypertexnames=false,pdftitle={What the occupation functional of a finite Markov chain determines},pdfauthor={Arjun Pemmasani}}

% ---------- theorem environments ----------
\newtheorem{theorem}{Theorem}[section]
\newtheorem{lemma}[theorem]{Lemma}
\newtheorem{proposition}[theorem]{Proposition}
\newtheorem{corollary}[theorem]{Corollary}
\newtheorem{introtheorem}{Theorem}
\renewcommand{\theintrotheorem}{\Alph{introtheorem}}
\theoremstyle{definition}
\newtheorem{definition}[theorem]{Definition}
\newtheorem{example}[theorem]{Example}
\theoremstyle{remark}
\newtheorem{remark}[theorem]{Remark}

% ---------- notation macros ----------
\newcommand{\R}{\mathbb{R}}
\newcommand{\E}{\mathbb{E}}
\newcommand{\Prob}{\mathbb{P}}
\newcommand{\supp}{\operatorname{supp}}
\newcommand{\ones}{\mathbf{1}}

\emergencystretch=1em

% Build switch. This file is the arXiv version, with the reversible design.
\newif\ifarxivbuild
\arxivbuildtrue


\title{What the occupation functional of a finite Markov chain determines:\\ principal minors, modal supports and cycle trees}
\author{Arjun Pemmasani\\[2pt]
\small Department of Mathematics, Harvey Mudd College\\
\small Claremont, CA 91711, USA\\
\small \texttt{apemmasani@hmc.edu}}
\date{}

\begin{document}
\maketitle

\begin{abstract}
What does the occupation measure of a finite Markov chain determine about its generator? For an
irreducible generator $Q$ on $d$ states, the Donsker and Varadhan functional determines exactly
the principal minors of $-Q$. Kim and Landim (Stochastic Process.\ Appl.\ 189 (2025) 104733)
showed that it determines a reversible generator and not a nonreversible one. The generators with
these minors form a finite set, which is $Q$ and its time reversal $Q^*$ when $-Q$ has no rank-one
cut. The exact law of the occupation times at one time determines the law at every time, from any
start. From a generic start it determines $Q$. From the stationary start it determines $Q$ up to
time reversal when $d\le4$, when $Q$ has no cut, and for generic $Q$, and then the second moments
suffice. This fails in general. Reversing $Q$ inside a block that every state leaves at one
common rate, and that every jump enters in proportion to the stationary law, keeps the stationary
occupation law, and on 5 states it gives generators other than $Q$ and $Q^*$. The first-order
support of the most likely counts of a rarely switching chain under killing at known rates
determines the same minors. For an Arrhenius family with positive
rates, the functionals determine the depth of every face and the Freidlin and Wentzell cycle tree,
without reversibility, and a threshold design fixed in advance reads the tree.\ifarxivbuild\ A
final section reads a reversible generator from a minimal design.\fi

\end{abstract}

\medskip
\noindent\textbf{Keywords:} identifiability; principal minors; occupation measure; large
deviations; Perron root; metastability; Freidlin--Wentzell cycles.

\noindent\textbf{Mathematics Subject Classification (2020):} 60J27, 60F10, 60J10, 15A15, 15A18,
15A29, 15B48.

\section{Introduction}\label{sec:intro}

Let $Q$ be an irreducible generator on $[d]=\{1,\dots,d\}$, and let $J_V$ be the
Donsker and Varadhan functional of its occupation measure~\cite{dv1975}. We prove
that $J_V$ determines exactly the principal minors of $-Q$
(Theorem~\ref{thm:A}). The generators with these minors form a finite set. It
consists of $Q$ and its time reversal $Q^*$ when $-Q$ has no rank-one cut, and it is
an orbit of cut-transposes in general. Kim and Landim~\cite[Section~5]{kl2024} show
that $J_V$ determines the minors of orders 1 and 2, and they observe that the
uniform 3-cycle and its reversal share $J_V$. Theorem~\ref{thm:A} says exactly
which generators share it.

The functional enters through the occupation counts of a chain that rarely
switches. Fix a start law and run the chain for $N$ steps with transition matrix
$I+(T/N)Q$. With $T=\tau\log N$ the chain makes about $\tau\log N$ switches. Its
composition is the vector of occupation counts, and the support of a composition is
a face of the simplex. In~\cite{paperC1} we describe the law of the composition on
every face at this scale. On a face $F$ the most likely composition has probability
$N^{-\Phi_\tau(F)+o(1)}$, with the score $\Phi_\tau(F)=|F|-1+\tau\lambda_F$ and
$\lambda_F$ the principal exit rate of $F$. One power of $N$ per dimension competes
with the cost of staying in the face, and the mode moves to larger faces as $\tau$
grows. Inside a face the logarithm of the law is $T$ times $J_V$.

We ask which coarse functionals of the law of the counts determine $Q$, and what
each of them loses. We compare 3, from fine to coarse. The first is the law itself,
at every $N$ and $T$. The second is its logarithmic scale on single compositions,
the functional $J_V$. The third is a finite label, the support of the most likely
composition when the chain is killed at known rates. The third determines exactly
the same minors as the second (Theorem~\ref{thm:B}). We then read a reversible
generator from finitely many boundaries of a design fixed in advance
(Theorem~\ref{thm:C}), and the Freidlin and Wentzell cycle tree of an Arrhenius
family from the same labels (Theorem~\ref{thm:D}).

\subsection{The Donsker and Varadhan functional}\label{sec:intro-dv}

Write $Q^*=D_\pi^{-1}Q^\top D_\pi$ for the time reversal of an irreducible
generator with stationary law $\pi$. Call a partition $[d]=X\cup Y$ with
$|X|,|Y|\ge2$ a cut of $Q$ if both blocks $Q[X,Y]$ and $Q[Y,X]$ have rank at most
$1$. Across a cut a cut-transpose replaces the 2 rank-one blocks by the products
of their factors taken crosswise and transposes the block on $Y$ (Section~\ref{sec:fiber}).

\begin{introtheorem}[what the Donsker and Varadhan functional determines]\label{thm:A}
Let $Q$ and $Q'$ be irreducible generators on $[d]$.
\begin{enumerate}[(i)]
\item $J_V=J'_V$ on the simplex if and only if $-Q$ and $-Q'$ have the same
principal minors.
\item The irreducible generators with the principal minors of $-Q$ form a finite
set. It contains $Q$ and $Q^*$. It is $\{Q,Q^*\}$ when $Q$ has no cut, in
particular when $d\le3$, and $\{Q\}$ when $Q$ is reversible. In general it is the
set of generator points of the cut-transpose orbit of $-Q$, and it can contain
generators other than $Q$ and $Q^*$.
\item From the stationary start the law of the composition is the same under $Q$
and under $Q^*$, for every $N$ and every $T$ with $Tq_{\max}\le N$.
\end{enumerate}
\end{introtheorem}

Part (i) is Theorem~\ref{thm:dv}, part (ii) is Theorem~\ref{thm:nonreversible}
with the example after it, and part (iii) is Proposition~\ref{prop:stationary}.
Part (i) has 3 inputs. The first is the variational formula
$\lambda_V(v)=\min_\alpha(J_V(\alpha)-v\cdot\alpha)$ for
$\lambda_V(v)=-\varrho(Q+\operatorname{diag}v)$. The second is the irreducibility
of the multiaffine polynomial $\det(-Q_F-\operatorname{diag}w)$ for every
irreducible face, which is Lemma~5.1 of Al Ahmadieh~\cite{alahmadieh2023}. The third
is the Perron-branch lemma (Lemma~\ref{lem:perron-branch}). It says that the
Perron root on an open set of fields determines that polynomial, hence every
principal minor. Kim and Landim use the minors of orders 1 and 2, with the
stationary law, to fix a reversible generator. By their Corollary~5.8 the joint
functional of the empirical measure and the empirical flow determines the
generator, so adding the flow removes the whole fiber. Part (iii) shows that the
fiber is not an artefact of the logarithmic scale. From the stationary start even
the exact law of the counts cannot separate $Q$ from $Q^*$. For a start outside a
finite union of proper subspaces the second-order amplitudes of the local law do
separate every member of the fiber (Proposition~\ref{prop:amplitudes}).

\subsection{Modal supports}\label{sec:intro-modal}

Weight each composition by $e^{Tv\cdot n/N}$ for a field $v\le0$. This is the
chance that the chain survives when it is killed at state $i$ at the known rate
$-v_i$. We record only which face carries the most likely surviving composition,
as a function of $(\tau,v)$, and call this the modal-support map. The data are the
boundaries between the open regions where one face wins strictly
(Definition~\ref{def:controlled}).

\begin{introtheorem}[what the modal supports determine]\label{thm:B}
Let $Q$ and $Q'$ be irreducible generators and let the start have full support.
They have the same modal-support map, outside a closed nowhere dense set of
$(\tau,v)$, if and only if $-Q$ and $-Q'$ have the same principal minors. If the
minors differ, there is a nonempty open set of $(\tau,v)$ on which both maps have
a strict winner and the 2 winners differ.
\end{introtheorem}

This is Theorem~\ref{thm:modal}, with Proposition~\ref{prop:identifiability} for
the last sentence. The proof is a descent through the faces. For every irreducible
face $F$ and a generic field on $F$, the exposure lemma (Lemma~\ref{lem:exposure})
finds a proper subface $G$ that shares a boundary of strict regions with $F$. Along
that boundary $\lambda_F=\lambda_G-(|F|-|G|)/\tau$. Induction on $|F|$ then traces
$\lambda_F$ on an open set of fields, and the Perron-branch lemma extends it to all
fields.

Theorem~\ref{thm:B} is a statement about retention, not about sampling. The
survival probability under calibrated killing is a finer observation. It equals
$e^{-T\lambda_V(v)}$ up to a bounded factor, so it gives every principal minor at
once, with no descent (Remark~\ref{rem:survival}). The modal support is not easier
to sample. The most likely composition carries probability $N^{-\Phi+o(1)}$, so
estimating which face carries it needs far more independent runs than estimating
the survival probability. We know of no experiment in which the modal support is
observed and the survival probability is not. But the forward theory
of~\cite{paperC1} outputs the support of the most likely composition, and so does
any exact optimizer that reports only a maximizer of the weighted law.
Theorem~\ref{thm:B} says that this output loses nothing that $J_V$ keeps.

\emph{What is borrowed and what is new.} The principal minors are not observed.
Theorem~\ref{thm:dv} reads all of them off $J_V$, and Theorem~\ref{thm:modal} reads
all of them off the modal-support map, through the exposure lemma and the
Perron-branch lemma. This is our main contribution. The fiber of the
principal-minor map of an irreducible matrix is known, up to diagonal similarity and
transposition, from Loewy~\cite{loewy1986}, Hartfiel and
Loewy~\cite{hartfielloewy1984}, Al Ahmadieh~\cite{alahmadieh2023} and Chatterjee,
Ghosh, Gurjar and Raj~\cite{cggr2024}, who describe it by sequences of
cut-transposes (Section~\ref{sec:fiber}). Our second contribution, and the shortest
step, translates this fiber to generators. A Perron
and Frobenius normalization picks exactly one generator in each class of diagonal
similarity. So the generator fiber is finite, every member is irreducible with the
spectrum of $Q$, and a cut-transpose acts as $Q$ on one side of the cut and as $Q^*$
on the other (Theorem~\ref{thm:nonreversible} and
Proposition~\ref{prop:amplitudes}). The third is the stationary obstruction of
part (iii) of Theorem~\ref{thm:A}.

\subsection{A reversible design and finite horizons}\label{sec:intro-reversible}

For a reversible chain on a known connected graph with $m$ edges the fields can be
fixed before anything is observed.

\begin{introtheorem}[a reversible design]\label{thm:C}
Let $Q$ be reversible and irreducible on a known connected graph with $m$ edges,
and suppose that known constants $R$ and $b$ bound the norm of the symmetrized
generator above and every edge product $\sqrt{q_{ij}q_{ji}}$ below.
\begin{enumerate}[(i)]
\item A design whose fields and time scales depend only on $R$, $b$ and the
labels reads $d-1$ singleton switches and $m$ pair-entry times off the
modal-support map, and these $m+d-1$ numbers determine $Q$.
\item No family of fewer than $m+d-1$ scalar $C^1$ observations, with finitely many
locally constant labels, identifies $Q$ on any open set of such generators.
\item The reconstruction is Lipschitz in the observed coordinates, with constants
that depend only on $d$ and $b/R$.
\item At the finite horizon $N$, with $L=\log N$, the observed pair-entry times are
$t_{ij}(1-\log L/(2L))+c_{ij}/L+O(\log L/L^2)$, uniformly over the generators of
the class and the starts with $\min_i\mu_i\ge\mu_*$, and the start bound cannot be
dropped.
\end{enumerate}
\end{introtheorem}

Parts (i) to (iii) are Theorems~\ref{thm:reversible-reconstruction},
\ref{thm:minimality} and~\ref{thm:stability}, and part (iv) is
Theorem~\ref{thm:finite-horizon} with Remark~\ref{rem:finite-start}. The coordinates
are those of Wolfer and Watanabe~\cite{wolferwatanabe2021}, and Kim and
Landim~\cite[Lemma~5.1 and Corollary~5.3]{kl2024} recover the same quantities from
$J_V$. Theorem~\ref{thm:C} adds a design fixed in advance, its minimality, its
stability with constants from $R$ and $b$ alone, and its finite-horizon behaviour.

\subsection{Low temperature and the cycle tree}\label{sec:intro-lowtemp}

Let the rates be $q_{ij}=e^{-V_{ij}/\eta}$ with fixed costs $V_{ij}\ge0$ and a
temperature $\eta\downarrow0$, and let the time scale be $\tau=e^{h/\eta}$ for a
height $h$. Then dimension and exit rate compete on the Freidlin and Wentzell cycles
of the chain. The depth $D(F)$ of a face is the exponential order of its exit rate,
and it is the largest depth of a cycle inside $F$ (Theorem~\ref{thm:depth}). Away
from finitely many heights the winner is a cycle of least size among those deeper
than $h$ (Theorem~\ref{thm:noncritical}). At a critical height explicit prefactors
decide, and a noncycle can win on an open window of the rescaled time scale
(Theorem~\ref{thm:critical-window} and Example~\ref{ex:four-state}).

So modal winners need not be cycles. Do the modal supports determine the cycles at
all? The 2 obvious experiments do not. The unperturbed modal path shows the depth of
each face it selects and nothing about the cycles that are never the smallest
eligible face. Killing the complement of a face at a large rate does not show its
depth either. The whole space then has an exit rate of the same order, and the
winner leaves the face only at a larger height. The cure is to reward the face
instead. For a proper face $G$, a constant $c>0$ and a real $\gamma$ we put the field
$v^{G,\gamma}=c\,e^{-\gamma/\eta}\ones_G$. The threshold $\vartheta(G)$ is the
supremum of the $\gamma$ for which $G$ wins under $v^{G,\gamma}$ at some height above
$\gamma$.

\begin{introtheorem}[the cycle tree from calibrated rewards]\label{thm:D}
Let $(Q_\eta)$ be an Arrhenius family with all rates positive, and let the start have
full support.
\begin{enumerate}[(i)]
\item $\vartheta(G)\le D(G)$ for every proper face $G$, with equality when $G$ is a
cycle.
\item $D(F)=\max_{G\subseteq F}\vartheta(G)$ for every proper face $F$.
\item A nonempty proper set $C$ is a cycle if and only if $D(C')<D(C)$ for every
nonempty proper subset $C'$ of $C$.
\end{enumerate}
So the modal-support map, restricted to the fields $v^{G,\gamma}$, determines the
depth of every face, the cycles, their depths and the cycle tree. No step uses
reversibility, and no nondegeneracy of the depths is assumed.
\end{introtheorem}

Parts (i) and (ii) are Theorem~\ref{thm:attraction} and part (iii) is
Theorem~\ref{thm:cycles-from-depths}. The design needs 3 calibrated quantities,
the time unit, the amplitude of the reward and the labels of the states, and no rate
of $Q$. The reward keeps a face whose rate is known, namely the whole space with
$\lambda_V(v)\le0$, and the killing design destroys exactly this. Criterion (iii)
holds with ties between the depths of different faces, because Catoni's cycles are
strict and depths increase strictly along nested cycles (Lemma~\ref{lem:nested}). Its
weak form, with $D(C')\le D(C)$, fails already on 3 states. Finally the endpoints
of the critical windows determine the prefactors of the faces on the lower envelope,
and only bound the others (Theorem~\ref{thm:window-prefactors}).

\subsection{Related work}\label{sec:related}

The rate inside a face is the Donsker and Varadhan functional of the killed
chain~\cite{dv1975}, and the exact law of the counts goes back to
Good~\cite{good1961}. Bertini, Gabrielli and Landim~\cite{bgl2024} and
Landim~\cite{landim2023} expand the level-two rate functions of finite chains in a
parameter, and Kim and Landim~\cite{kl2024} treat the joint rate of the empirical
measure and the flow~\cite{bfg2014,bfg2015}. Their Section~5 is the starting point of
Theorem~\ref{thm:A}. The matrix theory behind the fiber is credited above and in
Section~\ref{sec:fiber}, together with Engel and Schneider~\cite{engelschneider1980}.
Cohen~\cite{cohen1981} and Friedland~\cite{friedland1981} study the convexity of
$v\mapsto\varrho(A+\operatorname{diag}v)$, and Delmas, Dronnier and Zitt~\cite{ddz2024}
prove a multiplicative analogue of the Perron-branch lemma. The cycle framework is that of
Freidlin and Wentzell~\cite{freidlinwentzell2012}, Catoni~\cite{catoni1999}, Olivieri and
Vares~\cite{olivierivares2005}, Trouv\'e~\cite{trouve1996} and Cirillo, Nardi and
Sohier~\cite{cirillonardisohier2015}. The eigenvalue asymptotics behind the prefactors go
back to Wentzell~\cite{wentzell1972}. New at low temperature are the score of dimension
against exit rate, the competition of envelopes at a critical depth with its noncycle window,
and the reading of the cycle tree from calibrated rewards.

\subsection{Plan of the paper}\label{sec:plan}

Section~\ref{sec:prelim} fixes the notation and states the results of~\cite{paperC1}
that we use. They are the first-order results, the face maxima with a field, their
uniform version on a compact class and the constant of the local law.
Section~\ref{sec:minors} proves Theorem~\ref{thm:A}. Section~\ref{sec:modal} sets up
the modal-support map and proves Theorem~\ref{thm:B}. Section~\ref{sec:reversible}
treats the reversible design and the finite horizon, and Section~\ref{sec:lowtemp} the
low temperature. Section~\ref{sec:verify} lists the computations that check the paper
and the open questions.

\section{Setting and imported results}\label{sec:prelim}

We fix the notation and state, without proof, the results of~\cite{paperC1} that we
use, each with a pointer to its place there.

\subsection{Setting}\label{sec:notation}

\emph{The chain.} Fix $d\ge2$ and a generator $Q$ on $[d]=\{1,\dots,d\}$. Its rates
are $q_{ij}\ge0$ for $i\ne j$, the exit rates are $q_i=\sum_{j\ne i}q_{ij}$,
$Q_{ii}=-q_i$, and $q_{\max}=\max_iq_i$. We write $q$ for the vector $(q_i)$. The
start $\mu$ is a probability law on $[d]$, and $V=[d]$ is the face of all states.
For $N\ge2$ and $T>0$ put $L=\log N$ and $\tau=T/L$. The chain
$X_1,\dots,X_N$ has $X_1\sim\mu$ and transition matrix $I+(T/N)Q$, which is
stochastic when $Tq_{\max}\le N$. The composition is the vector $n$ of visit counts,
$n_i=\#\{t\le N:X_t=i\}$, so $n_1+\dots+n_d=N$. We write $P_N(n)$ for its law,
$\supp n=\{i:n_i>0\}$ and $\alpha=n/N$ for the visit fractions.

\emph{Faces.} A face is a nonempty set $F\subseteq[d]$. Let $\Delta_F$ be the set of
probability vectors on $[d]$ that vanish off $F$, and $\Delta_F^\circ$ the set of
those that are positive on $F$. The killed generator of $F$ is
$Q_F=(Q_{ij})_{i,j\in F}$. Its diagonal keeps the full exit rates $q_i$, so it
generates the chain killed when it leaves $F$. For $A\subseteq F$ the principal
submatrix of $-Q_F$ on $A$ is $-Q_A$. The switching graph $G_F$ has an edge $i\to j$
whenever $i,j\in F$, $i\ne j$ and $q_{ij}>0$. We call $F$ irreducible if $G_F$ is
strongly connected, so every singleton is irreducible. We call $F$ admissible if
some path of positive probability has support exactly $F$, and we say that $F$
satisfies (H) if it is irreducible and $\mu(F)>0$.

\emph{Perron roots and fields.} A square matrix $B$ is Metzler if its off-diagonal
entries are nonnegative, and its Perron root $\varrho(B)$ is then its largest real
eigenvalue. If $B$ is irreducible, $\varrho(B)$ is simple and has positive left and
right eigenvectors~\cite{bermanplemmons1994}. For a face $F$ and a field
$v\in\R^d$ put
\[
 \lambda_F(v)=-\varrho\bigl(Q_F+\operatorname{diag}v_F\bigr),\qquad
 \Phi_{\tau,v}(F)=|F|-1+\tau\lambda_F(v),
\]
and $\lambda_F=\lambda_F(0)$, $\Phi_\tau=\Phi_{\tau,0}$. The number $\lambda_F$ is the
exit rate of $F$, and $\lambda_{\{i\}}(v)=q_i-v_i$. If $F$ is reducible, $Q_F$ is block
triangular in an order of its strongly connected components, and $\lambda_F(v)$ is the
least of the numbers $\lambda_C(v)$ over these components $C$. The shift rule
$\lambda_F(v+c\ones)=\lambda_F(v)-c$ holds for every face. If $F$ is irreducible, let
$r$ and $l$ be positive right and left Perron vectors of $Q_F+\operatorname{diag}v_F$
with $l\cdot r=1$, and put $\alpha^*_F(v)=l\circ r$, the quasi-ergodic law of $F$ under
the field, where $x\circ y$ is the entrywise product. Finally
\[
 M_F(T,v)=\max\bigl\{P_N(n)\,e^{Tv\cdot n/N}:\supp n=F\bigr\},\qquad\max\emptyset=0 .
\]

\emph{The functional.} For $\alpha\in\Delta_F$ let
\[
 J_F(\alpha)=\sup_{u\in(0,\infty)^F}\sum_{i\in F}\alpha_i\Bigl(-\frac{(Q_Fu)_i}{u_i}\Bigr),
\]
the Donsker and Varadhan functional of the chain killed on leaving $F$~\cite{dv1975},
and $J_V$ that of the whole chain.

\subsection{Imported results}\label{sec:imported}

\begin{lemma}[admissible faces, {\cite[Lemma~2.2]{paperC1}}]\label{lem:admissible}
\textup{(a)} An irreducible face is admissible if and only if $\mu(F)>0$. So
\textup{(H)} holds exactly for the irreducible admissible faces, and if $\mu$ has full
support every irreducible face is admissible. \textup{(b)} If $Q$ has symmetric
support, that is $q_{ij}>0$ exactly when $q_{ji}>0$, then every admissible face is
irreducible. \textup{(c)} A face that is not admissible carries no composition of
positive probability.
\end{lemma}

Here (a) and (b) are~\cite[Lemma~2.2(b) and~(c)]{paperC1}, and (c) is immediate from the
definition, as noted in~\cite[Theorem~4.2(ii)]{paperC1}.

\begin{lemma}[monotonicity in the face, {\cite[Lemma~3.2]{paperC1}}]\label{lem:face-monotone}
If $F\subsetneq G$ are faces and $G$ is irreducible, then
$\lambda_G(v)<\lambda_F(v)$ for every field $v$. In particular every proper face of an
irreducible chain has $\lambda_F>\lambda_V=0$.
\end{lemma}

\begin{proposition}[the variational formula, {\cite[Propositions~3.1 and~3.4]{paperC1}}]\label{prop:variational}
Let $F$ be a face.
\begin{enumerate}[(a)]
\item $J_F=J_V$ on $\Delta_F$.
\item If $F$ is irreducible, then for every $v\in\R^F$
\[
 \lambda_F(v)=\min_{\alpha\in\Delta_F}\bigl(J_F(\alpha)-v\cdot\alpha\bigr),
\]
with the minimum attained only at $\alpha^*_F(v)$. So $-\lambda_F$ is the Legendre
transform of $J_F$, and $J_F$ is the Legendre transform of $v\mapsto-\lambda_F(v)$ on
$\Delta_F^\circ$. The map $v\mapsto\lambda_F(v)$ is concave and real-analytic, with
$\partial\lambda_F/\partial v_i=-\alpha^*_F(v)_i$.
\item If $F$ is irreducible, $Q_F$ is reversible with respect to some $\nu\in(0,\infty)^F$ and
$S_F=D_{\sqrt\nu}(-Q_F)D_{\sqrt\nu}^{-1}$, then
$\lambda_F(v)=\lambda_{\min}(S_F-\operatorname{diag}v_F)$.
\end{enumerate}
\end{proposition}

The derivative in (b) is the first-order perturbation $\partial\varrho/\partial v_i=l_ir_i$
of a simple eigenvalue, with $l\cdot r=1$.

\begin{lemma}[a uniform upper bound, {\cite[Lemma~4.1 and~(6)]{paperC1}}]\label{lem:upper}
Let $F$ be irreducible with $k=|F|\ge2$ and let $\mathcal V\subset\R^d$ be compact. There
is $C$, depending on $Q$, $\mu$ and $\mathcal V$, such that for all $v\in\mathcal V$, all
$N\ge2$ and $T>0$ with $Tq_{\max}\le N/2$, and every $n$ with $\supp n=F$,
\begin{equation}\label{eq:upper-field}
 P_N(n)\,e^{Tv\cdot n/N}\le C\Bigl(\frac TN\Bigr)^{k-1}e^{-T\lambda_F(v)+CT^2/N}.
\end{equation}
The bound comes from the Perron tilt $\theta=q-v-\lambda_F(v)\ones$, and it holds for a path
of any length in place of $N$, hence for each block of a path.
\end{lemma}

For a singleton the only composition is $Ne_i$, so by~\cite[(1)]{paperC1}
\begin{equation}\label{eq:vertexmax}
 M_{\{i\}}(T,v)=\mu_i\Bigl(1-\frac{Tq_i}N\Bigr)^{N-1}e^{Tv_i}
 =\mu_ie^{-T(q_i-v_i)}\bigl(1+O(T^2/N)\bigr).
\end{equation}

\begin{theorem}[first-order selection with a field, {\cite[Theorem~4.2 and Remark~4.3]{paperC1}}]\label{thm:first}
Fix $Q$, $\mu$ and a field $v$.
\begin{enumerate}[(a)]
\item Let $F$ be admissible with $k=|F|$. There are $\kappa$, $T_0$ and $N_0$ such that
$\log M_F(T,v)\le-(k-1)L+(k-1)\log T-T\lambda_F(v)+\kappa$ for $N\ge N_0$ and
$T_0\le T\le N^{1/3}$. If $F$ is not admissible, then $M_F(T,v)=0$.
\item Let $T=\tau L$ with $\tau>0$ fixed, and let $v=0$ or let $\mu$ have full support. Then
\[
 \log\max_nP_N(n)e^{Tv\cdot n/N}=-L\min_F\Phi_{\tau,v}(F)+O(\log L)
\]
over the admissible faces, and for large $N$ the support of every maximizer minimizes
$\Phi_{\tau,v}$ among them.
\end{enumerate}
\end{theorem}

At $v=0$, part (a) is the upper bound of~\cite[Theorem~4.2(ii)]{paperC1} and (b)
is~\cite[Theorem~4.2(iii)]{paperC1}. The field forms of both are~\cite[Remark~4.3]{paperC1}, which
derives them from the same proofs with $Q_F+\operatorname{diag}v_F$ in place of $Q_F$. We use (b)
only at a fixed $\tau$.

\begin{definition}[tilt and constant, {\cite[Definition~5.1]{paperC1}}]\label{def:tilt}
Let $F$ be irreducible with $k=|F|\ge2$, with its states numbered $1,\dots,k$. For $\alpha\in\Delta_F^\circ$ the supremum
defining $J_F(\alpha)$ is attained at a vector $r_\alpha$, unique up to a factor. Put
$l_\alpha=\alpha/r_\alpha$ entrywise, and let $\widehat Q_\alpha$ be the generator on $F$
with off-diagonal entries $q_{ij}r_{\alpha,j}/r_{\alpha,i}$, whose stationary law is
$\alpha$. With $D=\operatorname{diag}\alpha$ let
$Z_\alpha=(\ones\alpha^\top-\widehat Q_\alpha)^{-1}-\ones\alpha^\top$ and
$\Gamma_\alpha=DZ_\alpha+Z_\alpha^\top D$, and let $\Sigma_\alpha$ be $\Gamma_\alpha$ with
the row and column of state $1$ deleted. Then
\[
 C_F(\alpha)=\frac{(\mu\cdot r_\alpha)(l_\alpha\cdot\ones)}{(2\pi)^{(k-1)/2}\sqrt{\det\Sigma_\alpha}},
 \qquad K_F(v)=C_F\bigl(\alpha^*_F(v)\bigr),
\]
and $K_{\{i\}}(v)=\mu_i$.
\end{definition}

\begin{proposition}[the constant, {\cite[Propositions~3.3 and~5.2]{paperC1}}]\label{prop:constant}
Let $F$ be irreducible with $k\ge2$. The maps $\alpha\mapsto r_\alpha$ (normalized),
$l_\alpha$, $\Sigma_\alpha$, $J_F(\alpha)$ and $C_F(\alpha)$ are real-analytic on
$\Delta_F^\circ$, $\Sigma_\alpha$ is positive definite, and $\nabla^2J_F(\alpha)=\Sigma_\alpha^{-1}$ in the
coordinates $(\alpha_2,\dots,\alpha_k)$. At $\alpha=\alpha^*_F(v)$ the vector $r_\alpha$ is
a multiple of the right Perron vector of $Q_F+\operatorname{diag}v_F$, and $l_\alpha$ is
the matching multiple of the left one.
\end{proposition}

The first sentence is~\cite[Proposition~5.2(a)]{paperC1}. For the second let $r$ and $l$ be the
right and left Perron vectors of $Q_F+\operatorname{diag}v_F$ with $l\cdot r=1$, so that
$\alpha^*_F(v)=l\circ r$. The vector $u=r$ gives $-(Q_Fr)_i/r_i=v_i+\lambda_F(v)$, so its value in
the supremum defining $J_F(\alpha^*_F(v))$ is $\lambda_F(v)+v\cdot\alpha^*_F(v)$, which is
$J_F(\alpha^*_F(v))$ by~\cite[Proposition~3.3(a)]{paperC1}. So $r$ attains the supremum, and by
the uniqueness in~\cite[Proposition~3.3(b)]{paperC1} it is a multiple of $r_\alpha$. Then
$l_\alpha=\alpha/r_\alpha$ is the matching multiple of $l$. At $v=0$ this is the text after
\cite[Definition~5.1]{paperC1}.

\begin{theorem}[face maxima with a field, {\cite[Theorem~5.7]{paperC1}}]\label{thm:facemax}
Let $F$ satisfy \textup{(H)} with $k=|F|$, and let $\mathcal V\subset\R^d$ be compact.
Uniformly for $v\in\mathcal V$ and $1\le T\le N^{1/3}$,
\begin{equation}\label{eq:logmax}
 \log M_F(T,v)=-(k-1)L+\tfrac{k-1}2\log T+\log K_F(v)-T\lambda_F(v)
 +O_{\mathcal V}\Bigl(\frac1T+\frac{T^2}N\Bigr),
\end{equation}
and for $k\ge2$ every maximizer $\hat n$ satisfies $\hat n/N=\alpha^*_F(v)+O_{\mathcal V}(1/T)$.
\end{theorem}

\begin{theorem}[the stratified local law, {\cite[Theorem~5.6]{paperC1}}]\label{thm:stratified}
For every face $F$ satisfying \textup{(H)} with $k=|F|\ge2$ and every compact
$\mathcal K\subset\Delta_F^\circ$, uniformly for $\supp n=F$, $n/N\in\mathcal K$ and
$1\le T\le N^{1/3}$,
\[
 -\log P_N(n)=(k-1)L+T\,J_V(n/N)-\frac{k-1}2\log T-\log C_F(n/N)
 +O_{\mathcal K}\Bigl(\frac1T+\frac{T^2}N\Bigr).
\]
\end{theorem}

By Theorem~\ref{thm:stratified}, $J_V$ is the logarithmic scale of the law on single
compositions.\ifarxivbuild\ Section~\ref{sec:finite-horizon} also imports the uniform version
of Theorem~\ref{thm:facemax} on a compact class of reversible chains
(Proposition~\ref{prop:uniform-facemax}). It is stated there because it needs the design of
Section~\ref{sec:reversible}.\fi\ Nothing else from~\cite{paperC1} enters.

\section{The Donsker and Varadhan functional and the principal minors}\label{sec:minors}

This section proves Theorem~\ref{thm:A}. In Section~\ref{sec:perron-branch} the Perron
root of a face on an open set of fields determines the principal minors of the face.
Section~\ref{sec:dv-minors} deduces that $J_V$ determines exactly the principal minors of
$-Q$. Section~\ref{sec:fiber} describes the generators with given minors, and
Section~\ref{sec:stationary} records 2 first facts on the exact law, which
Section~\ref{sec:exactlaw} completes.

\subsection{The polynomial of a face and the Perron branch}\label{sec:perron-branch}

The principal minors of $-Q_F$ are the numbers $\det(-Q_A)$ for $A\subseteq F$, with
$\det(-Q_\emptyset)=1$. They are the coefficients of the multiaffine polynomial
\[
 R_F(w)=\det\bigl(-Q_F-\operatorname{diag}w\bigr)
 =\sum_{A\subseteq F}\det(-Q_A)\prod_{i\in F\setminus A}(-w_i),\qquad w\in\R^F .
\]
For every field $v$, $\lambda_F(v)$ is a real root of $t\mapsto R_F(v+t\ones)$, and if
$F$ is irreducible it is the smallest real root, since the other real eigenvalues of
$-Q_F-\operatorname{diag}v_F$ exceed the Perron eigenvalue.

\begin{lemma}[irreducibility]\label{lem:RF-irreducible}
If $F$ is irreducible, then $R_F$ is irreducible over $\mathbb C$. If $F$ is reducible,
$R_F$ is the product of the polynomials of the strongly connected components of $G_F$.
\end{lemma}

\begin{proof}
Al Ahmadieh~\cite[Lemma~5.1]{alahmadieh2023} proves, over any field and for every
order, that a square matrix $A$ is irreducible exactly when
$\det(\operatorname{diag}x+A)$ is an irreducible polynomial in $x$. Take the field
$\mathbb C$ and $A=-Q_F$. Then $R_F(w)=\det(\operatorname{diag}(-w)+A)$, and the linear
change of variables $x=-w$ preserves irreducibility. For a reducible face, $-Q_F$ is
block triangular with the blocks of the components on the diagonal, and the determinant
factors accordingly.
\end{proof}

The proof of Al Ahmadieh rests on a theorem of Hartfiel and
Loewy~\cite[Theorem~1]{hartfielloewy1984}, that every entry of the adjugate of
$\operatorname{diag}x+A$ is a nonzero polynomial when $A$ is irreducible.

\begin{lemma}[the Perron branch]\label{lem:perron-branch}
Let $F$ be irreducible, let $Q'$ be another generator whose face $F$ is irreducible,
and let $c\in\R$. Suppose that $\lambda'_F(v)=\lambda_F(v)+c$ for all $v$ in a nonempty
open set $\mathcal U\subseteq\R^F$. Then $-Q'_F$ and $-Q_F+cI$ have the same principal
minors, and $\lambda'_F=\lambda_F+c$ on all of $\R^F$. In particular, for $c=0$,
$-Q'_F$ and $-Q_F$ have the same principal minors.
\end{lemma}

\begin{proof}
Let $k=|F|$ and number the states of $F$ as $1,\dots,k$. Put
$\widetilde R(w)=\det(-Q_F+cI-\operatorname{diag}w)=R_F(w-c\ones)$, which is
irreducible over $\mathbb C$ by Lemma~\ref{lem:RF-irreducible}. For $k=1$ the hypothesis
says $q'_1=q_1+c$, which is the claim. Let $k\ge2$, and put
$\psi(v)=v+\lambda'_F(v)\ones$. For $v\in\mathcal U$ the number
$\lambda'_F(v)=\lambda_F(v)+c$ is an eigenvalue of $-Q'_F-\operatorname{diag}v$ and of
$-Q_F+cI-\operatorname{diag}v$, so $R'_F(\psi(v))=0=\widetilde R(\psi(v))$.

Both polynomials are affine in $w_1$. Write $\widetilde R=p_1w_1+p_0$ and
$R'_F=p'_1w_1+p'_0$ with $p_1,p_0,p'_1,p'_0$ polynomials in $w'=(w_2,\dots,w_k)$. The
coefficient of $\prod_{i\ge2}(-w_i)$ in $p_1$ is $-1$, so $p_1\ne0$. The polynomial
$\mathcal D=p_1R'_F-p'_1\widetilde R=p_1p'_0-p'_1p_0$ does not involve $w_1$ and vanishes
on $\psi(\mathcal U)$. We claim that the projection of $\psi(\mathcal U)$ to $w'$
contains a nonempty open set. Fix $v^0\in\mathcal U$ and consider
$y\mapsto(\psi(v^0+(0,y)))_{i\ge2}$ for small $y\in\R^{k-1}$. By
Proposition~\ref{prop:variational}(b), $\lambda'_F$ is real-analytic with
$\partial\lambda'_F/\partial v_i=-\beta_i$, where $\beta=\alpha^{\prime*}_F(v)>0$ and
$\sum_i\beta_i=1$. The derivative of the map at $y=0$ is $I-\ones\beta'^\top$ with
$\beta'=(\beta_i)_{i\ge2}$, of determinant $\beta_1>0$, and the inverse function theorem
proves the claim. So $\mathcal D=0$ and $p_1R'_F=p'_1\widetilde R$. The irreducible
$\widetilde R$ divides $p_1R'_F$, and it involves $w_1$, so it does not divide $p_1\ne0$,
which is free of $w_1$. Hence it divides $R'_F$. Both have total degree $k$ and the
coefficient $1$ at $\prod_i(-w_i)$, so $R'_F=\widetilde R$. Finally $\lambda'_F(v)$ is the
smallest real root of $t\mapsto R'_F(v+t\ones)=\widetilde R(v+t\ones)$, which is
$\lambda_F(v)+c$.
\end{proof}

Delmas, Dronnier and Zitt~\cite[Theorem~2.1]{ddz2024} prove a multiplicative analogue
for $\eta\mapsto\varrho(K\operatorname{diag}\eta)$ on $\{0,1\}^n$, by an induction over the
cube. Cohen~\cite{cohen1981} proved that $v\mapsto\varrho(A+\operatorname{diag}v)$ is convex
for a Metzler matrix $A$, and Friedland~\cite[Theorem~4.1]{friedland1981} showed that for
irreducible $A$ it is affine only along $\ones$. Lemma~\ref{lem:perron-branch} is the inverse
statement, that the map on an open set determines the principal minors of $A$. We did not
find it in these papers.

\subsection{The functional determines exactly the principal minors}\label{sec:dv-minors}

\begin{theorem}[the functional and the minors]\label{thm:dv}
For two irreducible generators $Q$ and $Q'$ on $[d]$ the following are equivalent.
\begin{enumerate}[(i)]
\item $J_V=J'_V$ on $\Delta_V$.
\item $\lambda_V=\lambda'_V$ on a nonempty open set of fields.
\item $-Q$ and $-Q'$ have the same principal minors.
\end{enumerate}
If they hold, then $\lambda_F=\lambda'_F$ on $\R^F$ and $J_F=J'_F$ on $\Delta_F$ for every
face $F$, and $F$ is irreducible for $Q$ exactly when it is irreducible for $Q'$.
\end{theorem}

\begin{proof}
(i) implies (ii). By Proposition~\ref{prop:variational}(b),
$\lambda_V(v)=\min_{\Delta_V}(J_V-v\cdot\alpha)=\lambda'_V(v)$ for every $v$.

(ii) implies (iii). This is Lemma~\ref{lem:perron-branch} with $F=V$ and $c=0$.

(iii) implies (i). Equal principal minors give $R'_F=R_F$ for every face $F$, since the
minors of $-Q_F$ are among those of $-Q$. By Lemma~\ref{lem:RF-irreducible}, $F$ is
irreducible for $Q$ exactly when $R_F$ is irreducible, which is the last claim. On an
irreducible face $\lambda'_F(v)$ and $\lambda_F(v)$ are the smallest real roots of the same
polynomial $t\mapsto R_F(v+t\ones)$, so they agree. For a reducible face, the factors of
$R_F=R'_F$ given by Lemma~\ref{lem:RF-irreducible} are irreducible, so by unique
factorization the components of $G_F$ have the same vertex sets for both generators, and
$\lambda_F(v)$, the least $\lambda_C(v)$ over the components, agrees too. By Proposition~\ref{prop:variational}(b), $J_F$ is the Legendre transform of
$v\mapsto-\lambda_F(v)$ on $\Delta_F^\circ$, so $J'_F=J_F$ there for irreducible $F$. On the
boundary we use that $J_F$ is convex and lower semicontinuous, as a supremum of affine
functions of $\alpha$. For $\alpha\in\Delta_F$ and $\beta\in\Delta_F^\circ$ the points
$\alpha_t=(1-t)\alpha+t\beta$ with $t\in(0,1]$ are interior. Convexity gives
$\limsup_{t\downarrow0}J_F(\alpha_t)\le J_F(\alpha)$, and lower semicontinuity the reverse
inequality. So $J_F(\alpha)=\lim_{t\downarrow0}J_F(\alpha_t)$, likewise for $J'_F$, and
$J'_F=J_F$ on $\Delta_F$. With $F=V$ this is (i). For a reducible face,
Proposition~\ref{prop:variational}(a) gives $J_F=J_V=J'_V=J'_F$ on $\Delta_F$.
\end{proof}

\begin{remark}[the survival probability]\label{rem:survival}
A finer observation than $J_V$ gives the same minors more directly. Let $v$ be a field
and let the start have full support. The weighted mass
$\sum_nP_N(n)e^{Tv\cdot n/N}$ is the probability that the chain survives when it is
killed at each visit to $i$ with probability $1-e^{Tv_i/N}$, for $v\le0$. As $N\to\infty$ at
fixed $T$ it tends to $\mu^\top e^{T(Q+\operatorname{diag}v)}\ones$. For any generator $Q'$
the limit $-T^{-1}\log\mu^\top e^{T(Q'+\operatorname{diag}v)}\ones$ as $T\to\infty$ is
$\lambda'_V(v)=-\varrho(Q'+\operatorname{diag}v)$, the least exit rate of the components of
$Q'$, because the start and the vector $\ones$ are positive. So the survival probabilities
on all fields give $\lambda'_V$ on $\R^d$. Suppose that they agree with those of the
irreducible $Q$. Then $Q'$ is irreducible. Otherwise $\lambda'_V=\lambda'_C$ for some
component $C\ne V$ on an open set of fields, and there $\lambda'_V$ does not depend on
$v_j$ for $j\notin C$, while $\partial\lambda_V/\partial v_j=-\alpha^*_V(v)_j<0$ by
Proposition~\ref{prop:variational}(b). Now Lemma~\ref{lem:perron-branch} applies to $Q$
and the irreducible $Q'$ and gives every principal minor at once, with no descent through
the faces.
\end{remark}

\subsection{The generators with given minors}\label{sec:fiber}

Without reversibility the exit rates and the edge products $q_{ij}q_{ji}$ do not determine
$Q$, since $Q^*$ has the same exit rates, the same products and the same stationary law.
The next theorem describes all generators with the principal minors of $-Q$.

Let $C$ be a Z-matrix, that is $-C$ is Metzler, and let $X$ be a cut of $C$ with
$Y=V\setminus X$. If $C$ is irreducible, the blocks $C[X,Y]$ and $C[Y,X]$ are nonzero with
nonpositive entries and rank one, so $C[X,Y]=-ab^\top$ and $C[Y,X]=-cu^\top$ with
$a,b,c,u\ge0$. The cut-transpose of $C$ across $X$ is
\begin{equation}\label{eq:cut-transpose}
 C=\begin{pmatrix}C_{XX}&-ab^\top\\-cu^\top&C_{YY}\end{pmatrix}\mapsto
 C'=\begin{pmatrix}C_{XX}&-ac^\top\\-bu^\top&C_{YY}^\top\end{pmatrix}.
\end{equation}
Another choice of the factors changes $C'$ by a positive diagonal similarity
$\operatorname{diag}(\alpha I_X,I_Y)$. Transposition commutes with this operation. The
transpose $C^\top$ has the same cuts, and with the factors $(u,c,b,a)$ in place of $(a,b,c,u)$
its cut-transpose across $X$ is $(C')^\top$, as~\eqref{eq:cut-transpose} shows by inspection.

\begin{definition}[generator points and the cut-transpose orbit]\label{def:orbit}
Let $Q$ be an irreducible generator and $C$ an irreducible Z-matrix with the principal minors of
$-Q$. A positive diagonal matrix $D$ such that $-D^{-1}CD$ is a generator exists and is unique up
to a factor (step (N) of the proof of Corollary~\ref{thm:nonreversible}), and $-D^{-1}CD$ is the
\emph{generator point} of $C$. The \emph{cut-transpose orbit} of $-Q$ is the set of matrices
obtained from $-Q$ or from $-Q^\top$ by finitely many cut-transposes, each across a cut of the
current matrix, taken up to positive diagonal similarity.
\end{definition}

If $Q$ has no cut, the orbit consists of the classes of $-Q$ and $-Q^\top$, whose generator
points are $Q$ and $Q^*$.

\begin{corollary}[the generator fiber]\label{thm:nonreversible}
Let $Q$ be an irreducible generator.
\begin{enumerate}[(a)]
\item Every generator $Q'$ with the principal minors of $-Q$ is irreducible and has the
spectrum of $Q$. The set of such generators is finite, and it contains $Q$ and its time
reversal $Q^*$.
\item If $Q$ has no cut, that is, for every partition $V=X\cup Y$ with $|X|,|Y|\ge2$ one of
the blocks $Q[X,Y]$, $Q[Y,X]$ has rank at least $2$, then the set in \textup{(a)} is
$\{Q,Q^*\}$. It is a single point exactly when $Q$ is reversible.
\item In general the set in \textup{(a)} is the set of generator points of the
cut-transpose orbit of $-Q$ (Definition~\ref{def:orbit}). Every member of the orbit is an
irreducible Z-matrix with the principal minors of $-Q$, and fewer than $2d$ cut-transposes
reach every point of the set. If $Q$ is reversible, the set is $\{Q\}$.
\item If $d\le3$, the set in \textup{(a)} is $\{Q,Q^*\}$ with no condition.
\end{enumerate}
\end{corollary}

Clause (c) needs the principal minors of $-Q$. For a general irreducible Z-matrix no
diagonal similarity need produce a generator. The matrix $3I-P$ with $P$ irreducible and
stochastic is nonsingular, while the negative of every generator has the eigenvalue $0$,
and similarity preserves the spectrum.

\emph{What is borrowed.} This is a corollary of the theory of the principal-minor map. The
matrix theorems used are Loewy~\cite[Theorem~1]{loewy1986}
for order $n\ge4$ without a cut, Hartfiel and Loewy~\cite[Theorem~3]{hartfielloewy1984} for
$n\le3$, combined in~\cite[Lemma~2.8]{cggr2024}, Al Ahmadieh~\cite[Theorem~4.1 and
Lemma~4.5]{alahmadieh2023} for the converse, and Chatterjee, Ghosh, Gurjar and
Raj~\cite[Theorem~1.1]{cggr2024} for the whole fiber. So (b), (d) and the reversible case of
(c) follow from them together with the normalization (N) below. What we add is (N), a short
translation to generators, and its consequences, the finiteness, the irreducibility of every
member, and the blockwise description of Proposition~\ref{prop:amplitudes}.

\begin{proof}
We first record 2 facts used in (a) to (d).

\emph{(N) The generator point.} Let $C$ be an irreducible Z-matrix with the principal
minors of $-Q$. The characteristic polynomial
$\det(tI-C)=\sum_{s=0}^d(-1)^st^{d-s}E_s(C)$, with $E_s$ the sum of the principal minors of
order $s$, is that of $-Q$. So $-C$ has the spectrum of $Q$, whose eigenvalue of largest real
part is the simple eigenvalue $0$. The theorem of Perron and Frobenius for the irreducible
Metzler matrix $-C$ gives a vector $x>0$, unique up to a factor, with $Cx=0$. With
$D=\operatorname{diag}x$ the matrix $-D^{-1}CD$ has off-diagonal entries $-c_{ij}x_j/x_i\ge0$
and zero row sums, so it is an irreducible generator. If $D'$ is real, diagonal and
invertible and $-D'^{-1}CD'$ has zero row sums, then $C(D'\ones)=0$, so $D'$ is a multiple of
$D$ and gives the same matrix. So the generator point of Definition~\ref{def:orbit} exists and is
unique. For $C=-Q$ it
is $Q$, and for $C=-Q^\top$, whose null vector is $\pi$, it is $Q^*$.

\emph{(R) Real similarities.} If $-Q'=E(-Q)E^{-1}$ or $-Q'=E(-Q)^\top E^{-1}$ for a complex
invertible diagonal $E$, then $E$ may be taken real. In the first case
$q'_{ij}=e_iq_{ij}/e_j$, so $e_i/e_j=q'_{ij}/q_{ij}$ is real on every edge of the graph of
$Q$. That graph is strongly connected, so every ratio $e_i/e_1$ is a product of such ratios
and is real, and $E/e_1$ is a real diagonal similarity. The second case is the same with
$q_{ji}$ in place of $q_{ij}$.

(a) Let $Q'$ be a generator with the principal minors of $-Q$, and suppose that $-Q'$ is
reducible. Then there is a proper nonempty $A\subset V$ with $Q'[V\setminus A,A]=0$, and
$\det(-Q')=\det(-Q'_A)\det(-Q'_{V\setminus A})$. Both factors are proper principal minors,
so they equal $\det(-Q_A)$ and $\det(-Q_{V\setminus A})$. By Lemma~\ref{lem:face-monotone}
every irreducible component of a proper face has a positive exit rate, so every eigenvalue
of $-Q_A$ has positive real part and $\det(-Q_A)>0$, and likewise
$\det(-Q_{V\setminus A})>0$. This contradicts $\det(-Q')=\det(-Q)=0$. So $Q'$ is irreducible,
and the same argument applies to every real matrix with the principal minors of $-Q$. This also
follows from Lemma~\ref{lem:RF-irreducible}, since such a matrix has the polynomial $R_V$ of
$-Q$. The spectra agree by (N). The time reversal $Q^*$ is a generator, since
$Q^*\ones=D_\pi^{-1}Q^\top\pi=0$ and its off-diagonal entries are $\pi_jq_{ji}/\pi_i\ge0$,
and $-Q^*=D_\pi^{-1}(-Q)^\top D_\pi$ has the principal minors of $-Q$, since transposition
and diagonal similarity preserve them. Finiteness follows from (c).

(b) For $d\le3$ the claim is (d). Let $d\ge4$. Loewy's theorem~\cite[Theorem~1]{loewy1986}
applies to the irreducible matrix $-Q$ without a cut, with no condition on $-Q'$. It gives
an invertible diagonal $E$ with $-Q'=E(-Q)E^{-1}$ or $-Q'^\top=E(-Q)E^{-1}$, and by (R) we may
take $E$ real. In the first case (N) with $C=-Q$ gives $Q'=Q$, and in the second case
$Q'=E^{-1}Q^\top E$ and (N) with $C=-Q^\top$ gives $Q'=Q^*$. Finally $Q^*=Q$ exactly when
$\pi_iq_{ij}=\pi_jq_{ji}$ for all $i,j$, that is when $Q$ is reversible.

(c) Let $C$ and $X$ be as in (c). The cut-transpose $C'$ is a Z-matrix. It has the same
principal minors as $C$. Indeed, for $U\subseteq V$ write $U_X=U\cap X$ and $U_Y=U\cap Y$. If
$U_X$ or $U_Y$ is empty, both minors agree since a determinant is invariant under
transposition. Otherwise replace $C_{XX}$ by $C_{XX}+sI$ with $s$ large, so that the block
of $U_X$ is invertible. The Schur complement gives
\[
 \det C_U=\det C_{U_X}\det\bigl(C_{U_Y}-\kappa\,c_{U_Y}b_{U_Y}^\top\bigr),\qquad
 \det C'_U=\det C_{U_X}\det\bigl(C_{U_Y}^\top-\kappa\,b_{U_Y}c_{U_Y}^\top\bigr),
\]
with the same scalar $\kappa=u_{U_X}^\top C_{U_X}^{-1}a_{U_X}$. The two right sides are
equal by transposition. Both sides are polynomials in $s$ that agree for large $s$, so they
agree at $s=0$. By the argument of (a), $C'$ is irreducible, and by (N) it has exactly one
generator point.

Now let $Q'$ be a generator in the fiber, irreducible by (a). The theorem of Chatterjee,
Ghosh, Gurjar and Raj~\cite[Theorem~1.1]{cggr2024} gives a sequence
$-Q=C_0,C_1,\dots,C_k$ with $k<2d$, each $C_{s+1}$ a cut-transpose of $C_s$ across a cut of
$C_s$~\cite[Definition~2.10]{cggr2024}, such that $-Q'$ is diagonally similar to $C_k$ or to
$C_k^\top$. Their cut-transpose uses a particular rank-one factorization, and another choice
changes it by a real diagonal similarity $\operatorname{diag}(\alpha I_X,I_Y)$~\cite[Remark~2.11]{cggr2024}.
A real invertible diagonal similarity of $C$ by $\operatorname{diag}(E_X,E_Y)$ changes the
cut-transpose by $\operatorname{diag}(E_X,E_Y^{-1})$, as a short computation with the blocks
shows, and does not change the cuts. So, up to real invertible diagonal similarity, which
(N) accepts, every $C_s$ is obtained by our cut-transposes, and $Q'$ is the generator point
of $C_k$ or of $C_k^\top$. Since transposition commutes with cut-transposes, $C_k^\top$ is
reached from $-Q^\top$ by the same sequence, so $-Q'$ lies in the orbit of
Definition~\ref{def:orbit}. Conversely every generator point of the orbit lies in the fiber, by
the first paragraph of (c) and the invariance of the minors under transposition. There are
fewer than $2^d$ cuts at each step and fewer than $2d$ steps, so the fiber is finite, which
completes (a).

If $Q$ is reversible, $S=D_{\sqrt\pi}(-Q)D_{\sqrt\pi}^{-1}$ is an irreducible symmetric
Z-matrix, and a symmetric matrix is its own cut-transpose. Indeed, if $S[X,Y]=-ab^\top$, then
$S[Y,X]=-ba^\top$, so we may take $c=b$ and $u=a$. By the rule for diagonal similarities
above, the cut-transpose of a matrix diagonally similar to $S$ is again diagonally similar
to $S$, and so is the transpose $E^{-1}SE$ of $ESE^{-1}$. Since $C_0=-Q$ is diagonally similar
to $S$, so is every $C_s$. Hence $-Q'$ is diagonally similar to $-Q$, and (N) with $C=-Q$
gives $Q'=Q$. See also Engel and Schneider~\cite[Theorem~3.5]{engelschneider1980} and Al
Ahmadieh~\cite[Lemma~6.1 and Theorem~6.2]{alahmadieh2023}.

(d) For $d=2$ a generator is fixed by $q_1$ and $q_2$, which are principal minors. Let
$d=3$ and let $Q'$ be in the fiber, so $Q'$ is irreducible by (a). The theorem of Hartfiel
and Loewy~\cite[Theorem~3]{hartfielloewy1984}, as restated in~\cite[Lemma~2.8]{cggr2024},
says that for an irreducible matrix $A$ of order $2$ or $3$ over a field, every matrix with
the principal minors of $A$ is diagonally similar to $A$ or to $A^\top$ over that field. So
$-Q'=E(-Q)E^{-1}$ or $-Q'=E(-Q)^\top E^{-1}$ with $E$ diagonal and invertible, real by (R),
and (N) gives $Q'=Q$ or $Q'=Q^*$, as in (b).
\end{proof}

A cut-transpose need not be $Q$ or $Q^*$. Take $d=4$, $X=\{1,2\}$ and
\[
 Q=\begin{pmatrix}-4.7&0.7&1&3\\1.3&-9.3&2&6\\2&2&-4.4&0.4\\1&1&2&-4\end{pmatrix},
\]
whose blocks $Q[X,Y]=(1,2)^\top(1,3)$ and $Q[Y,X]=(2,1)^\top(1,1)$ have rank one. To 4
decimals, the generator point of the cut-transpose of $-Q$ across $X$ is
\[
 Q'=\begin{pmatrix}-4.7&0.7&2.3810&1.6190\\1.3&-9.3&4.7619&3.2381\\
 0.84&0.84&-4.4&2.72\\1.8529&1.8529&0.2941&-4\end{pmatrix}.
\]
Its entries are rational, for instance $q'_{13}=50/21$, $q'_{41}=63/34$ and
$q'_{43}=5/17$. It has the 15 principal minors of $-Q$, with $Q'_X=Q_X$ and
$Q'_Y=Q^*_Y$, but $Q'\ne Q^*$, since $q^*_{12}=351/530\approx0.6623$.

\begin{remark}[three states by hand]\label{rem:three-states}
For $d=3$ the minors give the exit rates, the products $q_{ij}q_{ji}$, and the sum and
product of the two oriented $3$-cycle products $\omega=q_{12}q_{23}q_{31}$ and
$\omega'=q_{13}q_{32}q_{21}$. So $Q'$ has the cycle products of $Q$ or of $Q^*$, hence differs
from it by a potential, and (N) gives (d). Example~5.4 of Kim and Landim~\cite{kl2024} is the
case $\omega=1$, $\omega'=0$.
\end{remark}

\subsection{What the exact law adds}\label{sec:stationary}

The law of the counts is finer than $J_V$, and for a generic start it separates the fiber.
From the stationary start it cannot separate $Q$ from $Q^*$ at all. Section~\ref{sec:exactlaw}
determines what the exact law sees.

\begin{remark}[the stationary start]\label{rem:stationary}
Let $Q$ be irreducible with stationary law $\pi$. With the start $\pi$, the law of the
composition is the same under $Q$ and under $Q^*$, for every $N$ and every $T$ with
$Tq_{\max}\le N$, and so is the weighted law $P_N(n)e^{Tv\cdot n/N}$ for every field $v$. This is
path reversal, and we include the proof only to fix the discrete-time form. The matrix $P=I+(T/N)Q$ has stationary law $\pi$, and $P^*=D_\pi^{-1}P^\top D_\pi=I+(T/N)Q^*$
is stochastic, since $Q^*$ has the exit rates of $Q$. For a path $x_1\cdots x_N$ we have
$\pi_{x_1}\prod_tP(x_t,x_{t+1})=\pi_{x_N}\prod_tP^*(x_{t+1},x_t)$. Reversal is a bijection of
the paths that keeps the composition. Summing over the paths with composition $n$ gives the
claim, since the factor $e^{Tv\cdot n/N}$ depends only on $n$.
\end{remark}

For a face $F$ that is irreducible for $Q$ and a field $v$, the face maximum has the constant
$K_F(v)$ of Definition~\ref{def:tilt}. Let $r$ and $l$ be the right and left Perron vectors of
$Q_F+\operatorname{diag}v_F$ with $l\cdot r=1$. By Proposition~\ref{prop:constant},
\begin{equation}\label{eq:amplitude}
 K_F(v)=\frac{a_F(v)}{(2\pi)^{(k-1)/2}\sqrt{\det\Sigma_{\alpha^*_F(v)}}},\qquad
 a_F(v)=(\mu_F\cdot r)(l\cdot\ones),
\end{equation}
and we call $a_F(v)$ the amplitude.

\begin{proposition}[the amplitudes]\label{prop:amplitudes}
Let $Q$ be irreducible, let $Q'\ne Q$ be a generator in the fiber of
Corollary~\ref{thm:nonreversible}, and let $F$ be a face that is irreducible for $Q$.
\begin{enumerate}[(a)]
\item $F$ is irreducible for $Q'$, the two generators have the same $\lambda_F$,
$\alpha^*_F$, $J_F$ and $\det\Sigma_{\alpha^*_F(v)}$, and $K'_F(v)/K_F(v)=a'_F(v)/a_F(v)$.
\item There is a finite union $\mathcal Z$ of proper linear subspaces of $\R^V$ with the
following property. For every start $\mu$ of full support outside $\mathcal Z$ and every two
distinct generators $Q'$ and $Q''$ in the fiber, $K'_V(v)\ne K''_V(v)$ for all $v$ outside a
closed nowhere dense set of fields. That set always contains $v=0$, where every generator has $r=\ones$, $l=\pi$ and
$a_V(0)=1$. If $Q$ is not reversible, then $\pi\in\mathcal Z$.
\item If $Q'$ is the generator point of one cut-transpose of $-Q$ across $X$, then for every
start $K'_F=K_F$ on every face $F\subseteq X$ and $K'_F=K^*_F$ on every face $F\subseteq Y$,
where $K^*$ belongs to $Q^*$.
\end{enumerate}
\end{proposition}

\begin{proof}
(a) The first claims are Theorem~\ref{thm:dv}. So $\alpha^*_F(v)$ and, by
Proposition~\ref{prop:constant}, $\det\Sigma$ at that point agree for $Q$ and $Q'$, and
\eqref{eq:amplitude} gives the ratio.

(b) Take $F=V$ and put $g(v)=r(v)\,(l(v)\cdot\ones)$, which does not depend on the scale of
$r$, so that $a_V(v)=\mu\cdot g(v)$. The map $g$ is real-analytic on $\R^V$. Let $W$ be the
span of the vectors $g'(v)-g(v)$ over $v\in\R^V$. If $\mu\notin W^\perp$, the analytic
function $a'_V-a_V=\mu\cdot(g'-g)$ is not identically zero, so its zero set is closed and
nowhere dense. We show that $W\ne0$. Suppose $g'=g$. Then $r'(v)$ is a multiple of $r(v)$,
and both matrices have the eigenvalue $\lambda_V(v)$, so $(Q'-Q)r(v)=0$ for every $v$. Take
$v=te_j$ with $t\to\infty$. Then $\lambda_V\le q_j-t\to-\infty$, and row $i\ne j$ of the
eigenvalue equation gives $(q_i-\lambda_V)r_i\le q_{\max}\max r$, so $r(te_j)/\max r\to e_j$.
Hence $(Q'-Q)e_j=0$ for every $j$ and $Q'=Q$. So for $Q'\ne Q$ the space $W^\perp$ is proper.
Every member $Q''$ of the fiber is irreducible and has the same fiber as $Q$, by
Corollary~\ref{thm:nonreversible}(a), so the same argument with $Q''$ in place of $Q$ gives a
proper space $W^\perp_{Q'',Q'}$ for every ordered pair of distinct members. Let $\mathcal Z$ be
the union of these finitely many spaces.
At $v=0$ every generator has $r=\ones$ and $l=\pi$, so $a_V(0)=\mu(V)=1$. For the time
reversal, $-Q^*-\operatorname{diag}v=D_\pi^{-1}(-Q-\operatorname{diag}v)^\top D_\pi$, so
$r^*=D_\pi^{-1}l$ and $l^*=D_\pi r$, with $l^*\cdot r^*=1$. At $\mu=\pi$ this gives
$a^*_V(v)=(\ones\cdot l)(\pi\cdot r)=a_V(v)$ for all $v$, so $\pi\in\mathcal Z$ when
$Q^*\ne Q$, in line with Remark~\ref{rem:stationary}.

(c) Let $C=-Q$ and let $C'$ be its cut-transpose across $X$, with the factors of
\eqref{eq:cut-transpose}, and put $s_0=(b\cdot\ones)/(c\cdot\pi_Y)$. The rows $X$ of
$Q\ones=0$ give $-Q_X\ones=(b\cdot\ones)a$, the columns $Y$ of $\pi^\top Q=0$ give
$-Q_Y^\top\pi_Y=(\pi_X\cdot a)b$, and summing the columns $X$ of $\pi^\top Q=0$ gives
$(\pi_X\cdot a)(b\cdot\ones)=(\pi_Y\cdot c)(u\cdot\ones)$. With these identities,
$x=(\ones_X,s_0\pi_Y)$ satisfies $C'x=0$. So $Q'=-D_x^{-1}C'D_x$, and its killed blocks are
$Q'_X=Q_X$ and $Q'_Y=D_{\pi_Y}^{-1}Q_Y^\top D_{\pi_Y}=Q^*_Y$. The constant $K_F(v)$ of a face
depends only on $Q_F$ and $\mu_F$, which gives the claim.
\end{proof}

Proposition~\ref{prop:amplitudes} concerns the law, not the labels of
Section~\ref{sec:modal} (see the open questions in Section~\ref{sec:verify}).

% ---------------------------------------------------------------------
% RESERVED SECTION: the exact occupation law (theorem O1).
% \section{The exact occupation law}\label{sec:exactlaw}

This section answers the question of Section~\ref{sec:intro} at its first resolution, the exact
law, and proves Theorem~\ref{thm:law}. We work in continuous time and transfer the results to the
counts of $I+hQ$ at the end of Section~\ref{sec:law-reductions}. The law at one time contains the
principal minors (Theorem~\ref{thm:one-horizon}), so it determines everything that $J_V$
determines. It determines $Q$ from a generic start (Theorem~\ref{thm:generic-start}). From the
stationary start it determines $Q$ up to time reversal in the cases of
Theorem~\ref{thm:stationary-law}, and Section~\ref{sec:law-cuts} shows that it does not do so in
general.

\subsection{Reductions}\label{sec:law-reductions}

Let $(\xi_t)_{t\ge0}$ be the chain with generator $Q$ and start $\mu$, and let
$O(T)\in[0,T]^d$ be its vector of occupation times,
$O_i(T)=\int_0^T\ones\{\xi_t=i\}\,dt$. For $x\in\mathbb C^d$ put
\[
 F^\mu_Q(x)=\mu\bigl(\operatorname{diag}x-Q\bigr)^{-1}\ones ,\qquad F_Q=F^\pi_Q
 \quad\text{when $Q$ is irreducible,}
\]
a rational function of $x$. We write $W=\operatorname{diag}w$ for $w\in\R^d$, $D=D_\pi$ for the
diagonal matrix of $\pi$, and $\Sigma_m=DQ^m+(DQ^m)^\top$ for $m\ge0$.

\begin{lemma}[reductions]\label{lem:reductions}
Let $Q$ be a generator and $\mu$ a start.
\begin{enumerate}[(a)]
\item $\E_\mu e^{-w\cdot O(T)}=\mu e^{T(Q-W)}\ones$ for $w\in\R^d$. The laws of $O(T)$ for all
$T\ge0$ determine $F^\mu_Q$, and $F^\mu_Q$ determines them.
\item Let $h>0$ with $hq_{\max}\le1$, and let $n(N)$ be the composition of $X_1,\dots,X_N$, where
$X_1\sim\mu$ and the transition matrix is $I+hQ$. For small $|y|$ and nonzero $z\in\mathbb C^d$,
\[
 \sum_{N\ge1}y^N\,\E_\mu\prod_iz_i^{n_i(N)}=h^{-1}F^\mu_Q(x),\qquad x_i=\frac{1/(yz_i)-1}h .
\]
So at a fixed $h$ the laws of $n(N)$ for all $N\ge1$ determine $F^\mu_Q$, and $F^\mu_Q$
determines them.
\item If $Q$ is irreducible, then
\[
 \E_\pi\bigl[O(T)O(T)^\top\bigr]=\sum_{m\ge0}\frac{T^{m+2}}{(m+2)!}\,\Sigma_m .
\]
\end{enumerate}
\end{lemma}

\begin{proof}
(a) The function $u(T)=\E_\cdot\,e^{-w\cdot O(T)}$ solves $u'=(Q-W)u$ with $u(0)=\ones$, by the
Feynman and Kac formula, which for a finite chain is the backward equation. For real $s$ larger
than the real parts of the spectrum of $Q-W$,
$\int_0^\infty e^{-sT}\mu e^{T(Q-W)}\ones\,dT=\mu(sI-Q+W)^{-1}\ones=F^\mu_Q(s\ones+w)$. The law of
$O(T)$ lives on $[0,T]^d$, so it is determined by its Laplace transform on $\R^d$. Hence the laws
for all $T$ give the continuous functions $T\mapsto\mu e^{T(Q-W)}\ones$, these give $F^\mu_Q$ on
a nonempty open subset of $\R^d$, and a rational function is determined there. Conversely, for each
$w$ the function $s\mapsto F^\mu_Q(s\ones+w)$ for large $s$ is the Laplace transform of
$T\mapsto\mu e^{T(Q-W)}\ones$, which it determines.

(b) With $Z=\operatorname{diag}z$ and $P=I+hQ$ we have $\E_\mu\prod_iz_i^{n_i(N)}=\mu Z(PZ)^{N-1}\ones$.
Summing the geometric series gives $\mu\,yZ(I-yPZ)^{-1}\ones=\mu\bigl((yZ)^{-1}-P\bigr)^{-1}\ones$,
and $(yZ)^{-1}-I-hQ=h(\operatorname{diag}x-Q)$. The laws for all $N$ give the power series, hence
$F^\mu_Q$ on an open set of $x$, since $z\mapsto x$ is open for fixed $y$. Conversely $F^\mu_Q$ gives
the power series.

(c) By stationarity, with $P(r)=e^{rQ}=\sum_mr^mQ^m/m!$,
\[
 \E_\pi[O_i(T)O_j(T)]=\int_0^T(T-r)\bigl(\pi_iP_{ij}(r)+\pi_jP_{ji}(r)\bigr)\,dr,\qquad
 \int_0^T(T-r)\frac{r^m}{m!}\,dr=\frac{T^{m+2}}{(m+2)!}. \qedhere
\]
\end{proof}

By (a) and (b), each result of this section that assumes equal laws of $O(T)$ for all $T$, or at
one $T$ (which is the same by Theorem~\ref{thm:one-horizon}), holds verbatim for equal laws of
the counts of $I+hQ$ for all $N$ at one fixed $h$. Since $\Sigma_0=2D$, the second moments of
$O(T)$ to order $T^k$ are $\pi$ together with $\Sigma_1,\dots,\Sigma_{k-2}$.

\subsection{One time determines every time}\label{sec:law-one-time}

Put $\mathcal P(s,w)=\det(sI-Q+W)$. In the variables $(s,s\ones+w)$ it is the polynomial
$R_V(-(s\ones+w))$ of Section~\ref{sec:perron-branch}, which does not involve $s$. So if $Q$ is
irreducible, $\mathcal P$ is irreducible over $\mathbb C$ by Lemma~\ref{lem:RF-irreducible}, and it
is monic of degree $d$ in $s$.

\begin{lemma}[vanishing on a graph]\label{lem:graph}
Let $\mathcal P\in\mathbb C[s,w]$ be irreducible and monic in $s$ of degree $d\ge1$. Let
$U\subseteq\R^d$ be nonempty and open, and let $\theta$ be a function on $U$ with
$\mathcal P(\theta(w),w)=0$. If $\mathcal R\in\mathbb C[s,w]$ and $\mathcal R(\theta(w),w)=0$ for
all $w\in U$, then $\mathcal P$ divides $\mathcal R$.
\end{lemma}

\begin{proof}
Divide by the monic $\mathcal P$ in $s$, so $\mathcal R=\mathcal A\mathcal P+\mathcal B$ with
$\deg_s\mathcal B<d$ and $\mathcal B(\theta(w),w)=0$ on $U$. Suppose $\mathcal B\ne0$. Then
$\mathcal P$ does not divide $\mathcal B$, so by Gauss's lemma they have no common factor over
$\mathbb C(w)$, and the resultant $\rho=\operatorname{Res}_s(\mathcal P,\mathcal B)\in\mathbb C[w]$
is not zero. Since $\mathcal P$ is monic, $\rho(w)$ is the product of $\mathcal B(\sigma,w)$ over
the roots $\sigma$ of $\mathcal P(\cdot,w)$, and one of them is $\theta(w)$. So $\rho$ vanishes on
$U$, which is impossible for a nonzero polynomial.
\end{proof}

\begin{theorem}[one time]\label{thm:one-horizon}
Let $Q$ be an irreducible generator and $Q'$ any generator on $[d]$, with starts $\mu$ and $\mu'$.
If $O(T_0)$ has the same law under $(Q,\mu)$ and under $(Q',\mu')$ for one $T_0>0$, then $-Q'$
has the principal minors of $-Q$, so $Q'$ is irreducible, and $F^{\mu'}_{Q'}=F^\mu_Q$. So the laws
agree at every $T$.
\end{theorem}

The idea is to read the exponential of the leading eigenvalue of $Q-W$ off the Laplace
transform at one time, as $w$ grows along a ray. Its exponent is algebraic in $w$, and the
exponential of a nonconstant algebraic function is not algebraic. So the leading eigenvalues of
the 2 generators agree up to a constant, and the constant is $0$.

\begin{proof}
\emph{Step 1, the data.} By Lemma~\ref{lem:reductions}(a) the hypothesis says
$\mu e^{T_0(Q-W)}\ones=\mu'e^{T_0(Q'-W)}\ones$ for all $w\in\R^d$.

\emph{Step 2, one exponential.} Let $q=\max(q_{\max},q'_{\max})>0$, and for $v\in\R^d$ with
distinct entries let $\delta(v)$ be the least distance between 2 of them. By the theorem of
Gershgorin the eigenvalues of $V-\epsilon Q$, $V=\operatorname{diag}v$, lie within $2|\epsilon|q$
of the entries of $v$. So for $|\epsilon|<\epsilon_0(v)=\delta(v)/(4q)$ they are simple, given by
analytic functions $\kappa_i(\epsilon)$ with $\kappa_i(0)=v_i$, and so are their spectral
projectors $E_i(\epsilon)$. The same holds for $Q'$. Fix $v^0$ with distinct entries and $v^0_1$
the smallest, and a ball $\mathcal B$ around $v^0$ whose closure consists of such vectors. Then
$\epsilon_*=\min\epsilon_0>0$ on $\mathcal B$, because $\delta$ is continuous and positive on the
closure. This is what makes the set $U$ below open.

For $v\in\mathcal B$ and $t>1/\epsilon_*$ the eigenvalues of $Q-tV$ are $-t\kappa_i(1/t)$ and
\[
 e^{T_0tv_1}\mu e^{T_0(Q-tV)}\ones=\phi_v(1/t)+O(e^{-T_0t\delta(v)/2}),\qquad
 \phi_v(\epsilon)=\mu E_1(\epsilon)\ones\,\exp\Bigl(-T_0\,\frac{\kappa_1(\epsilon)-v_1}{\epsilon}\Bigr),
\]
where $\phi_v$ is analytic for $|\epsilon|<\epsilon_*$. With $\phi'_v$ formed from $Q'$ and
$\mu'$, Step 1 gives $\phi_v(\epsilon)-\phi'_v(\epsilon)=O(e^{-c/\epsilon})$ as $\epsilon\downarrow0$.
Every Taylor coefficient of $\phi_v-\phi'_v$ at $0$ is therefore $0$, and $\phi_v=\phi'_v$. On the
open cone $U=\{tv\mid v\in\mathcal B,\ t>1/\epsilon_*\}$ let $\theta(w)$ be the eigenvalue of $Q-W$
within $2q$ of $-w_1$, which is simple and real since its conjugate is also an eigenvalue there,
and let $\chi(w)=\mu E(w)\ones$ with $E(w)$ its spectral projector. Define $\theta'$ and $\chi'$
from $Q'$ and $\mu'$. We have shown that
\begin{equation}\label{eq:one-exponential}
 \chi(w)\,e^{T_0\theta(w)}=\chi'(w)\,e^{T_0\theta'(w)},\qquad w\in U .
\end{equation}

\emph{Step 3, the residue.} Write $\mu(sI-Q+W)^{-1}\ones=\mathcal N_\mu/\mathcal P$ with
$\mathcal N_\mu=\mu\operatorname{adj}(sI-Q+W)\ones$. Then
$\chi=\mathcal N_\mu(\theta,w)/\partial_s\mathcal P(\theta,w)$, the residue at the simple pole
$s=\theta(w)$. The polynomial $\mathcal N_\mu$ has degree $d-1$ in $s$ and leading coefficient
$\mu\cdot\ones=1$, so $\mathcal P$ does not divide it, and by Lemma~\ref{lem:graph}
$\mathcal N_\mu(\theta(w),w)$ does not vanish identically on any open subset of $U$. So there is a
ball $U_1\subseteq U$ on which $\chi\ne0$, and then $\chi'\ne0$ by~\eqref{eq:one-exponential}.

\emph{Step 4, transcendence.} Fix $w^0\in U_1$ and $\zeta\in\R^d$, and put
$f(r)=T_0(\theta'-\theta)(w^0+r\zeta)$ and $g(r)=(\chi/\chi')(w^0+r\zeta)$ for real $r$ near $0$.
Both are real-analytic, and both are algebraic over $\mathbb C(r)$, since $\theta$ and $\theta'$
are roots of polynomials in $s$ with coefficients in $\mathbb C[r]$, and $\chi$, $\chi'$ are
rational in them and $r$. By~\eqref{eq:one-exponential}, $e^f=g$. We claim that $f$ is constant.
The field $\mathbb C(r,f,g)$ is the field of meromorphic functions of a compact Riemann surface,
on which $f$ and $g$ are meromorphic, and $e^f=g$ continues to the surface minus the poles of $f$
and $g$. If $f$ were not constant it would have a pole $p$, with $f=a\tau^{-k}(1+O(\tau))$, $a\ne0$,
$k\ge1$, in a local coordinate $\tau$ at $p$. Along the directions where $a\tau^{-k}>0$ the
function $|e^f|$ grows faster than every power of $1/|\tau|$, while $g$ is meromorphic at $p$.
So $f$ is constant. As $\zeta$ is arbitrary, the gradient of $\theta'-\theta$ vanishes on $U_1$,
and $\theta'=\theta+c$ there for a real constant $c$.

\emph{Step 5, the minors.} Put $\mathcal P'(s,w)=\det(sI-Q'+W)$. The polynomial
$\mathcal P'(s+c,w)$ vanishes at $s=\theta(w)$ for every $w\in U_1$. This holds only for real $w$,
and Lemma~\ref{lem:graph} turns it into an identity of polynomials. It gives that $\mathcal P$
divides $\mathcal P'(s+c,w)$, and both are monic of degree $d$ in $s$, so they are equal. At $s=0$
this says that $-Q'+cI$ and $-Q$ have the same principal minors, and at $w=0$ that the spectrum
of $Q'$ is that of $Q$ shifted by $c$. Every eigenvalue of a generator has real part at most $0$.
Since $0$ is an eigenvalue of $Q'$, $-c$ is one of $Q$, so $c\ge0$, and since $0$ is an eigenvalue
of $Q$, $c$ is one of $Q'$, so $c\le0$. Hence $c=0$, $-Q'$ has the principal minors of $-Q$, and
$Q'$ is irreducible by Corollary~\ref{thm:nonreversible}(a).

\emph{Step 6, the transforms.} Now $\mathcal P'=\mathcal P$, $\theta'=\theta$, and $\chi=\chi'$ on
$U_1$ by~\eqref{eq:one-exponential}. So $\mathcal N_\mu-\mathcal N'_{\mu'}$ vanishes at
$s=\theta(w)$ on $U_1$, and by Lemma~\ref{lem:graph} it is divisible by $\mathcal P$. Its degree
in $s$ is less than $d$, so it is $0$. Hence $F^\mu_Q(s\ones+w)=F^{\mu'}_{Q'}(s\ones+w)$, and
Lemma~\ref{lem:reductions}(a) gives the laws at every $T$.
\end{proof}

\subsection{A generic start}\label{sec:law-generic}

For a generator $Q$ put $v_Q(x)=(\operatorname{diag}x-Q)^{-1}\ones$, a vector of rational
functions, so that $F^\mu_Q=\mu\cdot v_Q$. Let $\mathrm{Fib}(Q)$ be the finite set of generators with
the principal minors of $-Q$ (Corollary~\ref{thm:nonreversible}).

\begin{theorem}[a generic start]\label{thm:generic-start}
Let $Q$ be irreducible. For $Q'\in\mathrm{Fib}(Q)\setminus\{Q\}$ the set
$Z_{Q'}=\{\mu\in\R^d\mid\mu\cdot(v_{Q'}-v_Q)=0\}$ is a proper linear subspace. Let $Z$ be their
union. If $\mu$ is a start outside $Z$ and $Q'$ is a generator such that $O(T_0)$ has the same law
under $(Q,\mu)$ and $(Q',\mu)$ for one $T_0>0$, then $Q'=Q$. The set $Z$ meets the simplex in a
closed nowhere dense set of measure zero, and $\mu$ need not have full support.
\end{theorem}

\begin{proof}
With $m_i=1/x_i$, the expansion $(\operatorname{diag}x-Q)^{-1}=\sum_{k\ge0}(\operatorname{diag}m\,Q)^k\operatorname{diag}m$
for large $x$ gives $v_Q=m+\operatorname{diag}m\,Qm+O(|m|^3)$. The coefficient of $m_im_j$ in the
component $i$ is $q_{ij}$, also for $j=i$. So $v_Q$ determines $Q$, $v_{Q'}\ne v_Q$, and $Z_{Q'}$ is
proper. Under the hypothesis, Theorem~\ref{thm:one-horizon} puts $Q'$ in $\mathrm{Fib}(Q)$ and
gives $\mu\cdot v_{Q'}=\mu\cdot v_Q$, so $Q'=Q$ or $\mu\in Z_{Q'}$. A proper subspace cannot
contain the hyperplane $\{\mu\cdot\ones=1\}$, which spans $\R^d$. So each $Z_{Q'}$ meets it in a
proper affine subspace or not at all, and there are finitely many.
\end{proof}

The stationary law lies in $Z_{Q^*}$ when $Q^*\ne Q$, by Theorem~\ref{thm:stationary-law} below.
Proposition~\ref{prop:amplitudes}(b) separates the fiber by the constants of the local law, and
Theorem~\ref{thm:generic-start} by the exact law at one time.

\subsection{Second moments}\label{sec:law-moments}

Let $Q$ be irreducible, let $e$ have the entries $e_i=\pi_i^{1/2}$, and put $G=D^{1/2}QD^{-1/2}$, so that $Ge=0=e^\top G$ and
the time reversal $Q^*$ corresponds to $G^\top$. Write $S=(G+G^\top)/2$ and $K=(G-G^\top)/2$,
so $Ke=0$. Then $\Sigma_m=2D^{1/2}\operatorname{sym}(G^m)D^{1/2}$, where
$\operatorname{sym}A=(A+A^\top)/2$. A word in $S$ and $K$ is symmetric or antisymmetric according
to the parity of the number of letters $K$. So
\begin{equation}\label{eq:sym-powers}
 \operatorname{sym}G=S,\qquad\operatorname{sym}G^2=S^2+K^2,\qquad
 \operatorname{sym}G^3=S^3+SK^2+K^2S+KSK .
\end{equation}
The matrix $K^2$ is symmetric and negative semidefinite. If a nonzero eigenvalue $-\sigma^2$ of
$K^2$ has multiplicity $2$, its eigenplane $E$ is invariant under $K$, which acts on it as $\pm\sigma$
times a rotation by a right angle. Let $P_E$ be the orthogonal projection on $E$.

\begin{theorem}[second moments]\label{thm:second-moments}
Let $Q$ be irreducible with stationary law $\pi$, and let $Q'$ be a generator with $\pi Q'=0$.
\begin{enumerate}[(i)]
\item If $d\le4$ and $\Sigma'_m=\Sigma_m$ for $m=1,2$, then $Q'$ is $Q$ or $Q^*$.
\item Suppose that every nonzero eigenvalue of $K^2$ has multiplicity exactly $2$, and that the
graph on the eigenplanes, with an edge between $E$ and $E'$ when $P_ESP_{E'}\ne0$, is connected.
If $\Sigma'_m=\Sigma_m$ for $m=1,2,3$, then $Q'$ is $Q$ or $Q^*$.
\item The hypothesis of \textup{(ii)} holds for all generators with positive rates outside a
relatively closed nowhere dense set of measure zero. For each $\pi$ the same holds inside the
generators with positive rates and stationary law $\pi$.
\end{enumerate}
\end{theorem}

\begin{proof}
Define $G'$, $S'$ and $K'$ from $Q'$ with the same $D$, so $G'e=0=e^\top G'$ and $K'e=0$. By
\eqref{eq:sym-powers} the data give $S'=S$ and $K'^2=K^2$, and in (ii) also $K'SK'=KSK$. For real
antisymmetric $K'$ and every $y$ we have $|K'y|^2=-y^\top K'^2y=-y^\top K^2y$, so $K'$ vanishes on
the kernel of $K^2$. It commutes with $K'^2=K^2$, so it maps each eigenspace of $K^2$ to itself.

(i) Both $K$ and $K'$ act on $e^\perp$, of dimension $d-1\le3$, so they have rank at most $2$. If
$K=0$, then $K'=0$. Otherwise $K^2$ has one nonzero eigenvalue $-\sigma^2$ with an eigenplane $E$,
and on $E$ both $K$ and $K'$ are $\pm\sigma$ times the same rotation, since an
antisymmetric map of a plane with square $-\sigma^2I$ is $\pm\sigma$ times the rotation by a right
angle. So $K'=\pm K$.

(ii) By the same argument $K'=\epsilon_EK$ on each eigenplane $E$, with $\epsilon_E=\pm1$. Then
$P_EK'SK'P_{E'}=\epsilon_E\epsilon_{E'}\,KP_ESP_{E'}K$ and $P_EKSKP_{E'}=KP_ESP_{E'}K$. Since $K$ is
invertible on $E$ and on $E'$, $P_ESP_{E'}\ne0$ forces $\epsilon_E=\epsilon_{E'}$. By connectedness
$K'=\pm K$.

In both cases $G'=S\pm K\in\{G,G^\top\}$, and $Q'=D^{-1/2}G'D^{1/2}$ is $Q$ or
$D^{-1}Q^\top D=Q^*$.

(iii) Let $r=\lfloor(d-1)/2\rfloor$. Since $K$ is antisymmetric with $Ke=0$,
$\det(tI-K)=t^{d-2r}\psi(t^2)$ with $\psi$ monic of degree $r$, whose coefficients are polynomials
in $G$. If the discriminant of $\psi$ and $\psi(0)$ do not vanish, the roots $u_1,\dots,u_r$ of
$\psi$ are distinct and negative, and they are the nonzero eigenvalues of $K^2$, each with an
eigenplane $E_k$ and the projection
$P_k=(K^2/u_k)\prod_{j\ne k}(K^2-u_j)/(u_k-u_j)$. The product
$\prod_{k<l}\operatorname{tr}(P_kSP_lS)$, multiplied by $\bigl(\prod_ku_k\psi'(u_k)\bigr)^{r-1}$, is
a symmetric polynomial in $u_1,\dots,u_r$ with coefficients polynomial in $G$, hence a polynomial
in $G$. Where it and the first 2 do not vanish, every 2 eigenplanes are coupled, since
$\operatorname{tr}(P_kSP_lS)=\|P_kSP_l\|^2$. So the hypothesis of (ii) fails only on the zero set
of a polynomial $\Phi$ in $G$. On irreducible generators $\pi$ is a positive rational function
of $Q$, so $Q\mapsto\Phi(G)$ is real-analytic on the convex open set of generators with positive
rates. The zero set of a real-analytic function on a connected open set that is not identically
zero is closed, nowhere dense and of measure zero, and inside the slice of a fixed $\pi$ the map
$Q\mapsto G$ is linear. So it suffices to give, for each $\pi$, one generator with positive rates
and stationary law $\pi$ with $\Phi\ne0$.

Let $f_1,\dots,f_{d-1}$ be an orthonormal basis of $e^\perp$, put $u=\sum_if_i$, and take
\[
 G=-(I-ee^\top)+\epsilon\,uu^\top+\delta\sum_{k=1}^r\sigma_k\bigl(f_{2k}f_{2k-1}^\top-f_{2k-1}f_{2k}^\top\bigr)
\]
with distinct $\sigma_k>0$ and small $\epsilon,\delta>0$. Then $Ge=0=e^\top G$, the off-diagonal
entries are close to $e_ie_j>0$, and $Q=D^{-1/2}GD^{1/2}$ is a generator with positive rates and
stationary law $\pi$. The eigenplanes of $K^2$ are spanned by $f_{2k-1},f_{2k}$, with the distinct
eigenvalues $-\delta^2\sigma_k^2$. The first term is $-I$ on $e^\perp$, so for $k\ne l$
$P_kSP_l=\epsilon(f_{2k-1}+f_{2k})(f_{2l-1}+f_{2l})^\top\ne0$.
\end{proof}

\subsection{The stationary start}\label{sec:law-stationary}

\begin{theorem}[the stationary start]\label{thm:stationary-law}
Let $Q$ and $Q'$ be irreducible, and let $O(T_0)$ have the same law under $Q$ and $Q'$, each
started from its stationary law, for one $T_0>0$. Then $Q'$ is $Q$ or $Q^*$ in each of the
following cases.
\begin{enumerate}[(i)]
\item $d\le4$.
\item $Q$ has no cut.
\item $Q$ satisfies the hypothesis of Theorem~\ref{thm:second-moments}(ii).
\end{enumerate}
In cases (i) and (iii) the first and second moments of $O(T)$ to order $T^4$, respectively $T^5$,
suffice. Conversely, $O(T)$ has the same stationary law under $Q^*$ as under $Q$ for every $T$.
\end{theorem}

\begin{proof}
By Theorem~\ref{thm:one-horizon} the laws agree at every $T$. The first moments $T\pi$ and $T\pi'$
give $\pi'=\pi$, and Lemma~\ref{lem:reductions}(c) gives $\Sigma'_m=\Sigma_m$ for every $m$. Cases
(i) and (iii) are Theorem~\ref{thm:second-moments}, which uses only the moments named. In case
(ii) Theorem~\ref{thm:one-horizon} gives the principal minors, and
Corollary~\ref{thm:nonreversible}(b) gives $Q'\in\{Q,Q^*\}$. For the converse,
$F_{Q^*}(x)=\pi D^{-1}(\operatorname{diag}x-Q^\top)^{-1}D\ones=\ones^\top(\operatorname{diag}x-Q^\top)^{-1}\pi^\top=F_Q(x)$.
\end{proof}

The four-state cut-transpose $Q'$ after Corollary~\ref{thm:nonreversible} has the principal minors
and the stationary law of $Q$, and it is neither $Q$ nor $Q^*$. So by
Theorem~\ref{thm:stationary-law}(i) the stationary laws of $O(T)$ differ. In fact
$\pi(Q-W)^n\ones$ agrees for $n=1,2$, and
\[
 \pi(Q-W)^3\ones-\pi(Q'-W)^3\ones=\frac{29(w_1-w_2)(w_3-w_4)}{9673}
 =\frac12\,w^\top(\Sigma_1-\Sigma'_1)w ,
\]
so the symmetrized fluxes already differ.

\begin{remark}[one fixed $N$]\label{rem:one-N}
For the counts of $I+hQ$ from $\pi$ at one fixed $N$, the same computation as in
Lemma~\ref{lem:reductions}(c), with
\[
 P^r=\sum_k\binom rkh^kQ^k\qquad\text{and}\qquad\sum_{r}(N-r)\binom rk=\binom{N+1}{k+2},
\]
gives
\[
 \E_\pi[nn^\top]=N^2D+\sum_{k=1}^{N-1}h^k\binom{N+1}{k+2}\Sigma_k .
\]
The term $N^2D$ collects the diagonal $ND$ and the $k=0$ terms $N(N-1)D$. With $\E_\pi n=N\pi$,
the laws at one $N$ for $N-1$ distinct values of $h$ give $\pi$ and $\Sigma_1,\dots,\Sigma_{N-1}$.
So for $N\ge3$ they determine $Q$ up to time reversal when $d\le4$, and for $N\ge4$ under the
hypothesis of Theorem~\ref{thm:second-moments}(ii). One $N$ is not enough in general. At $N=2$
the law is $\pi_i(1-hq_i)$ on $2e_i$ and $h(\pi_iq_{ij}+\pi_jq_{ji})$ on $e_i+e_j$. On 3 states
with the uniform stationary law, $Q_k=S+kC$ with $S$ symmetric with off-diagonal entries at least
$1$ and $C$ the unit circulation on the 3-cycle is a generator for $|k|<1$, and its law at $N=2$
does not depend on $k$ for any $h$. Since $Q_k^*=Q_{-k}$, 2 of them with different $|k|$ are
neither equal nor time reversals of each other. We do not know what one pair $(N,T)$
determines.
\end{remark}

\subsection{Cuts, balanced sides and memoryless blocks}\label{sec:law-cuts}

Let $Q$ be irreducible with a cut $[d]=X\cup Y$, and write $Q[X,Y]=ab^\top$ and $Q[Y,X]=cu^\top$
with $a,b,c,u\ge0$, as in~\eqref{eq:cut-transpose} with $C=-Q$. Replacing $(a,b)$ by $(ta,b/t)$
we may assume $\pi_X\cdot a=u\cdot\ones$. Let $D_Y$ be the diagonal matrix of $\pi_Y$, let
$R_X=(\operatorname{diag}x_X-Q_X)^{-1}$ and $R_Y=(\operatorname{diag}x_Y-Q_Y)^{-1}$, and put
\[
\begin{aligned}
 p_1&=\pi_XR_Xa, & p_2&=u^\top R_X\ones, & \alpha&=u^\top R_Xa, & E_X&=\pi_XR_X\ones,\\
 q_1&=b^\top R_Y\ones, & q_2&=\pi_YR_Yc, & \beta&=b^\top R_Yc, & E_Y&=\pi_YR_Y\ones .
\end{aligned}
\]

\begin{proposition}[gluing across a cut]\label{prop:gluing}
In this setting $b\cdot\ones=\pi_Y\cdot c$. The generator point $Q'$ of the cut-transpose of $-Q$
across $X$ is given by $Q'_X=Q_X$, $Q'[X,Y]=a(D_Yc)^\top$, $Q'[Y,X]=D_Y^{-1}bu^\top$ and
$Q'_Y=Q^*_Y$, and it has the stationary law $\pi$. Moreover
\[
 F_Q=E_X+E_Y+\frac{p_1q_1+p_2q_2+\alpha q_1q_2+\beta p_1p_2}{1-\alpha\beta},\qquad
 F_Q-F_{Q'}=\frac{(p_1-p_2)(q_1-q_2)}{1-\alpha\beta}.
\]
\end{proposition}

\begin{proof}
The rows $X$ and $Y$ of $Q\ones=0$ give $Q_X\ones=-(b\cdot\ones)a$ and $Q_Y\ones=-(u\cdot\ones)c$.
The columns $X$ and $Y$ of $\pi Q=0$ give $\pi_XQ_X=-(\pi_Y\cdot c)u^\top$ and
$\pi_YQ_Y=-(\pi_X\cdot a)b^\top$. Applying the first and third to $\pi_X$ and $\ones$ gives
$(b\cdot\ones)(\pi_X\cdot a)=(\pi_Y\cdot c)(u\cdot\ones)$, hence $b\cdot\ones=\pi_Y\cdot c$. With
these 4 identities the rows of $Q'$ sum to $0$ and $\pi Q'=0$. For instance the rows $Y$ give
$D_Y^{-1}\bigl((u\cdot\ones)b+Q_Y^\top\pi_Y^\top\bigr)=0$. The off-diagonal entries of $Q'$ are
nonnegative, and $-Q'=E^{-1}C'E$ for the cut-transpose $C'$ of~\eqref{eq:cut-transpose} and
$E=\operatorname{diag}(I_X,D_Y)$. So $Q'$ is the generator point of $C'$.

For $x$ with large entries, $\operatorname{diag}x-Q=K_0-UV^\top$ with
$K_0=\operatorname{diag}(R_X^{-1},R_Y^{-1})$ and the $d\times2$ matrices
\[
 U=\begin{pmatrix}a&0\\0&c\end{pmatrix},\qquad V=\begin{pmatrix}0&u\\b&0\end{pmatrix},\qquad
 \text{so that}\qquad V^\top K_0^{-1}U=\begin{pmatrix}0&\beta\\\alpha&0\end{pmatrix}.
\]
Also $\pi K_0^{-1}U=(p_1,q_2)$ and $V^\top K_0^{-1}\ones=(q_1,p_2)^\top$, and the formula of Sherman,
Morrison and Woodbury gives the first identity. For $Q'$ the factors are $(a,D_Yc,D_Y^{-1}b,u)$,
which satisfy the same normalization, and $(\operatorname{diag}x_Y-Q^*_Y)^{-1}=D_Y^{-1}R_Y^\top D_Y$.
So $Q'$ has the same $p_1,p_2,\alpha,E_X$, the same $\beta$ and $E_Y$, and $q_1$ and $q_2$
exchanged. Subtracting gives the second identity.
\end{proof}

We call the side $X$ \emph{balanced} if $p_1=p_2$ as rational functions. This does not depend on
the choice of the factors, since the normalization ties the scale of $a$ to that of $u$.

\begin{corollary}[the balance criterion]\label{cor:balance}
The generator point of the cut-transpose of $-Q$ across $X$ gives the same stationary law of
$O(T)$ as $Q$, for one $T>0$ or equivalently for all $T$, if and only if $X$ or $Y$ is balanced.
\end{corollary}

\begin{proof}
Rational functions form an integral domain, and $1-\alpha\beta\to1$ as $x\to\infty$. Apply
Theorem~\ref{thm:one-horizon} and Lemma~\ref{lem:reductions}(a).
\end{proof}

\begin{remark}[necessary conditions]\label{rem:balance-necessary}
Let $X$ be balanced. Expanding $p_1-p_2$ in $m_i=1/x_i$ as in the proof of
Theorem~\ref{thm:generic-start}, the first order gives $u_i=\pi_ia_i$ for $i\in X$, so the jumps
into $X$ land in proportion to $\pi_ia_i$. The second order then gives
$(a_j-a_i)(\pi_iq_{ij}-\pi_jq_{ji})=0$ for $i\ne j$ in $X$. So an edge of $X$ that is not in
detailed balance joins 2 states with the same rate of leaving $X$. If $|X|=2$, either $a_1\ne a_2$
and the edge is in detailed balance, or $a_1=a_2$, and then $Q_X+(b\cdot\ones)a_1I$ is a
generator on 2 states with stationary law proportional to $\pi_X$, hence reversible. In both cases
$Q^*_X=Q_X$.
\end{remark}

\begin{definition}[memoryless block]\label{def:memoryless}
A set $B$ of states with $2\le|B|\le d-1$ is a \emph{memoryless block} of the irreducible $Q$ if
there are $\kappa>0$, a law $\nu$ on $B^c$ and $c\in[0,\infty)^{B^c}$ such that
$q_{iy}=\kappa\nu_y$ and $q_{yi}=c_y\pi_i/\pi(B)$ for all $i\in B$ and $y\notin B$. That is, every
state of $B$ leaves $B$ at the same rate and with the same landing law, and every jump into $B$
lands with the law $\pi_B/\pi(B)$. We write $Q^{[B]}$ for $Q$ with the block $Q_B$ replaced by
$Q^*_B=D_B^{-1}Q_B^\top D_B$, where $D_B$ is the diagonal matrix of $\pi_B$, and the other entries
kept.
\end{definition}

\begin{theorem}[reversal of a memoryless block]\label{thm:memoryless}
Let $B$ be a memoryless block of the irreducible $Q$. Then $Q^{[B]}$ is an irreducible generator
with stationary law $\pi$, and $F_{Q^{[B]}}=F_Q$. So $O(T)$ has the same stationary law under $Q$
and $Q^{[B]}$ for every $T$, and so do the counts of $I+hQ$ and $I+hQ^{[B]}$ for every $N$ and every
admissible $h$. Moreover $Q^{[B]}=Q$ exactly when $Q^*_B=Q_B$, which holds when $|B|=2$, and
$Q^{[B]}=Q^*$ when $|B^c|=1$.
\end{theorem}

\begin{proof}
Put $\rho=\pi_B/\pi(B)$, so $Q[B,B^c]=\kappa\ones\nu^\top$ and $Q[B^c,B]=c\rho$. The flux out of
$B$ equals the flux in, so $\kappa\pi(B)=\sum_y\pi_yc_y$. The columns $B$ of $\pi Q=0$ then give
$\pi_BQ_B=-\kappa\pi_B$, and the rows give $Q_B\ones=-\kappa\ones$. Hence
$Q^*_B\ones=D_B^{-1}Q_B^\top\pi_B^\top=-\kappa\ones$, so the rows of $Q^{[B]}$ sum to $0$, and
$\pi_BQ^*_B=\ones^\top Q_B^\top D_B=-\kappa\pi_B$, which the inflow $\kappa\pi_B$ cancels in the
columns $B$ of $\pi Q^{[B]}$. The other columns do not change. The off-diagonal entries
$\pi_jq_{ji}/\pi_i$ of $Q^*_B$ are nonnegative. The generator $Q^{[B]}$ is irreducible. It has the
same jumps as $Q$ out of $B$, into $B$ and inside $B^c$, and $\kappa>0$ and $c\ne0$ since $Q$ is
irreducible. Take states $s$ and $t$. If $s\in B$, jump out of $B$ first, which is possible from
every state of $B$. A path of $Q$ between two states of $B^c$ stays a path of $Q^{[B]}$ when each of
its sojourns in $B$ is cut to its entry state, since the jumps out of $B$ have the same rates from
every state of $B$. So we reach every state of $B^c$. If $t\in B$, go to a state $y$ with $c_y>0$
and jump to $t$, which is possible since $\rho_t>0$.

Solve $(\operatorname{diag}x-Q)z=\ones$ for $x$ with large entries, and put
$g=\rho(\operatorname{diag}x_B-Q_B)^{-1}\ones$. The rows $B$ give
$z_B=(1+\kappa\,\nu\cdot z_{B^c})(\operatorname{diag}x_B-Q_B)^{-1}\ones$, so
$\rho z_B=(1+\kappa\,\nu\cdot z_{B^c})g$. The rows $B^c$ give
$(\operatorname{diag}x_{B^c}-Q_{B^c}-\kappa g\,c\nu^\top)z_{B^c}=\ones+gc$, and
$F_Q=\pi(B)(1+\kappa\,\nu\cdot z_{B^c})g+\pi_{B^c}z_{B^c}$. So $F_Q$ depends on $Q_B$ only through
$g$, and on the other data $\pi$, $\kappa$, $\nu$, $c$ and $Q_{B^c}$, which $Q^{[B]}$ shares. For
$Q^{[B]}$ the function is
$\rho D_B^{-1}(\operatorname{diag}x_B-Q_B^\top)^{-1}D_B\ones=\pi_B(\operatorname{diag}x_B-Q_B)^{-1}\ones/\pi(B)=g$.
Hence $F_{Q^{[B]}}=F_Q$. The statement on the counts is Lemma~\ref{lem:reductions}(b).

If $|B|=2$, then $Q_B+\kappa I$ is a generator on 2 states with stationary law $\rho$, hence
reversible. If $B^c=\{y\}$, then $q^*_{iy}=\pi_yc_y/\pi(B)=\kappa=q_{iy}$ and
$q^*_{yi}=\pi_i\kappa/\pi_y=c_y\pi_i/\pi(B)=q_{yi}$, so $Q^{[B]}=Q^*$.
\end{proof}

The proof is algebraic, and the probability behind it is simple. The sojourns in $B$ have
exponential lengths of rate $\kappa$ and exit laws that do not depend on what happens inside
them, and each starts in the law $\rho$, which is stationary for the chain $Q_B+\kappa I$ inside.
Reversing a stationary path keeps its occupation times, and the moves into and out of $B$ do not
see the state. The path law does change.

\begin{remark}[memoryless blocks are balanced sides]\label{rem:memoryless-cut}
If $|B^c|\ge2$, then $B\cup B^c$ is a cut. With $X=B$ the factors of
Proposition~\ref{prop:gluing} are $a=\ones$, $b=\kappa\nu$, the vector $c/\pi(B)$ in the role of
$c$, and $u=\pi_B^\top$. They satisfy $\pi_B\cdot a=u\cdot\ones$, and then
$p_1=p_2=\pi_BR_B\ones$. So $B$ is a balanced side. By Proposition~\ref{prop:gluing}, the
generator point $Q''$ of the cut-transpose across $B$ has $Q''_B=Q_B$ and $Q''_{B^c}=Q^*_{B^c}$,
and a direct computation from its blocks gives $(Q'')^*=Q^{[B]}$. So
Theorem~\ref{thm:memoryless} is a case of Corollary~\ref{cor:balance}, and $Q^{[B]}$ is a generator
point of the cut-transpose orbit of $-Q$.
\end{remark}

\begin{example}[five states]\label{ex:five-state}
Let $d=5$ and $B=\{1,2,3\}$. Inside $B$ put $q_{12}=q_{23}=q_{31}=2$ and $q_{21}=q_{32}=q_{13}=1$.
Each state of $B$ jumps to $4$ at rate $1$ and to $5$ at rate $2$. State $4$ jumps to each state
of $B$ at rate $1$ and to $5$ at rate $1$, and state $5$ jumps to each state of $B$ at rate $2$
and to $4$ at rate $4$. The balance equations give $\pi=(4,4,4,6,3)/21$. So $B$ is a memoryless
block with $\kappa=3$, $\nu=(1/3,2/3)$ and $c=(3,6)$, and since $\pi_B$ is uniform, $Q^{[B]}$ is $Q$
with the block $Q_B$ transposed. It differs from $Q$, since $Q_B$ is not reversible, and from
$Q^*$, since $q^{[B]}_{14}=1$ while $q^*_{14}=\pi_4q_{41}/\pi_1=3/2$. By
Theorem~\ref{thm:memoryless} the stationary laws of $O(T)$ agree for every $T$, and so do the laws
of the counts of $I+hQ$ and $I+hQ^{[B]}$ from $\pi$ for every $N$ and every $h\le1/10$. From the
start $e_1$ the transforms differ. The hypothesis of Theorem~\ref{thm:second-moments}(ii) fails,
as it must, since the 2 generators have the same $\Sigma_m$ for every $m$. Here $K^2$ has the
eigenvalues $-7/8$ and $-3/4$, each twice, and $S$ does not couple the 2 eigenplanes.
\end{example}

By Theorem~\ref{thm:stationary-law}(i) there is no such example with $d\le4$. The stationary law
of the occupation times therefore determines $Q$ up to time reversal in the cases of
Theorem~\ref{thm:stationary-law}, and not in general. We expect the reversals of memoryless blocks
to be the only obstruction.

\begin{conjecture}[the stationary fiber]\label{conj:fiber}
Let $Q$ be irreducible. An irreducible generator $Q'$ gives the same stationary law of $O(T)$ as
$Q$ if and only if it is reached from $Q$ or from $Q^*$ by finitely many reversals of memoryless
blocks, each a memoryless block of the current generator.
\end{conjecture}

The \emph{if} part holds by Theorems~\ref{thm:memoryless} and~\ref{thm:stationary-law}. The
\emph{only if} part holds in the cases of Theorem~\ref{thm:stationary-law}, where the fiber is
$\{Q,Q^*\}$.

\begin{conjecture}[balanced sides]\label{conj:balanced}
A side $X$ of a cut, with the factors normalized as above, is balanced if and only if $u=D_Xa$, with
$D_X$ the diagonal matrix of $\pi_X$, and the matrix $(\pi_iq_{ij}-\pi_jq_{ji})_{i,j\in X}$ is a sum
of matrices each supported on $B\times B$ for a memoryless block $B\subseteq X$ of $Q$.
\end{conjecture}

Remark~\ref{rem:balance-necessary} proves that $u=D_Xa$ is necessary, and
Remark~\ref{rem:memoryless-cut} gives the case where $X$ is itself a memoryless block.

% Empty stub. No such section exists in this version and nothing is
% claimed about the exact law beyond Remark rem:stationary and
% Proposition prop:amplitudes.
% ---------------------------------------------------------------------
\section{Modal supports}\label{sec:modal}

This section proves Theorem~\ref{thm:B}. The observable is artificial, and we say so first.
The survival probability under calibrated killing is a finer observation, it is easier to
sample, and it gives every principal minor at once (Remark~\ref{rem:survival}). We know of no
experiment in which the modal support is observed and the survival probability is not. We
study the label because the forward theory of~\cite{paperC1} describes the law of the counts
through it, and so does any exact optimizer that reports only a maximizer. Theorem~\ref{thm:B}
says that its first-order form loses nothing that $J_V$ keeps.

\subsection{The killed chain and two labels}\label{sec:observation}

Fix a field $v\in\R^d$ and $\tau>0$, put $T=\tau L$, and weight each composition by
$e^{Tv\cdot n/N}$. If $v\le0$, the weight of a path is the probability that the chain survives
when it is killed at each visit to $i$ with probability $1-e^{Tv_i/N}$. So the maximizers of
$P_N(n)e^{Tv\cdot n/N}$ are the modes of the composition of the surviving chain. On a face the
killing shifts the diagonal, $Q_F+\operatorname{diag}v_F$ in place of $Q_F$. The shift
$v\mapsto v+c\ones$ multiplies the weighted law by $e^{Tc}$ and moves every score
$\Phi_{\tau,v}(F)$ by $-\tau c$, so it changes no winner and every field may be taken negative.
The field and the time scale are calibrated, that is, known in the units of $Q$.

\begin{proposition}[the controlled score law]\label{prop:controlled-score}
Fix $Q$, a start $\mu$ with full support, a field $v$ and $\tau>0$, and let $T=\tau L$. For every
admissible face $F$, $\log M_F(T,v)\le-\Phi_{\tau,v}(F)\,L+O(\log L)$, with equality up to
$O(\log L)$ if $F$ is irreducible. A reducible admissible face $F$ has
$\Phi_{\tau,v}(F)\ge\Phi_{\tau,v}(C)+1$ for one of its strongly connected components $C$. Hence if
$F$ is the unique minimizer of $\Phi_{\tau,v}$ over the admissible faces, then $F$ is irreducible
and the support of every maximizer of the weighted law is $F$ for all large $N$.
\end{proposition}

\begin{proof}
For large $N$ we have $T_0\le\tau L\le N^{1/3}$. Theorem~\ref{thm:first}(a) gives the upper bound,
and Theorem~\ref{thm:facemax} the equality, since an irreducible face satisfies (H) by
Lemma~\ref{lem:admissible}(a). Let $F$ be reducible and admissible. Then $\lambda_F(v)$ is the
least $\lambda_C(v)$ over the components, and a component $C$ where it is attained is
irreducible with $|C|\le|F|-1$, so $\Phi_{\tau,v}(F)=\Phi_{\tau,v}(C)+|F|-|C|$. If $F$ is the unique
minimizer and $\varepsilon>0$ is its least gap to the other admissible faces, their weighted maxima
are at most $N^{-\Phi_{\tau,v}(F)-\varepsilon+o(1)}$, that of $F$ is at least
$N^{-\Phi_{\tau,v}(F)-o(1)}$, and a face that is not admissible carries no mass.
\end{proof}

Let $\Omega=(0,\infty)\times\R^d$ and fix a start of full support. We distinguish two labels.
The \emph{finite-horizon label} $\mathcal W_N(\omega)$ at $\omega=(\tau,v)$ is the set of supports
of the maximizers of $n\mapsto P_N(n)e^{Tv\cdot n/N}$. The \emph{first-order modal-support map}
is defined on the strict set $\mathcal S_Q$ of the points $\omega$ at which $\Phi_{\tau,v}$ has a
unique minimizer over the admissible faces, and $W_Q(\omega)$ is that minimizer.

\begin{lemma}[strict regions]\label{lem:strict}
Let $Q$ be a generator and $\mu$ a start of full support.
\begin{enumerate}[(a)]
\item $\mathcal S_Q$ is open and dense in $\Omega$, and $W_Q$ is locally constant on it.
\item At every $\omega\in\mathcal S_Q$ the face $W_Q(\omega)$ is irreducible, and
$\mathcal W_N(\omega)=\{W_Q(\omega)\}$ for all $N\ge N_0(\omega)$.
\item $\mathcal S_Q$ and $W_Q$ are invariant under $(\tau,v)\mapsto(\tau,v+c\ones)$.
\end{enumerate}
\end{lemma}

\begin{proof}
(b) is Proposition~\ref{prop:controlled-score} and (c) is the shift rule. For (a), the scores are
continuous in $(\tau,v)$ and there are finitely many faces, so $\mathcal S_Q$ is open and $W_Q$
locally constant. A reducible face never attains the least score, so the complement of
$\mathcal S_Q$ lies in the union of the closed sets $Z_{FH}$ where
$\Phi_{\tau,v}(F)=\Phi_{\tau,v}(H)$, over pairs of distinct irreducible faces. If $|F|\ne|H|$, the
difference vanishes for at most one $\tau$ at fixed $v$. If $|F|=|H|$, pick $j\in F\setminus H$.
Then $\lambda_H(v)$ does not depend on $v_j$, while
$\partial\lambda_F/\partial v_j=-\alpha^*_F(v)_j<0$ by Proposition~\ref{prop:variational}(b). So
each $Z_{FH}$ has empty interior, and so has their finite union.
\end{proof}

By Lemma~\ref{lem:strict}(b) the finite-horizon labels at all $N$ determine $W_Q$ on its strict
set. They may determine more, since the positions of the label changes at finite $N$ involve the
constants $K_F(v)$ (Section~\ref{sec:verify}). Theorem~\ref{thm:B} concerns $W_Q$ only, and only
for $W_Q$ do we say \emph{exactly}.

\begin{definition}[controlled boundaries]\label{def:controlled}
Fix a start $\mu$ of full support. Two generators $Q$ and $Q'$ have the same controlled
boundaries if $W_Q=W_{Q'}$ on $\mathcal S_Q\cap\mathcal S_{Q'}$.
\end{definition}

The set $\mathcal S_Q\cap\mathcal S_{Q'}$ is open and dense, and both maps are locally constant on
it. So $Q$ and $Q'$ fail to have the same controlled boundaries exactly when there is a nonempty
open set of $\omega$ on which both have a strict winner and the winners differ. By
Lemma~\ref{lem:strict}(c) it suffices to compare the maps on $(0,\infty)\times(-\infty,0)^d$,
where every field is a killing rate. Ties are never observed and never used.

\subsection{Suppression and exposure}\label{sec:exposure}

For a Metzler matrix $B$ and a positive vector $u$, the Perron root of $B$ is at most
$\max_i(Bu)_i/u_i$, by the bound of Collatz and Wielandt~\cite{bermanplemmons1994}.

\begin{lemma}[suppression]\label{lem:suppression-nr}
Let $Q$ be a generator, $F$ a nonempty face, $\mathcal K\subset\R^F$ compact and $\varepsilon>0$.
There is $\bar v_0$, depending only on $Q$, $F$, $\mathcal K$ and $\varepsilon$, such that for every
field $v$ with $v_F\in\mathcal K$ and $v_j\le-\bar v_0$ for all $j\notin F$, and every face
$G\not\subseteq F$, we have $\lambda_G(v)\ge\lambda_{G\cap F}(v)-\varepsilon$ if
$G\cap F\ne\emptyset$, and $\lambda_G(v)\ge\bar v_0$ if $G\cap F=\emptyset$. Every larger value of
$\bar v_0$ also serves.
\end{lemma}

\begin{proof}
Put $P=G\cap F$, $S=G\setminus F$ and $B=Q_G+\operatorname{diag}v_G$. If $P=\emptyset$, every row
sum of $B$ is at most $-\bar v_0$, and the bound with $u=\ones$ gives $\lambda_G(v)\ge\bar v_0$.
Let $P\ne\emptyset$, $\varrho_P(v)=-\lambda_P(v)$, and let $\mathbf J_P$ have ones off the diagonal
of $P$. For $\zeta>0$ the Perron root of the irreducible matrix
$Q_P+\operatorname{diag}v_P+\zeta\mathbf J_P$ decreases to $\varrho_P(v)$ as $\zeta\downarrow0$,
continuously in $v_P$. By Dini's theorem one $\zeta$ makes it at most $\varrho_P(v)+\varepsilon/2$
for all $v_F\in\mathcal K$. Let $x$ be its Perron vector with $\max x=1$, let $x_*>0$ be the least
entry of $x$ over $\mathcal K$, and test $B$ with $u=x$ on $P$ and $u=y_0=\varepsilon
x_*/(2q_{\max})$ on $S$. For $p\in P$, since $\mathbf J_Px\ge0$,
$(Bu)_p\le(\varrho_P+\frac\varepsilon2)x_p+y_0q_{\max}\le(\varrho_P+\varepsilon)x_p$. For $s\in S$,
the off-diagonal entries of row $s$ inside $S$ sum to at most $q_s$, so
$(Bu)_s\le q_{\max}+y_0v_s\le q_{\max}-y_0\bar v_0\le(\varrho_P+\varepsilon)y_0$ once
$\bar v_0\ge q_{\max}/y_0+\max_{\mathcal K}|\varrho_P|+\varepsilon$. There are finitely
many $P$ and $S$, so one $\bar v_0$ serves for all, and so does every larger value.
\end{proof}

For an irreducible face $F$ with $k=|F|\ge2$ and an irreducible face $H\subsetneq F$,
singletons included, put
\[
 \rho_H(v)=\frac{k-|H|}{\lambda_H(v)-\lambda_F(v)},\qquad v\in\R^F .
\]
The denominator is positive by Lemma~\ref{lem:face-monotone}, $\rho_H$ is real-analytic, and
$\Phi_{\tau,v}(F)<\Phi_{\tau,v}(H)$ exactly when $\tau>\rho_H(v)$. So as $\tau$ decreases, the
first proper subface to tie with $F$ maximizes $\rho_H(v)$.

\begin{lemma}[exposure by descent]\label{lem:exposure}
Let $Q$ be irreducible with a start of full support, and let $F$ be an irreducible face with
$|F|\ge2$. There is a closed nowhere dense set $\mathcal E\subset\R^F$ such that every
$v^0\in\R^F\setminus\mathcal E$ has a unique partner $G$, the irreducible face $G\subsetneq F$
that maximizes $\rho_H(v^0)$. Put $\tau_1=\rho_G$. There are an open neighbourhood
$\mathcal U$ of $v^0$ with compact closure and numbers $\delta>0$ and $\bar v_0$, depending on
$Q$, $F$ and $v^0$, such that on the open set
\[
 O=\bigl\{(\tau,v)\in\Omega:\ v_F\in\mathcal U,\ v_j<-\bar v_0\ \text{for } j\notin F,\
 |\tau-\tau_1(v_F)|<\delta\bigr\}
\]
the following hold.
\begin{enumerate}[(i)]
\item The points of $O$ with $\tau>\tau_1(v_F)$ lie in $\mathcal S_Q$ with $W_Q=F$, and those
with $\tau<\tau_1(v_F)$ lie in $\mathcal S_Q$ with $W_Q=G$.
\item On the hypersurface $\tau=\tau_1(v_F)$ we have
$\lambda_F(v)=\lambda_G(v)-(|F|-|G|)/\tau_1(v_F)$.
\end{enumerate}
For $F=V$ the condition on $v_j$ is empty.
\end{lemma}

\begin{proof}
\emph{Generic partners.} Let $H_1\ne H_2$ be irreducible proper subfaces of $F$ with some
$j\in H_1\setminus H_2$, and let $v=te_j$ with $t\to\infty$. Then $\lambda_{H_2}$ stays fixed,
$-t\le\lambda_F\le\lambda_{H_1}\le q_j-t$, so $\rho_{H_2}\to0$ while
$\rho_{H_1}\ge(k-|H_1|)/q_j>0$. So $\rho_{H_1}-\rho_{H_2}$ is a nonzero real-analytic function, and
the set $\mathcal E$ where the maximum of $\rho_H$ is attained twice is closed and nowhere dense.

\emph{The boundary.} Let $G$ be the partner of $v^0\notin\mathcal E$. Choose $\delta>0$ with
$3\delta<\tau_1(v^0)$ and $\rho_H(v^0)<\tau_1(v^0)-3\delta$ for every other irreducible
$H\subsetneq F$, and take for $\mathcal U$ a ball around $v^0$ on whose closure these hold with
$2\delta$. Let $\tau_\pm$ be the largest and least $\tau$ on the closure of $O$, so $\tau_->\delta$.
There $\Phi(H)-\Phi(F)=(\lambda_H-\lambda_F)(\tau-\rho_H)\ge\kappa>0$ for every such $H$, while
$\Phi(F)-\Phi(G)=(\lambda_G-\lambda_F)(\tau_1(v_F)-\tau)$ changes sign on the hypersurface, and a
reducible face loses to a component. So among the faces inside $F$ the least score is attained
only by $F$ above and only by $G$ below. For the faces $G'\not\subseteq F$ apply
Lemma~\ref{lem:suppression-nr} with the closure of $\mathcal U$ and $\varepsilon=1/(2\tau_+)$. If
$P=G'\cap F\ne\emptyset$ and $C$ is a component of $P$ with $\lambda_C(v)=\lambda_P(v)$, then
$\Phi(G')\ge|P|+\tau\lambda_P(v)-\frac12\ge\Phi(C)+\frac12$. Enlarge $\bar v_0$ until $\tau_-\bar
v_0$ exceeds the largest score of $F$ on the closure of $O$ by $\frac12$. Then a face with $P$
empty also scores $\frac12$ above the least score, which proves (i). The tie
$\Phi(F)=\Phi(G)$ on the hypersurface is (ii).
\end{proof}

\begin{theorem}[the modal supports and the minors]\label{thm:modal}
Let $Q$ and $Q'$ be irreducible generators and let the start have full support. They have the
same controlled boundaries if and only if $-Q$ and $-Q'$ have the same principal minors. If $Q$ is
irreducible and $Q'$ is any generator with the controlled boundaries of $Q$, then $Q'$ is
irreducible.
\end{theorem}

\begin{proof}
\emph{Equal minors.} By Theorem~\ref{thm:dv}, $Q$ and $Q'$ have the same irreducible faces and
the same $\lambda_F$ on each. With a start of full support every irreducible face is admissible,
and a reducible admissible face scores above one of its components. So $\mathcal S_Q=\mathcal
S_{Q'}$ and $W_Q=W_{Q'}$.

\emph{Irreducibility.} At a fixed field, $V$ is the strict winner for $Q$ once $\tau$ is large
(Lemma~\ref{lem:face-monotone}), on an open set of $\Omega$. That set meets the dense set
$\mathcal S_{Q'}$, where $W_{Q'}=V$. A reducible face is never a strict winner, so $Q'$ is
irreducible.

\emph{Singletons.} Fix $i\ne1$ and let $\bar\tau=1/(2\max(q_{\max},q'_{\max}))$. Let $I$ be a
compact interval in $(0,\bar\tau]$ and $J$ a compact interval with $q_i-q_1$ and $q'_i-q'_1$ in its
interior. Consider the open set of $(\tau,v)$ with $\tau$ in the interior of $I$, $|v_1|<1$,
$v_i-v_1$ in the interior of $J$ and $v_j<-\bar v_0$ for $j\notin\{1,i\}$. Here $\bar v_0$ is the
larger threshold of Lemma~\ref{lem:suppression-nr} for $Q$ and $Q'$, with $F=\{1,i\}$, the compact
set $[-1,1]\times([-1,1]+J)$ of $(v_1,v_i)$ and $\varepsilon=1/(4\bar\tau)$, enlarged as in
Lemma~\ref{lem:exposure}. The pair $\{1,i\}$, if it is
irreducible, has exit rate at least $\min(q_1-v_1,q_i-v_i)-\sqrt{q_{1i}q_{i1}}$ and
$\sqrt{q_{1i}q_{i1}}\le q_{\max}$, so it scores at least $\frac12$ above the better singleton, for
both generators. Hence the strict winners are $\{1\}$ and $\{i\}$ on the two sides of
$v_i-v_1=q_i-q_1$ for $Q$, and of $v_i-v_1=q'_i-q'_1$ for $Q'$. If the hyperplanes differed, the
slab between them would carry different winners. So $q'_i-q'_1=q_i-q_1$, and with $c=q'_1-q_1$ we
get $\lambda'_{\{i\}}=\lambda_{\{i\}}+c$ for every $i$.

\emph{The descent.} We show by induction on $|F|$ that every irreducible face $F$ of $Q$ is
irreducible for $Q'$ with $\lambda'_F=\lambda_F+c$ on $\R^F$. Let $|F|\ge2$, choose
$v^0\notin\mathcal E$ with partner $G$, and take $O$ from Lemma~\ref{lem:exposure}. On the dense
set $O\cap\mathcal S_{Q'}$ we have $W_{Q'}=W_Q$, so $F$ and $G$ are irreducible for $Q'$, and by
continuity $\Phi'(F)=\Phi'(G)$ on the hypersurface. With (ii) and the induction hypothesis,
\[
 \lambda'_F(v)=\lambda'_G(v)-\frac{|F|-|G|}{\tau_1(v_F)}=\lambda_G(v)+c-\frac{|F|-|G|}{\tau_1(v_F)}
 =\lambda_F(v)+c\qquad(v_F\in\mathcal U),
\]
and Lemma~\ref{lem:perron-branch} extends this to $\R^F$. For $F=V$ both generators have
$\lambda_V(0)=0$, so $c=0$, and Lemma~\ref{lem:perron-branch} gives equal principal minors.
\end{proof}

The descent is an algorithm. For each face choose a field at random, decrease $\tau$ from a large
value, and read the first change of label and its position. Its open sets have thresholds that
depend on $Q$ and $Q'$, so no experiment is fixed in advance, and the statement is one of
identifiability.

\section{Low temperature and the cycle hierarchy}\label{sec:lowtemp}

We use the cycles of Freidlin and Wentzell~\cite{freidlinwentzell2012} in the
form of Catoni~\cite{catoni1999}. The depth of a face
(Theorem~\ref{thm:depth}) is a short consequence of Catoni's theorem on the
exit time of a domain, and the prefactors (Proposition~\ref{prop:prefactor})
are classical in substance~\cite{wentzell1972}. Section~\ref{sec:cycle-tree} proves
Theorem~\ref{thm:D}. Its main steps are the depth of a face under an added calibrated
killing (Lemma~\ref{lem:killed-depth}), the strict growth of depths along nested cycles
(Lemma~\ref{lem:nested}), and a comparison of the rewarded cycle with each kind of face.

\subsection{Arrhenius families}\label{sec:arrhenius}

Let $V_{ij}\ge0$ for $i\ne j$ be fixed costs and let
\[
 q_{ij}(\eta)=e^{-V_{ij}/\eta},\qquad \eta>0,
\]
so every off-diagonal rate is positive at every positive temperature. We
write $Q_\eta$ for the generator, and for a face $F$ we write $\lambda^\eta_F$
for the exit rate $\lambda_F$ of $F$ under $Q_\eta$, and $\lambda_F$ when $\eta$ is fixed. In this section the
letter $h$ is a height. For a proper face $F$ define
its depth
\[
 D(F)=\lim_{\eta\downarrow0}\bigl(-\eta\log \lambda^\eta_F\bigr),
\]
whose existence is part of Theorem~\ref{thm:depth}. Put $D(V)=\infty$. A fixed
height $h>0$ sets the time scale
\[
 \tau_\eta=e^{h/\eta},\qquad T=\tau_\eta L.
\]
Throughout the section we take $N\to\infty$ first at fixed $\eta$ and then
$\eta\downarrow0$, and every threshold in $N$ may depend on $\eta$. The start
$\mu_\eta$ may depend on $\eta$ but has full support and does not depend on
$N$. Since all rates are positive, every face is irreducible, and with a
full-support start every face satisfies (H) by Lemma~\ref{lem:admissible}(a).
So Theorem~\ref{thm:first}(b) applies at every fixed $\eta$ to every face.
At fixed $\eta$ the face maxima of Theorem~\ref{thm:facemax} give
\begin{equation}\label{eq:lowtemp-score}
 \log M_F=-\Phi_{\eta,h}(F)\,L+\frac{|F|-1}2\log L+O_\eta(1),\qquad
 \Phi_{\eta,h}(F)=|F|-1+e^{h/\eta}\lambda^\eta_F.
\end{equation}
So $\Phi_{\eta,h}$ is the exponent $\Phi_\tau$ of Theorem~\ref{thm:first} for the
generator $Q_\eta$ at $\tau=e^{h/\eta}$. The constants in~\eqref{eq:lowtemp-score}
may depend badly on $\eta$, and nothing below exchanges the two limits.

We use the cycles of Freidlin and Wentzell in Catoni's
convention~\cite[Definition~4.2]{catoni1999}. A nonempty
proper set $C$ is a cycle if it is a singleton or if, for all distinct
$x,y\in C$, the probability of leaving $C$ before hitting $y$, starting from
$x$, is exponentially small in $1/\eta$. Its depth $\Gamma(C)$ is the
exponential scale of its mean exit time, which does not depend on the
starting point in $C$~\cite[Proposition~4.6]{catoni1999}, and
$\Gamma(\{i\})=\min_{j\ne i}V_{ij}$.
For a reversible family write $V_{ij}=H_{ij}-U_i$ with $H$ symmetric, so that
$\pi_i(\eta)\propto e^{-U_i/\eta}$. Let $\beta(C)=\min_{i\in C,j\notin C}H_{ij}$ be the
boundary height of a proper set $C$ with at least 2 states. Then $C$ is a cycle exactly
when any 2 of its points are joined by a path whose saddles $H$ are all less than $\beta(C)$,
and then $\Gamma(C)=\beta(C)-\min_{i\in C}U_i$
\cite[Definition~3.5, Propositions~3.10 and~3.15]{cirillonardisohier2015}.
We call this the energy criterion. The inequality is strict, since a weak one
would admit sets at a common saddle level that are not cycles.

\begin{theorem}[the depth of a face]\label{thm:depth}
For every nonempty proper face $F$ the depth $D(F)$ exists and
\[
 D(F)=\max\{\Gamma(C):C\subseteq F\ \text{a cycle}\}.
\]
In particular $D(C)=\Gamma(C)$ for every cycle $C$, and $F$ contains a cycle
$C$ with $D(C)=D(F)$. Moreover, with $\mathcal G_F=(-Q_F)^{-1}$ and
$M^{\mathrm{exit}}_F=\max_{i\in F}\E_i\tau_{F^c}$,
\[
 \frac1{M^{\mathrm{exit}}_F}\le \lambda^\eta_F\le\frac{|F|}{M^{\mathrm{exit}}_F},
\]
so that $D(F)=\lim_{\eta\downarrow0}\eta\log M^{\mathrm{exit}}_F$.
\end{theorem}

\begin{proof}
Fix $\eta$ and a proper face $F$, and write $\lambda_F$ for $\lambda^\eta_F$.
By Lemma~\ref{lem:face-monotone} with $v=0$, $\lambda_F>\lambda_V=0$. So
$\mathcal G_F=\int_0^\infty e^{tQ_F}\,dt$ converges, and since
$(e^{tQ_F})_{ij}=\Prob_i(X_t=j,\ t<\tau_{F^c})$ for the chain $X$ with
generator $Q_\eta$,
\[
 (\mathcal G_F)_{ij}=\E_i\int_0^{\tau_{F^c}}\mathbf 1\{X_t=j\}\,dt
 =\Prob_i(\tau_j<\tau_{F^c})\,(\mathcal G_F)_{jj}\le(\mathcal G_F)_{jj}
\]
by the strong Markov property at $\tau_j$.

\emph{The exit-time bounds.} The row sums of $\mathcal G_F$ are the mean exit
times $\E_i\tau_{F^c}$. The spectral radius $\rho$ of a nonnegative matrix is at
most its largest row sum and at least each of its diagonal entries. Hence
\[
 \rho(\mathcal G_F)\le M^{\mathrm{exit}}_F\le\operatorname{tr}\mathcal G_F
 \le|F|\max_j(\mathcal G_F)_{jj}\le|F|\rho(\mathcal G_F).
\]
The Perron vector of $-Q_F$ is positive and is an eigenvector of
$\mathcal G_F$ for $1/\lambda_F$. A positive eigenvector of a nonnegative
matrix belongs to its spectral radius~\cite{bermanplemmons1994}, so
$\rho(\mathcal G_F)=1/\lambda_F$, which gives the bounds.

\emph{The depth.} Fix $\varkappa>d$ and put $P_\eta=I+Q_\eta/\varkappa$. Every
rate is at most $1$, so every exit rate is at most $d-1$, and $P_\eta$ is stochastic
with holding probabilities at least $1-(d-1)/\varkappa>0$. Then $-\eta\log P_\eta(i,j)\to V_{ij}$ for $i\ne j$ and
$-\eta\log P_\eta(i,i)\to0$. So $(P_\eta)$ has rare transitions in the sense
of~\cite[Definition~2.1]{catoni1999}, with $\beta=1/\eta$ and the rate function
equal to $V_{ij}$ off the diagonal and to $0$ on it, and this rate function is
irreducible since all rates are positive. The chain $P_\eta$
visits the same states in the same order as $X$, so the two chains have the
same cycles. The continuous chain is the discrete one run at the jump times of
a Poisson clock of rate $\varkappa$, so its mean exit times are the discrete ones
divided by $\varkappa$. For a domain $\mathcal O$ and $x\in\mathcal O$, Catoni's
Proposition~4.2~\cite{catoni1999} gives the limit $H_{\mathcal O}(x)$ of $\beta^{-1}\log$ of
the mean exit time from $\mathcal O$ started at $x$. Let $\mathcal M(\mathcal O)$ be the set
of maximal cycles contained in $\mathcal O$ (his Definition~4.3), and $\Gamma(C)$ the depth of
his Proposition~4.6. His Proposition~4.11 states that
$\max_{x\in\mathcal O}H_{\mathcal O}(x)=\max\{\Gamma(C):C\in\mathcal M(\mathcal O)\}$ for every
domain $\mathcal O$ other than the whole space. With $\mathcal O=F$ this reads
\[
 \lim_{\eta\downarrow0}\eta\log M^{\mathrm{exit}}_F=\max_{C\in\mathcal M(F)}\Gamma(C).
\]
These are cycles of the whole chain, not of a chain restricted to $F$. By the
exit-time bounds $-\eta\log\lambda_F$ differs from $\eta\log M^{\mathrm{exit}}_F$
by at most $\eta\log|F|$, so $D(F)$ exists and equals this limit. Every cycle
$C\subseteq F$ lies in some $C'\in\mathcal M(F)$, and $\Gamma(C)\le\Gamma(C')$
since $\tau_{C^c}\le\tau_{C'^c}$ pathwise. So the maximum over $\mathcal M(F)$
is the maximum over all cycles in $F$, attained at some $C\in\mathcal M(F)$. If
$F$ is a cycle, then $\mathcal M(F)=\{F\}$ and $D(F)=\Gamma(F)$.
\end{proof}

\begin{theorem}[noncritical heights]\label{thm:noncritical}
Let $\mathcal D$ be the finite set of depths of proper cycles and let
$h\notin\mathcal D$. Let $k(h)$ be the least size of a cycle of depth greater
than $h$, where $V$ counts as a cycle of infinite depth and size $d$. Then for
all sufficiently small $\eta$, every minimizer of $\Phi_{\eta,h}$ is a cycle of
size $k(h)$ and depth greater than $h$, and every eventual support of the
global mode, in the limit $N\to\infty$ at that $\eta$, is such a cycle.
\end{theorem}

\begin{proof}
By Theorem~\ref{thm:depth} the depth of every proper face lies in $\mathcal D$
and differs from $h$. If $D(F)<h$, then
$e^{h/\eta}\lambda^\eta_F=e^{(h-D(F)+o(1))/\eta}\to\infty$. If $D(F)>h$, then
$e^{h/\eta}\lambda^\eta_F\to0$. The face $V$ has the score $d-1$. Call a face
deep if it is $V$ or a proper face of depth greater than $h$.

A deep face of least size is a cycle. Indeed, a deep proper face $F$ contains a
cycle $C$ with $D(C)=D(F)>h$ by Theorem~\ref{thm:depth}, and $C$ is deep, so
$C=F$ if $F$ has least size. Every cycle of depth greater than $h$ is deep. So
the deep faces of least size are the cycles of size $k(h)$ and depth greater
than $h$.

There are finitely many faces, so for all sufficiently small $\eta$ every face
that is not deep has $\Phi_{\eta,h}>d$, and every deep face of size $k(h)$ has
$\Phi_{\eta,h}<k(h)$. A deep face of larger size has $\Phi_{\eta,h}\ge k(h)$,
since $\lambda_F\ge0$. Hence every minimizer is a cycle of size $k(h)$ and
depth greater than $h$, and Theorem~\ref{thm:first}(b) at $\tau=e^{h/\eta}$
puts the support of every global mode among the minimizers for large $N$.
\end{proof}

Theorem~\ref{thm:noncritical} is a theorem about least cardinality. It does
not say that every cycle of size $k(h)$ and depth greater than $h$ wins, and
the cycles that are minimal for inclusion among the deep ones can have other
sizes and are then never selected. Among faces of equal size, minimizing
$\Phi_{\eta,h}$ is minimizing $\lambda^\eta_F$ at fixed $\eta$, whatever $h$
is. So the deeper cycle wins first, and at equal depth the smaller prefactor of
Proposition~\ref{prop:prefactor}. An exponentially small difference of rates
still decides, because $N\to\infty$ first. If two scores agree exactly at some
$\eta$, the terms of~\eqref{eq:lowtemp-score} below $L$ may decide, but they need
not. In the symmetric case, for instance two disjoint cycles built from the same costs
with a start that is symmetric between them, the two faces have the same exit rate and
the same constant, and nothing at any order separates them. The theorem then says only
that the mode lies on one of the tied cycles, and Section~\ref{sec:cycle-tree} treats
such ties as unobserved. The full-support start is needed, since otherwise some cycles
are not admissible faces.

\begin{proposition}[prefactors]\label{prop:prefactor}
For every proper face $F$ there is $a_F\in(0,\infty)$ with
$\lambda^\eta_F=a_Fe^{-D(F)/\eta}(1+o(1))$, namely $a_F=1/\rho(B_F)$ with
$B_F=\lim_{\eta\downarrow0}e^{-D(F)/\eta}\mathcal G_F(\eta)$, and
\[
 D(F)=\nu\bigl(\det(-Q_F)\bigr)-\min_{i\in F}\nu\bigl(\det(-Q_{F\setminus\{i\}})\bigr),
\]
where $\nu(f)$ is the exponent $b$ of the first nonzero term $c\,e^{-b/\eta}$
of a finite sum of exponentials $f$, and the determinant of the empty matrix
is $1$.
\end{proposition}

\begin{proof}
Every minor of $-Q_F$ is a finite real combination of exponentials
$e^{-b/\eta}$. After equal exponents are collected, a combination that is not
identically zero has a first nonzero term $c\,e^{-\nu/\eta}$ and equals
$c\,e^{-\nu/\eta}(1+o(1))$. The determinant of $-Q_F$ is positive, since its
real eigenvalues are at least $\lambda_F>0$ and the others come in conjugate
pairs. The entries of $\mathcal G_F$ are positive by the identity in the proof
of Theorem~\ref{thm:depth}. By Cramer's rule each is a cofactor divided by
$\det(-Q_F)$, so
\[
 (\mathcal G_F)_{ij}=\beta_{ij}e^{e_{ij}/\eta}(1+o(1)),\qquad \beta_{ij}>0 .
\]
The diagonal cofactor of $j$ is $\det(-Q_{F\setminus\{j\}})$, so
$e_{jj}=\nu(\det(-Q_F))-\nu(\det(-Q_{F\setminus\{j\}}))$. The exit-time
bounds give $\max_j(\mathcal G_F)_{jj}\le1/\lambda_F\le|F|\max_j(\mathcal G_F)_{jj}$,
so $D(F)=\max_je_{jj}$, which is the determinant formula. Since
$(\mathcal G_F)_{ij}\le(\mathcal G_F)_{jj}$, we have $e_{ij}\le D(F)$ for all
$i,j$. Hence $e^{-D(F)/\eta}\mathcal G_F$ converges entrywise to the matrix
$B_F$ with entries $\beta_{ij}$ where $e_{ij}=D(F)$ and $0$ elsewhere. It has a
positive diagonal entry, so $\rho(B_F)>0$. The spectral radius is continuous
in the entries, and $\rho(e^{-D(F)/\eta}\mathcal G_F)=e^{-D(F)/\eta}/\lambda_F$.
Hence $e^{D(F)/\eta}\lambda_F\to1/\rho(B_F)$.
\end{proof}

The asymptotics of principal eigenvalues of matrices with entries of order
$e^{-V_{ij}/\eta}$, and their expressions through graphs and determinants, are due to
Wentzell~\cite{wentzell1972} and Freidlin and
Wentzell~\cite[Chapter~6]{freidlinwentzell2012}, and the proof above is a short
version for killed generators. The proposition uses the exact form of the rates. A family
that satisfies only the logarithmic limits $-\eta\log q_{ij}\to V_{ij}$ need not have finite
prefactors.

\begin{theorem}[critical windows]\label{thm:critical-window}
Fix a depth $\gamma\in\mathcal D$ and $x>0$, and let $T=xe^{\gamma/\eta}L$.
Define the limiting score
\[
 \mathcal J_{\gamma,x}(F)=\begin{cases}+\infty,&D(F)<\gamma,\\ |F|-1+xa_F,&D(F)=\gamma,\\ |F|-1,&D(F)>\gamma,\end{cases}
\]
with $V$ in the last case. If $\mathcal J_{\gamma,x}$ has a unique minimizer $F_x$, then
for all sufficiently small $\eta$, $F_x$ is the unique minimizer of the score
$|F|-1+xe^{\gamma/\eta}\lambda^\eta_F$ and the eventual support of the global mode
as $N\to\infty$ at that $\eta$. The same holds, with one threshold for $\eta$,
uniformly for $x$ in a compact interval on which one face has a positive gap.
The lower envelope
$x\mapsto\min_F\mathcal J_{\gamma,x}(F)$ is concave and piecewise linear, and its
minimizers need not be cycles.
\end{theorem}

\begin{proof}
Let $S_\eta(F)=|F|-1+xe^{\gamma/\eta}\lambda^\eta_F$, which is
$\Phi_\tau(F)$ for $Q_\eta$ at $\tau=xe^{\gamma/\eta}$. By
Proposition~\ref{prop:prefactor},
$e^{\gamma/\eta}\lambda^\eta_F=a_Fe^{(\gamma-D(F))/\eta}(1+o(1))$, which tends to
$\infty$, $a_F$ or $0$ as $D(F)<\gamma$, $D(F)=\gamma$ or $D(F)>\gamma$, and
$S_\eta(V)=d-1$. On a compact interval $I\subset(0,\infty)$ the convergence
$S_\eta(F)\to \mathcal J_{\gamma,x}(F)$ is uniform in $x$, divergence included, since
$x$ is bounded below on $I$.

Suppose that one face $F$ satisfies $\mathcal J_{\gamma,x}(G)\ge \mathcal J_{\gamma,x}(F)+g$ for
all $G\ne F$ and $x\in I$, with $g>0$. There are finitely many faces, so there
is $\eta_0$ with $S_\eta(G)>S_\eta(F)$ for all $G\ne F$, $x\in I$ and
$\eta<\eta_0$. A unique minimizer $F_x$ at a single $x$ has such a gap on a
small interval around $x$, because the scores are continuous in $x$.
Theorem~\ref{thm:first}(b) at $\tau=xe^{\gamma/\eta}$ then gives the support
of the global mode, with a threshold in $N$ that may depend on $\eta$ and $x$.

The function $x\mapsto \mathcal J_{\gamma,x}(F)$ is affine for every face with
$D(F)\ge\gamma$, and $\mathcal J_{\gamma,x}(V)=d-1$ is finite. So the envelope is the
minimum of finitely many affine functions, which is concave and piecewise
linear. Example~\ref{ex:four-state} shows a minimizer that is not a cycle.
\end{proof}

At a tie of the limiting lines, Theorem~\ref{thm:critical-window} identifies
the candidates only, and the exact scores at each $\eta$ decide. Along the
envelope the slope $a_F$ of the active line decreases with $x$, and its
intercept $|F|-1$ increases. The convex hull of first-order winners
in~\cite[Theorem~4.4]{paperC1} behaves in the same way.
Putting $h=\gamma$ and treating the faces of depth $\gamma$ as faces of greater
depth would discard the term $xa_F$, which is of order $1$ in the coefficient
of $L$.

\begin{example}[a noncycle winner on an open window]\label{ex:four-state}
Take $U=(0,1,1,1)$ and the symmetric saddles $H_{12}=2$,
$H_{13}=H_{23}=H_{34}=3$, $H_{14}=H_{24}=5$, so that with $t=e^{-1/\eta}$
\[
 (q_{ij})_{i\ne j}=\begin{pmatrix}*&t^2&t^3&t^5\\ t&*&t^2&t^4\\ t^2&t^2&*&t^2\\ t^4&t^4&t^2&*\end{pmatrix}.
\]
The only proper cycle with more than one state is $\{1,2\}$, of depth $3$. The
faces of depth $3$ are $\{1,2\}$, $\{1,2,3\}$ and $\{1,2,4\}$, with prefactors
$a_{12}=2$, $a_{123}=\frac23$ and $a_{124}=2$, so at $T=xt^{-3}L$ the limiting
scores are $1+2x$, $2+\frac23x$, $2+2x$ and $3$ for $V$. Hence the eventual
winner is $\{1,2\}$ for $0<x<\frac34$, the noncycle $\{1,2,3\}$ for
$\frac34<x<\frac32$, and $V$ for $x>\frac32$. At $h=3$ and $x=1$ the noncycle
beats every cycle. Its window has width $\log2$ in the rescaled height
$(h-3)/\eta$ and collapses onto the single height $h=3$, which is why it does
not contradict Theorem~\ref{thm:noncritical}.
\end{example}

\begin{proof}
Here $q_{ij}=t^{H_{ij}-U_i}$ and $\pi\propto(1,t,t,t)$. The exit rates are
\[
 q_1=t^2+t^3+t^5,\qquad q_2=t+t^2+t^4,\qquad q_3=3t^2,\qquad q_4=t^2+2t^4 .
\]

\emph{Cycles and depths.} The singletons have depths $\min_jV_{ij}=2,1,2,2$. We
test the sets with 2 or 3 states by the energy criterion. The set $\{1,2\}$
has boundary height $3$ and internal saddle $H_{12}=2<3$, so it is a cycle of
depth $3-U_1=3$. A set that contains exactly one of the states $1$ and $2$ has
the edge $H_{12}=2$ to its complement, and no two of its points communicate
below $2$. The remaining sets $\{3,4\}$, $\{1,2,3\}$ and $\{1,2,4\}$ have
boundary height $3$. In the first two the state $3$ is joined to the others
only by saddles of height $3$, and in the third the state $4$ is joined to the
others only by saddles of height $5$. So none is a cycle. By
Theorem~\ref{thm:depth}, a face that contains $\{1,2\}$ has depth $3$ and every
other proper face has depth at most $2$. So $\mathcal D=\{1,2,3\}$, and the
faces of depth $3$ are the three named.

\emph{Prefactors.} We apply Proposition~\ref{prop:prefactor}. For $\{1,2\}$,
$\det(-Q_{12})=q_1q_2-t^3=2t^4+O(t^5)$, and the first column of the adjugate is
$(q_2,t)=(t,t)+O(t^2)$, while the second column is $(t^2,q_1)=O(t^2)$. Hence
$t^3\mathcal G_{12}\to B_{12}$ with first column $(\frac12,\frac12)$ and second
column $0$, so $\rho(B_{12})=\frac12$ and $a_{12}=2$. For $\{1,2,3\}$,
\[
 -Q_{123}=\begin{pmatrix}t^2+t^3+t^5&-t^2&-t^3\\-t&t+t^2+t^4&-t^2\\-t^2&-t^2&3t^2\end{pmatrix}.
\]
Its cofactors along the first row are $3t^2q_2-t^4=3t^3+O(t^4)$, $3t^3+t^4$
and $t^3+t^2q_2=2t^3+O(t^4)$, so expansion along that row gives the determinant
$q_1(3t^2q_2-t^4)-t^2(3t^3+t^4)-t^3(t^3+t^2q_2)=2t^6+O(t^7)$. The cofactors
along the second row are $3t^4+t^5$, $3t^2q_1-t^5$ and $t^2q_1+t^4$, and
along the third row $t^4+t^3q_2$, $t^2q_1+t^4$ and $q_1q_2-t^3$, all
$O(t^4)$. So the first column of the adjugate is $(3t^3,3t^3,2t^3)+O(t^4)$, its
other entries are $O(t^4)$, and $B_{123}$ has first column
$(\frac32,\frac32,1)$ and zeros elsewhere. Its eigenvalues are $\frac32$ and
$0$, so $a_{123}=\frac23$. For $\{1,2,4\}$ put $q_4=t^2+2t^4$. The cofactors
along the first row are $q_2q_4-t^8=t^3+O(t^4)$, $tq_4+t^8=t^3+O(t^4)$ and
$t^5+t^4q_2=O(t^5)$, so the determinant is $2t^6+O(t^7)$. The other cofactors
are $t^2q_4+t^9$, $q_1q_4-t^9$, $t^4q_1+t^6$, $t^6+t^5q_2$, $t^4q_1+t^6$ and
$q_1q_2-t^3$, all $O(t^4)$. So $B_{124}$ has first column $(\frac12,\frac12,0)$
and zeros elsewhere, and $a_{124}=2$.

\emph{The winners.} At $\gamma=3$ the faces of depth less than $3$ have
$\mathcal J_{3,x}=\infty$. The lines are $1+2x$, $2+\frac23x$, $2+2x$ and $3$. The face
$\{1,2,4\}$ is always $1$ above $\{1,2\}$. Then $1+2x<2+\frac23x$ exactly when
$x<\frac34$, $2+\frac23x<3$ exactly when $x<\frac32$, and $1+2x<3$ exactly
when $x<1$. So the minimizer is unique and equals $\{1,2\}$ on $(0,\frac34)$,
$\{1,2,3\}$ on $(\frac34,\frac32)$ and $V$ on $(\frac32,\infty)$, and
Theorem~\ref{thm:critical-window} turns each into the eventual support of the
global mode, uniformly on compact subintervals. We claim nothing at the ties
$x=\frac34$ and $x=\frac32$.

\emph{The noncycle.} At $x=1$, that is at $h=3$, the score of $\{1,2,3\}$ is
$\frac83$. The cycles of depth less than $3$ have infinite limiting score, and
the only cycles of depth at least $3$ are $\{1,2\}$ and $V$, both of score $3$.
So the noncycle beats every cycle. In the rescaled height $s=(h-3)/\eta=\log x$ the window is
$\log\frac34<s<\log\frac32$, of width $\log2$. Its heights $h=3+\eta s$ tend to
$3$, so at every fixed $h\ne3$ the window is eventually left.
\end{proof}

The mechanism is that state $3$ costs one dimension but returns the chain to
$\{1,2\}$ before it is killed at state $4$, which lowers the leading exit rate
from $2t^3$ to $\frac23t^3$. On the middle interval this gain is worth more
than the dimension. The selected face is the support of the counts, not of the
limiting occupation fractions. For $\{1,2,3\}$ the right Perron vector tends to
$(1,1,\frac23)$ up to scale, and with $\pi\propto(1,t,t,t)$ the quasi-ergodic
law $\alpha^*_{123}$ tends to the point mass at state $1$.

\begin{example}[a branch switch]\label{ex:branch-switch}
Take $U=(2,0,0)$, $H_{12}=H_{13}=5$ and $H_{23}=1$. Then $\{2,3\}$ is a cycle
of depth $5$, the singleton depths are $3,1,1$, and the winners are $\{1\}$
for $0<h<3$, $\{2,3\}$ for $3<h<5$ and $V$ for $h>5$. The two winning cycles
are disjoint nodes of the same cycle tree, so the selector jumps between
branches and never selects the nodes $\{2\}$ and $\{3\}$.
\end{example}

\begin{proof}
With $t=e^{-1/\eta}$ the rates are $q_{12}=q_{13}=t^3$, $q_{21}=q_{31}=t^5$ and
$q_{23}=q_{32}=t$. So $\lambda_{\{1\}}=2t^3$, $\lambda_{\{2\}}=\lambda_{\{3\}}=t+t^5$
and $\lambda_{\{2,3\}}=t^5$. By the energy criterion $\{2,3\}$ is a cycle, and
$\{1,2\}$, $\{1,3\}$ are not, since each has an edge of height $1$ to its
complement. Hence $\mathcal D=\{1,3,5\}$. Theorem~\ref{thm:noncritical} gives
the winner $\{1\}$ on $(0,3)\setminus\{1\}$, since $2t^3<t+t^5$, the winner
$\{2,3\}$ on $(3,5)$, where $k(h)=2$, and $V$ above $5$. At $h=1$,
$\Phi_{\eta,1}(\{1\})=2t^2\to0$ while every other face has a score that tends
to at least $1$, so Theorem~\ref{thm:first}(b) applies directly. We claim
nothing at $h=3$ and $h=5$.
\end{proof}

\subsection{Reading the cycle tree}\label{sec:cycle-tree}

We observe only the winning face, as a function of the height and of a
field, and ask how much of the cycle structure it shows. A field is a vector $v=v(\eta)\in\R^d$, which may depend on $\eta$ but not on
$N$. As in Section~\ref{sec:observation} we write
$\lambda_F(v)=-\varrho(Q_{\eta,F}+\operatorname{diag}v_F)$, with $Q_{\eta,F}$ the
killed generator of $F$ under $Q_\eta$, and
\[
 \Phi_{\eta,h}(F;v)=|F|-1+e^{h/\eta}\lambda_F(v),
\]
so that $\Phi_{\eta,h}(F;0)=\Phi_{\eta,h}(F)$. In this subsection a height $h$
may be any real number. We say that a face $G$ \emph{wins at $(h,v)$} if there
is $\eta_0>0$ such that for every $\eta<\eta_0$ the face $G$ is the unique
minimizer of $\Phi_{\eta,h}(\cdot\,;v)$ over all nonempty faces. Where two
faces of equal size have identical rates, for instance two disjoint cycles
built from the same costs, no face wins and nothing is observed.

Proposition~\ref{prop:controlled-score} turns a win into an observation. At a
fixed $\eta$ it applies to $Q_\eta$, $\mu_\eta$, the vector $v(\eta)$ of either
sign and $\tau=e^{h/\eta}$. So if $G$ wins at $(h,v)$, then for every
$\eta<\eta_0$ the support of every maximizer of the weighted law is $G$ for all
large $N$. Nothing else below is probabilistic.

Two elementary facts are used throughout. The Perron root of a Metzler matrix is
nondecreasing in its entries~\cite[Ch.~2]{bermanplemmons1994}, so $v\le w$
entrywise gives $\lambda_F(w)\le\lambda_F(v)$, and
$\lambda_F(v+\kappa\ones)=\lambda_F(v)-\kappa$ for every real $\kappa$. So
adding a constant to the field changes no winner. And by
Lemma~\ref{lem:face-monotone}, $\lambda_G(v)<\lambda_{F}(v)$ whenever
$F\subsetneq G$. In particular $D(F)\le D(G)$ for $F\subseteq G$.

\begin{proposition}[the unperturbed path]\label{prop:unperturbed}
Let $v=0$, and let a proper face $W$ win at some height $h_0\notin\mathcal D$.
Then $W$ is a cycle with $D(W)>h_0$, $W$ wins at every
$h\in[h_0,D(W))\setminus\mathcal D$, and $W$ wins at no $h>D(W)$. Hence the
supremum of the heights at which $W$ wins is $D(W)$.
\end{proposition}

\begin{proof}
By Theorem~\ref{thm:noncritical}, $W$ is a cycle of size $k(h_0)$ and depth
greater than $h_0$. Let $h\in[h_0,D(W))\setminus\mathcal D$. A face of depth
greater than $h$ has depth greater than $h_0$, so $k(h)\ge k(h_0)=|W|$, and
$k(h)=|W|$ since $D(W)>h$. By Theorem~\ref{thm:noncritical} again, for small
$\eta$ every minimizer of $\Phi_{\eta,h}$ is a cycle $G$ with $|G|=|W|$ and
$D(G)>h_0$. For such $G\ne W$,
$\Phi_{\eta,h}(G)-\Phi_{\eta,h}(W)=e^{h/\eta}(\lambda_G-\lambda_W)$ has the sign
of $\Phi_{\eta,h_0}(G)-\Phi_{\eta,h_0}(W)$, which is positive for small $\eta$.
Hence $W$ is the unique minimizer. For $h>D(W)$ we have
$e^{h/\eta}\lambda_W\to\infty$, while $\Phi_{\eta,h}(V)=d-1$.
\end{proof}

So the path shows the depth of each face it selects, as the supremum of the
heights at which it wins, and nothing about the other faces. The heights at
which a face wins may skip depths of faces that never win, and a noncycle
winner of a critical window occupies a single height. In
Example~\ref{ex:four-state} the singleton $\{4\}$
wins on $(0,1)\cup(1,2)$, since $q_4<q_1<q_3$ and $q_4<q_2$ for small $\eta$,
then $\{1,2\}$ wins on $(2,3)$ and $V$ above $3$. The cut at $1$ is the depth
of the cycle $\{2\}$, which never wins. The path shows $D(\{4\})=2$ and
$D(\{1,2\})=3$ and nothing about the cycles $\{1\}$, $\{2\}$ and $\{3\}$. The
reference is absolute, since $\lambda_V(0)=0$ makes the score of $V$ exactly
$d-1$.

Killing the complement of a face at one fixed large rate does not show its
depth either. For two states, $F=\{1\}$ and the killing rate $e^{g/\eta}$ at state
$2$, the exit rate $\lambda$ of $V$ is the smaller root of
$(q_{12}-\lambda)(q_{21}+e^{g/\eta}-\lambda)=q_{12}q_{21}$, so
$\lambda_{\{1\}}-\lambda=q_{12}q_{21}/(q_{21}+e^{g/\eta}-\lambda)=e^{-(V_{12}+V_{21}+g)/\eta}(1+o(1))$.
The winner leaves $\{1\}$ at the height $V_{12}+V_{21}+g$ and not at
$D(\{1\})=V_{12}$. By the shift rule this killing is the same experiment as the
reward $e^{g/\eta}$ on $\{1\}$. So the sign of the field is not the issue. Both
readings fail for the same reason. They look at the height at which a face stops
winning, at one fixed field.

What works is to scan the size of the field on a calibrated scale and to read a
threshold in its exponent. For a proper face $G$, a fixed constant $c>0$ and a
real $\gamma$ put
\[
 v^{G,\gamma}=c\,e^{-\gamma/\eta}\ones_G .
\]
By the shift rule $v^{G,\gamma}$ and the killing field
$-c\,e^{-\gamma/\eta}\ones_{V\setminus G}$ have the same winners, even at finite
$N$, since the weights differ by the constant factor $e^{cTe^{-\gamma/\eta}}$. We
write the field as a reward only because then $v\ge0$, so $\lambda_V(v)\le0$ and
$\Phi_{\eta,h}(V;v)\le d-1$, which shortens the comparisons. The field is constant
on $G$, so $\lambda_G(v^{G,\gamma})=\lambda_G-c\,e^{-\gamma/\eta}$ exactly. Define
the threshold of $G$ as
\[
 \vartheta(G)=\sup\bigl\{\gamma\in\R:\ G\text{ wins at }(h,v^{G,\gamma})\text{ for some }h>\gamma\bigr\},
\]
with $\sup\emptyset=-\infty$. It is read from the map $(h,v)\mapsto$ winner alone,
at heights where a winner exists. The observer needs no list of critical
heights, which depend on the depths. The restriction $h>\gamma$ makes the field
visible on the time scale, since then $e^{h/\eta}c\,e^{-\gamma/\eta}\to\infty$.
Without it every face that wins on the unperturbed path would have an infinite
threshold. A strong field, as in the example above with $\gamma=-g$, keeps $G$
winning above its depth, but only the supremum over $\gamma$ is read.

The next 2 lemmas are standard facts of cycle theory. Laminarity is Catoni's
Proposition~4.4~\cite{catoni1999} and the tree of cycles of
Trouv\'e~\cite{trouve1996} and of Olivieri and Vares~\cite[Chapter~6]{olivierivares2005}.
The strict growth of depths along nested cycles can be deduced from Catoni's
Propositions~4.5 and~4.10, which express the depth of a cycle through those of its maximal
subcycles. We include the short proofs for completeness, in the convention of
cycles through hitting probabilities used here.

\begin{lemma}[laminarity]\label{lem:laminar}
Two cycles are disjoint or nested.
\end{lemma}

\begin{proof}
Let $C_1,C_2$ be cycles with a common state $z$, and suppose that neither
contains the other. Take $x\in C_1\setminus C_2$ and $y\in C_2\setminus C_1$. Since
$y\notin C_1$, $\Prob_z(\tau_x<\tau_y)\ge\Prob_z(\tau_x<\tau_{C_1^c})$, which
tends to $1$ because $C_1$ is a cycle. Since $x\notin C_2$,
$\Prob_z(\tau_{C_2^c}<\tau_y)\ge\Prob_z(\tau_x<\tau_y)$, which tends to $1$ and
contradicts that $C_2$ is a cycle.
\end{proof}

\begin{lemma}[depth with a calibrated killing]\label{lem:killed-depth}
Let $K\subseteq[d]$, $\gamma\in\R$ and $s=c\,e^{-\gamma/\eta}$. For every proper
face $G$ the limit
$D^K(G)=\lim_{\eta\downarrow0}\bigl(-\eta\log\lambda_G(-s\ones_K)\bigr)$ exists, and
\[
 D^K(G)=\max\{D^K(C):C\subseteq G\ \text{a cycle}\}.
\]
Moreover $D^K(C)=D(C)$ if $C\cap K=\emptyset$, and $D^K(C)\le\gamma$ if
$C\subseteq K$.
\end{lemma}

\begin{proof}
Fix $\Lambda$ larger than every element of $\mathcal D$. Add a cemetery state
$\partial$. Every state of $K$ jumps to $\partial$ at rate $s$, every other state
of $[d]$ jumps to $\partial$ at rate $\varepsilon=e^{-\Lambda/\eta}$, and $\partial$
jumps to every state of $[d]$ at rate $\varepsilon$. All rates of this enlarged
chain are positive. If $\gamma<0$ we multiply all its rates by $e^{\gamma/\eta}$,
which shifts every depth by $-\gamma$ and changes no hitting probability, hence
no cycle. The proof of Theorem~\ref{thm:depth} uses the rates only through the
limits of $-\eta\log$ of the rates and through the exit-time bounds, once the
uniformization constant $\varkappa$ exceeds every exit rate. After the rescaling every
rate between two states of $[d]$ is at most $1$, every rate into $\partial$ is at most
$\max(c,\varepsilon)$, and $\partial$ has the exit rate $d\varepsilon$. So every exit rate is
at most $d-1+\max(c,1)$, and $\varkappa>d+c$ serves. The bound covers the case where some
$V_{ij}=0$ and $\gamma=0$, in which the rates $1$ and $c$ occur exactly. So the proof applies
to the enlarged chain.

For $G\subseteq[d]$ the killed generator of $G$ in the enlarged chain is
$Q_G-\operatorname{diag}(s\ones_{G\cap K}+\varepsilon\ones_{G\setminus K})$. By the two
elementary facts its exit rate $\lambda'_G$ satisfies
$\lambda_G(-s\ones_K)\le\lambda'_G\le\lambda_G(-s\ones_K)+\varepsilon$. Since
$\lambda_G(-s\ones_K)\ge\lambda_G$ and $D(G)<\Lambda$, we get
$\lambda'_G=\lambda_G(-s\ones_K)(1+o(1))$. So Theorem~\ref{thm:depth} for the
enlarged chain gives the existence of $D^K(G)$ and the identity
$D^K(G)=\max D^K(C)$ over the cycles $C\subseteq G$ of the enlarged chain.

Each such $C$ is a cycle of $Q_\eta$. Indeed, the enlarged chain is the chain
$X$ of $Q_\eta$ run until an independent killing time, and then $\partial$. If $X$
leaves $C$ before it hits $y$, so does the enlarged chain. So for distinct
$x,y\in C$ the probability $\Prob_x(\tau_{C^c}<\tau_y)$ for $X$ is at most the
same probability for the enlarged chain, which is exponentially small. Hence
$D^K(G)$ is at most the maximum of $D^K(C)$ over the cycles $C\subseteq G$ of
$Q_\eta$, and the reverse inequality holds by Lemma~\ref{lem:face-monotone}. If
$C\cap K=\emptyset$ then $\lambda_C(-s\ones_K)=\lambda_C$. If $C\subseteq K$ then
$\lambda_C(-s\ones_K)=\lambda_C+s\ge s$, so $D^K(C)\le\gamma$.
\end{proof}

\begin{lemma}[strict growth along nested cycles]\label{lem:nested}
If $C'\subsetneq C$ are cycles, then $D(C')<D(C)$. Hence $D(S)<D(C)$ for every
nonempty proper subset $S$ of a cycle $C$.
\end{lemma}

\begin{proof}
Fix $z\in C\setminus C'$. Since $C$ is a cycle there is $\alpha>0$ with
$\Prob_w(\tau_{C^c}<\tau_y)\le e^{-\alpha/\eta}$ for all distinct $w,y\in C$ and
all small $\eta$. For each $\eta$ choose $y\in C'$ with
$\E_y\tau_{C'^c}=M^{\mathrm{exit}}_{C'}$. Start at $y$ and cut the path into
rounds. A round runs until the chain leaves $C'$. If it then sits at a state
$w\in C\setminus C'$, the round continues until the chain returns to $y$ or
leaves $C$. A round ends outside $C$ only if the chain leaves $C'$ straight into
$C^c$, which forces $\tau_{C^c}<\tau_z$, or if it leaves $C$ from $w$ before it
reaches $y$. By the strong Markov property each round ends outside $C$ with
probability at most $2e^{-\alpha/\eta}$, and the rounds are independent and
identically distributed. So the number $R$ of rounds up to the exit from $C$ has
$\E R\ge e^{\alpha/\eta}/2$. Each round lasts at least its first part, of mean
$M^{\mathrm{exit}}_{C'}$, and $R$ is a stopping time for the rounds. Wald's
identity gives
\[
 M^{\mathrm{exit}}_C\ge\E_y\tau_{C^c}\ge\E R\cdot M^{\mathrm{exit}}_{C'}
 \ge\tfrac12e^{\alpha/\eta}M^{\mathrm{exit}}_{C'} ,
\]
and $D(C)\ge D(C')+\alpha$ by the last statement of Theorem~\ref{thm:depth}.
For the second claim, Theorem~\ref{thm:depth} gives a cycle $C''\subseteq S$
with $D(C'')=D(S)$, and $C''\subsetneq C$.
\end{proof}

\begin{theorem}[depths from a calibrated attraction]\label{thm:attraction}
Fix $c>0$.
\begin{enumerate}
\item[(a)] Let $G$ be a proper face. If $G$ wins at $(h,v^{G,\gamma})$ for some
$h>\gamma$, then $\gamma\le D(G)$. Hence $\vartheta(G)\le D(G)$.
\item[(b)] Let $C$ be a cycle, and put $D^-(C)=\max_{i\in C}D(C\setminus\{i\})$,
with $D^-(C)=-\infty$ for a singleton. Then $D^-(C)<D(C)$, and for every
$\gamma<D(C)$ the cycle $C$ wins at $(h,v^{C,\gamma})$ for every
$h\in(\max(\gamma,D^-(C)),D(C))$. Hence $\vartheta(C)=D(C)$.
\item[(c)] For every proper face $F$,
\[
 D(F)=\max_{G\subseteq F}\vartheta(G).
\]
\end{enumerate}
\end{theorem}

\begin{proof}
Write $s=c\,e^{-\gamma/\eta}$ and $\sigma=e^{h/\eta}$, and $\Phi(G)$ for the score of
$G$ under the field.

(a) Let $\gamma>D(G)$ and $h>\gamma$. Then $s$ is of smaller order than
$\lambda_G$, so $\sigma(\lambda_G-s)\ge\sigma\lambda_G/2\to\infty$. Since
$v^{G,\gamma}\ge0$, the score of $V$ is at most $d-1$, so $G$ is not a
minimizer for small $\eta$.

(b) The inequality $D^-(C)<D(C)$ is Lemma~\ref{lem:nested}. Fix
$h\in(\max(\gamma,D^-(C)),D(C))$ and write $v=v^{C,\gamma}$. The score of $C$ is
$|C|-1+\sigma\lambda_C-\sigma s$, where $\sigma\lambda_C\to0$ because $h<D(C)$ and
$\sigma s=c\,e^{(h-\gamma)/\eta}\to\infty$ because $h>\gamma$. On another face
$G$ the field is $s\ones_{G\cap C}=s\ones_G-s\ones_{G\setminus C}$, so
\[
 \lambda_G(v)=\lambda_G(-s\ones_{G\setminus C})-s .
\]
\emph{If $G\cap C=\emptyset$}, then $\lambda_G(v)=\lambda_G>0$, so $\Phi(G)\ge0$
while $\Phi(C)\to-\infty$.
\emph{If $G\subsetneq C$}, then $\lambda_G(v)-\lambda_C(v)=\lambda_G-\lambda_C$, and
$D(G)\le D^-(C)<h<D(C)$, so $\sigma(\lambda_G-\lambda_C)\ge\sigma\lambda_G/2\to\infty$.
\emph{If $G\supsetneq C$}, then
$\lambda_C(v)-\lambda_G(v)=\lambda_C-\lambda_G(-s\ones_{G\setminus C})\le\lambda_C$, so
$\Phi(G)-\Phi(C)\ge|G|-|C|-\sigma\lambda_C\to|G|-|C|\ge1$. The face $V$ is of this
kind.
\emph{Otherwise} $G$ meets $C$ without containing it or lying in it. Put
$K=G\setminus C$. Then $\lambda_G(v)-\lambda_C(v)=\lambda_G(-s\ones_K)-\lambda_C$. A
cycle $C'\subseteq G$ that met both $C$ and $K$ would be nested with $C$ by
Lemma~\ref{lem:laminar}, which is impossible since $C'$ meets $K$ and
$C\not\subseteq G$. So $C'$ lies in $G\cap C$ or in $K$. In the first case
Lemma~\ref{lem:killed-depth} gives $D^K(C')=D(C')\le D(G\cap C)\le D^-(C)$, and in
the second $D^K(C')\le\gamma$. By Lemma~\ref{lem:killed-depth},
$D^K(G)\le\max(D^-(C),\gamma)<h$, so $\sigma\lambda_G(-s\ones_K)\to\infty$, while
$\sigma\lambda_C\to0$.

In every case $\Phi(G)-\Phi(C)$ tends to $+\infty$ or has limit inferior at
least $1$. There are finitely many faces, so $C$ wins at $(h,v)$. The interval
of $h$ is nonempty for every $\gamma<D(C)$, so $\vartheta(C)\ge D(C)$, and (a) gives
equality.

(c) For $G\subseteq F$, (a) and the monotonicity of $D$ give
$\vartheta(G)\le D(G)\le D(F)$. By Theorem~\ref{thm:depth} there is a cycle
$C\subseteq F$ with $D(C)=D(F)$, and $\vartheta(C)=D(F)$ by (b).
\end{proof}

The theorem needs no uniqueness. Under its own attraction the score of $C$
tends to $-\infty$, so $C$ wins even when other faces tie, as for two disjoint
cycles of equal rates. The unperturbed path has no such property.
Only the threshold in $\gamma$ is used, since above $D(C)$ the faces that
contain $C$ compete through $\sigma(\lambda_C(v)-\lambda_G(v))$, which can
diverge. A noncycle may win under its own attraction and its threshold may
reach $D(G)$, which is why (c) takes the maximum over all subfaces and needs no
prior knowledge of the cycles. In the deep-singleton family with
$V_{12}=V_{21}=1$, $V_{13}=V_{23}=3$ and $V_{31}=V_{32}=5$, the singleton
$\{3\}$ of depth $5$ wins on all of $(0,5)$ along the unperturbed path, so the
path never shows the cycle $\{1,2\}$ of depth $3$, while
Theorem~\ref{thm:attraction}(b) reads $\vartheta(\{1,2\})=3$.

\begin{theorem}[cycles from depths]\label{thm:cycles-from-depths}
For a nonempty proper face $C$ the following are equivalent.
\begin{enumerate}
\item[(i)] $C$ is a cycle.
\item[(ii)] $D(C')<D(C)$ for every nonempty proper subset $C'$ of $C$.
\item[(iii)] $|C|=1$, or $D(C\setminus\{i\})<D(C)$ for every $i\in C$.
\end{enumerate}
Moreover $\Gamma(C)=D(C)$ for every cycle. Hence the depth function determines
the set of cycles, their depths and the cycle tree, that is, the laminar
family of cycles ordered by inclusion with root $V$, and depths increase
strictly along every branch.
\end{theorem}

\begin{proof}
(i) implies (ii) by Lemma~\ref{lem:nested}, and (ii) and (iii) are equivalent
since $D$ is monotone. Suppose (ii). Theorem~\ref{thm:depth} gives a cycle
$C^*\subseteq C$ with $D(C^*)=D(C)$, and (ii) forces $C^*=C$. The identity
$\Gamma(C)=D(C)$ is in Theorem~\ref{thm:depth}, laminarity is
Lemma~\ref{lem:laminar}, and the strict growth is Lemma~\ref{lem:nested}.
\end{proof}

No nondegeneracy is assumed. Ties between the depths of different faces are
allowed, because the cycles are strict and Lemma~\ref{lem:nested} rules out a
tie along nested cycles. The weak form of the criterion fails. Take 3
states with all $V_{ij}=1$. Then $\lambda_{\{1\}}=2e^{-1/\eta}$ and
$\lambda_{\{1,2\}}=e^{-1/\eta}$, so $D(\{1\})=D(\{2\})=D(\{1,2\})=1$. The set
$\{1,2\}$ is not a cycle, since from state $1$ the chain leaves it before it
hits $2$ with probability $\frac12$. Criterion (ii) rejects it, and the weak
criterion $D(C')\le D(C)$ accepts it. In Example~\ref{ex:four-state} the faces
$\{1,2,3\}$ and $\{1,2,4\}$ tie in depth with the cycle $\{1,2\}$ and are
rejected for the same reason.

\begin{corollary}[the cycle tree from the functionals]\label{cor:tree-from-dv}
Let $(Q_\eta)$ and $(Q'_\eta)$ be Arrhenius families on $[d]$ with all rates
positive, and suppose that $J^\eta_V=J'^\eta_V$ on the simplex for every $\eta$
in a sequence that tends to $0$. Then $\lambda^\eta_F=\lambda'^\eta_F$ for every
face $F$ and every such $\eta$. So the two families have the same depths, the
same cycles, the same cycle tree and the same prefactors. The same holds with the
first-order modal-support maps of $Q_\eta$ and $Q'_\eta$, for a start of full
support, in place of $J^\eta_V$.
\end{corollary}

\begin{proof}
At each $\eta$ both generators are irreducible, since all rates are positive.
Theorem~\ref{thm:dv} gives $\lambda^\eta_F=\lambda'^\eta_F$ for every face, and
Theorem~\ref{thm:modal} gives the same conclusion from equal modal-support maps.
The depths and the prefactors are limits as $\eta\downarrow0$, which exist by
Theorem~\ref{thm:depth} and Proposition~\ref{prop:prefactor}, so they may be
taken along the sequence. The cycles and the tree follow by
Theorem~\ref{thm:cycles-from-depths}.
\end{proof}

The depths are even explicit in the principal minors, by the determinant formula
of Proposition~\ref{prop:prefactor}. Corollary~\ref{cor:tree-from-dv} uses the
whole functional, or the whole map, at each temperature. In the case of the maps
the descent of Section~\ref{sec:exposure} chooses its fields after $Q_\eta$ is
known. Theorem~\ref{thm:attraction} uses one family of fields, fixed in advance,
and reads the depths as thresholds. Theorems~\ref{thm:attraction}
and~\ref{thm:cycles-from-depths} and Corollary~\ref{cor:tree-from-dv} together
give Theorem~\ref{thm:D}.

The prefactors are read at a critical depth from the unperturbed path. Fix
$\gamma\in\mathcal D$, take $v=0$ and $T=xe^{\gamma/\eta}L$, and let
$\mathcal J_{\gamma,x}$ be the limiting score of Theorem~\ref{thm:critical-window}. This
is the height $h=\gamma+\eta\log x$, which moves with $\eta$, so here we
observe the winner as a function of the window coordinate $x$. Write
$e(x)=\min_F\mathcal J_{\gamma,x}(F)$ for the lower envelope, and put $\bar a_G=a_G$ if
$D(G)=\gamma$ and $\bar a_G=0$ if $D(G)>\gamma$, so that
$\mathcal J_{\gamma,x}(G)=|G|-1+x\bar a_G$ when $D(G)\ge\gamma$.

\begin{theorem}[prefactors from window endpoints]\label{thm:window-prefactors}
Suppose there are $0=x_0<x_1<\dots<x_m<x_{m+1}=\infty$ and faces $G_0,\dots,G_m$
with $G_{k-1}\ne G_k$ such that $G_k$ is the unique minimizer of $\mathcal J_{\gamma,x}$
for every $x\in(x_k,x_{k+1})$.
\begin{enumerate}
\item[(a)] $D(G_m)>\gamma$, and $D(G_k)=\gamma$ for $k<m$.
\item[(b)] $\bar a_{G_m}=0$, and for $k=m,\dots,1$,
\[
 \bar a_{G_{k-1}}=\bar a_{G_k}+\frac{|G_k|-|G_{k-1}|}{x_k}.
\]
Hence the points $x_k$ and the sizes determine $a_{G_k}$ for every $k<m$.
\item[(c)] Let $G$ be a face of depth $\gamma$ other than $G_0,\dots,G_m$, and put
$s_G=\sup_{x>0}(e(x)-|G|+1)/x$. Then $a_G\ge s_G$, and every value
$a>s_G$ in place of $a_G$ gives the same unique minimizers. If $G'\subseteq G$
has depth $\gamma$, then $a_G\le a_{G'}$.
\end{enumerate}
\end{theorem}

\begin{proof}
By Theorem~\ref{thm:critical-window}, $G_k$ is the eventual winner, uniformly on
compact subintervals of $(x_k,x_{k+1})$, so the faces $G_k$ and the points $x_k$
are observed. The envelope is continuous, so the lines of $G_{k-1}$ and $G_k$
meet at $x_k$, that is,
$|G_{k-1}|-1+x_k\bar a_{G_{k-1}}=|G_k|-1+x_k\bar a_{G_k}$. This is the recursion
in (b). Since $G_{k-1}$ lies below $G_k$ to the left of $x_k$ and above it to
the right, $\bar a_{G_{k-1}}>\bar a_{G_k}\ge0$, and a face with $\bar a_G>0$ has
depth $\gamma$. As $x\to\infty$ every line of positive slope diverges, while
$\mathcal J_{\gamma,x}(V)=d-1$. So the last piece has slope $0$, that is,
$D(G_m)>\gamma$ and $\bar a_{G_m}=0$. This proves (a) and (b). For (c), the
envelope is the minimum of the lines, so $|G|-1+xa_G\ge e(x)$ for all $x>0$,
which is $a_G\ge s_G$. If $a>s_G$, then $|G|-1+xa>e(x)$ for all $x>0$, and the unique
minimizers do not change. Finally $\lambda_G\le\lambda_{G'}$ for $G'\subseteq G$,
and at equal depths Proposition~\ref{prop:prefactor} turns this into
$a_G\le a_{G'}$.
\end{proof}

The reference is absolute because the last piece has slope exactly $0$, and
$V$, with $\lambda_V(0)=0$, guarantees a face with a flat line. In
Example~\ref{ex:four-state} at $\gamma=3$ the pieces are $G_0=\{1,2\}$,
$G_1=\{1,2,3\}$ and $G_2=V$, with $x_1=3/4$, $x_2=3/2$ and sizes
$2,3,4$. The recursion gives $a_{123}=(4-3)/(3/2)=2/3$ and
$a_{12}=2/3+(3-2)/(3/4)=2$, the values found in the example. The face
$\{1,2,4\}$ is never on the envelope. Here $s_{124}=2/3$, attained on the
whole interval $(3/4,3/2)$, and $\{1,2\}\subseteq\{1,2,4\}$ gives
$a_{124}\le a_{12}=2$. So $2/3\le a_{124}\le2$. At $a_{124}=2/3$ the
line of $\{1,2,4\}$ would coincide with that of $\{1,2,3\}$ on the middle
interval, so the hypothesis of Theorem~\ref{thm:window-prefactors} excludes
the lower end, but the observed winners alone do not. The window thus recovers
$a_{12}=2$ and $a_{123}=2/3$ and only bounds $a_{124}$, whose value is $2$.
A calibrated field need not help, because a field changes the rate of
$V$ to a value that depends on $Q$. In this example the limiting quasi-stationary
laws of the 3 faces of depth $3$ and the limit of $\pi$ all sit at state $1$. So a
field $c\,e^{-3/\eta}$ at state $1$ moves every line of the window by $-cx$ to first
order, and a field of this order at another state moves no line.

\begin{remark}[what the design needs]\label{rem:design}
The design of Theorem~\ref{thm:attraction} needs the temperature $\eta$, the time
unit, the amplitude $c$ and the labels of the states, and no rate of $Q_\eta$. It
observes the winner along a family $\eta\downarrow0$, with $N\to\infty$ first at
each $\eta$. No step of this subsection uses reversibility. Theorem~\ref{thm:depth} rests on
Catoni's theorem and the exit-time bounds, and the other lemmas use only the Green
matrix, the definition of a cycle and the strong Markov property.
Lemma~\ref{lem:killed-depth} applies Theorem~\ref{thm:depth} to a chain with a cemetery,
which is never reversible. The selection step is Proposition~\ref{prop:controlled-score},
which holds for a general generator.
\end{remark}

\section{Computations and open questions}\label{sec:verify}

Every statement is proved in the text, and no proof rests on a computation. The scripts below
check parts of the paper by assertions, with a tolerance in floating point.

The script \texttt{verify\_reconstruction.py} checks Sections~\ref{sec:minors}
to~\ref{sec:reversible}. It verifies the inversion formula of
Theorem~\ref{thm:reversible-reconstruction}(b) and the three-state identities of
Remark~\ref{rem:three-states} symbolically, and the family of Remark~\ref{rem:partial-data}. It
recovers $Q$ from the strict winners on $40$ random reversible chains with four states, with
Lemma~\ref{lem:suppression} checked on a grid of every design. It checks the constant
$20/3+28/9<10$ in the proof of Theorem~\ref{thm:stability} and the bounds $\phi_*$ and
$g_*$ on $2000$ chains. On $64$ random generators with a rank-one cut it checks that the generator
point of the cut-transpose has the principal minors of $-Q$ (in $49$ of them it is neither $Q$ nor
$Q^*$), and it checks the four-state example after Theorem~\ref{thm:nonreversible} in rational
arithmetic. It checks the derivative used in the proof of Lemma~\ref{lem:perron-branch}, and on
$200$ generators with three states it checks that every generator with the same minors is $Q$ or
$Q^*$. The script \texttt{verify\_finite\_horizon.py} computes the entry times of
Theorem~\ref{thm:finite-horizon} on a path with three states for $N=10^3$ to $10^6$ from the exact
run expansion of~\cite[Lemma~2.3]{paperC1}, and the numbers of Remark~\ref{rem:finite-size}. The
script \texttt{verify\_threshold.py} computes the least $N$ after Theorem~\ref{thm:finite-horizon}
for 2 chains and $b=2^{-1},\dots,2^{-8}$, with $Nb^2$ and $Nb^2/\log N$. The script
\texttt{verify\_cycles.py} redoes Examples~\ref{ex:four-state} and~\ref{ex:branch-switch} in exact
arithmetic. The script \texttt{verify\_attraction.py}, with the module \texttt{attraction\_core.py},
tests Theorem~\ref{thm:attraction} on $11$ cost matrices and $138$ proper faces at $t=10^{-32}$ and
$10^{-48}$, with the exact finite-$\eta$ scores. Every cycle has the grid threshold $D(G)-1/4$,
no noncycle exceeds it, and $\max_{G\subseteq F}\vartheta(G)$ gives $D(F)$ on every face.

Some finite and algebraic steps are also proved in Lean~4 with Mathlib, in
\nolinkurl{lean/PaperC} of the same repository. They cover the multiaffine expansion of $R_F$, the
normalization (N), the fact that a diagonal similarity keeping an irreducible Metzler matrix
Metzler has entries of one sign, Theorem~\ref{thm:nonreversible}(d) for positive rates, and the
determinants, adjugates and prefactors of Example~\ref{ex:four-state}. Mathlib has no Perron and
Frobenius theorem for irreducible Metzler matrices, so the positive null vector in (N) is a
hypothesis there.

\emph{Open questions.} The crossing equation of~\cite[Theorem~5.8]{paperC1} for the weighted law would read the
amplitudes of Proposition~\ref{prop:amplitudes} off the modal-support map and separate the fiber of
Theorem~\ref{thm:nonreversible} by labels alone, but it is not proved. We do not know whether a single
irreducible face is always traced against a singleton rather than a larger subface in the exposure
lemma. At low temperature it is open whether some calibrated field separates $a_{124}$ from
$a_{12}$ in Example~\ref{ex:four-state}, and a version of Theorem~\ref{thm:attraction} with $\eta$
tied to $N$ is open.

\section{A reversible design and finite horizons}\label{sec:reversible}

This section is not part of the journal version. It reads a reversible generator off the
first-order modal-support map by a design fixed in advance
(Theorem~\ref{thm:reversible-reconstruction}), counts the dimension
(Proposition~\ref{thm:minimality}), bounds the error (Theorem~\ref{thm:stability}) and treats a
finite horizon uniformly over a compact class (Theorem~\ref{thm:finite-horizon}).

\subsection{The design}\label{sec:reversible-reconstruction}

For a reversible chain the modal supports determine all of $Q$, and an explicit finite
experiment reads it. Let $Q$ be reversible and irreducible with stationary law $\pi$, on a known
connected graph $(V,E)$ with $m=|E|$ edges, where $\{i,j\}\in E$ exactly when $q_{ij}>0$. We
call it the graph $E$.
Put $w_{ij}=q_{ij}q_{ji}$ for the edge products and
\[
 S=D_{\sqrt\pi}(-Q)D_{\sqrt\pi}^{-1},\qquad S_{ii}=q_i,\qquad S_{ij}=-\sqrt{w_{ij}}\ \ (\{i,j\}\in E),
\]
the symmetrized matrix. We assume known constants $R,b>0$ with $\|S\|_2\le R$ and
$\sqrt{w_{ij}}\ge b$ on every edge. The field and the time are calibrated.

We use 2 facts throughout. First, $S$ is symmetric and similar to $-Q$, whose eigenvalues
have nonnegative real parts, so $S$ is positive semidefinite. Hence $0\le q_k\le R$ for every
$k$, $\sqrt{w_{kl}}\le R$ on every edge, and $|q_k-q_l|\le R$. Second, for a face $F$ on which
$E$ induces a connected graph, Proposition~\ref{prop:variational}(c) gives
\begin{equation}\label{eq:lambda-sym}
 \lambda_F(v)=\lambda_{\min}\bigl(S_F-\operatorname{diag}v_F\bigr).
\end{equation}
Since $Q$ has symmetric support, the admissible faces are exactly these connected faces
(Lemma~\ref{lem:admissible}).

The design uses, for a retained pair $(i,j)$, the field
\[
 v_i=\xi,\qquad v_j=0,\qquad v_k=-\bar v\ \ (k\ne i,j),
\]
which suppresses every other state, with the constants
\[
 \bar\xi=2R,\qquad \tau_0=\frac1{2R},\qquad \tau_{\max}=\frac{8R}{b^2},\qquad \bar v=4R+4\tau_{\max}R^2,
\]
and the ranges $|\xi|\le3R$ and $\frac1{4R}\le\tau\le \tau_{\max}$. Every constant is fixed by $R$ and $b$.

\begin{lemma}[suppression in the reversible design]\label{lem:suppression}
Under the design for the pair $(i,j)$, every irreducible face that contains a suppressed state
has score at least $\frac34$ larger than the best face inside $\{i,j\}$, uniformly on the ranges
of $\xi$ and $\tau$. Hence the strict winner is one of $\{i\}$, $\{j\}$ and $\{i,j\}$, and the last
only if $\{i,j\}$ is an edge.
\end{lemma}

\begin{proof}
Let $F$ be connected and let $A=F\setminus\{i,j\}$ be nonempty.

\emph{Wholly suppressed faces.} If $F\cap\{i,j\}=\emptyset$, then
$S_F-\operatorname{diag}v_F=S_F+\bar vI\ge\bar vI$, so $\lambda_F(v)\ge\bar v$
by~\eqref{eq:lambda-sym}. The singleton $\{j\}$ has score $\tau q_j\le\tau R$. Hence
$\Phi_{\tau,v}(F)-\Phi_{\tau,v}(\{j\})\ge\tau(\bar v-R)\ge(3R+4\tau_{\max}R^2)/(4R)\ge\frac34$.

\emph{Mixed faces.} Let $G_0=F\cap\{i,j\}$ be nonempty. In the order $G_0$, $A$,
\[
 S_F-\operatorname{diag}v_F=\begin{pmatrix}B&Z\\Z^\top&S_A+\bar vI\end{pmatrix},\qquad
 B=S_{G_0}-\operatorname{diag}v_{G_0},\quad Z=S[G_0,A].
\]
Here $\|Z\|_2\le R$ and $S_A\ge0$. Let $e=\lambda_{\min}(B)$, so
$|e|\le\|S_{G_0}\|_2+|\xi|\le4R$. The test vector $(x,0)$, with $x$ a unit eigenvector of $B$
for $e$, gives $\lambda_F(v)\le e$. For a unit vector $(x,y)$,
\[
 \Bigl\langle\binom xy,(S_F-\operatorname{diag}v_F)\binom xy\Bigr\rangle
 \ge e|x|^2-2R|x||y|+\bar v|y|^2\ge\lambda_{\min}\begin{pmatrix}e&-R\\-R&\bar v\end{pmatrix}
 \ge e-\frac{R^2}{\bar v-e},
\]
by $\sqrt{u^2+R^2}\le u+R^2/(2u)$ with $u=(\bar v-e)/2>0$. Since $\bar v-e\ge\bar v-4R=4\tau_{\max}R^2$,
we get $e-1/(4\tau_{\max})\le\lambda_F(v)\le e$. Now $e=\lambda_{G_*}(v)$ for an irreducible face
$G_*\subseteq\{i,j\}$. If $G_0$ is connected, take $G_*=G_0$ and use~\eqref{eq:lambda-sym}. If
$G_0=\{i,j\}$ is not an edge, $B$ is diagonal and $G_*$ is the singleton with the smaller
diagonal entry. In both cases $|G_*|\le|F|-1$, so
$\Phi_{\tau,v}(F)-\Phi_{\tau,v}(G_*)\ge1-\tau/(4\tau_{\max})\ge\frac34$ as $\tau\le \tau_{\max}$.

So the unique minimizer of $\Phi_{\tau,v}$, if there is one, is an irreducible face inside
$\{i,j\}$, and Proposition~\ref{prop:controlled-score} identifies it with the strict winner.
\end{proof}

\begin{theorem}[reversible reconstruction]\label{thm:reversible-reconstruction}
Let $Q$ be reversible and irreducible on the known connected graph $E$ with $\|S\|_2\le R$ and
$\sqrt{w_{ij}}\ge b$ on edges, and let the start have full support.
\begin{enumerate}[(a)]
\item (Exit differences.) For each $i\ne1$, under the design for the pair $(i,1)$ at
$\tau=\tau_0$, the strict winner changes from $\{1\}$ to $\{i\}$ at $\xi=x_i=q_i-q_1$, and
$|x_i|\le R$. Put $x_1=0$.
\item (Products.) For each edge $\{i,j\}$, under the design for $(i,j)$ at $\xi=\bar\xi$, the
strict winner changes from a singleton to the pair at
\[
 \tau=t_{ij}=\frac{\sqrt{s_{ij}^2+4w_{ij}}+s_{ij}}{2w_{ij}},\qquad
 s_{ij}=\bar\xi-x_i+x_j,
\]
with $1/R\le t_{ij}<\tau_{\max}/2$, and $w_{ij}=s_{ij}/t_{ij}+t_{ij}^{-2}$.
\item (Reconstruction.) Let $B$ be the symmetric matrix with diagonal $x_i$, entries
$-\sqrt{w_{ij}}$ on the edges and zero elsewhere, let $\lambda=\lambda_{\min}(B)$ and let $\phi$ be
its positive unit eigenvector. Then
\[
 S=B-\lambda I,\qquad Q=-D_\phi^{-1}SD_\phi,\qquad \pi_i=\phi_i^2 .
\]
So the $m+d-1$ limiting boundary coordinates $(x_i)_{i\ne1}$ and $(t_{ij})_{\{i,j\}\in E}$
determine $Q$, and the design is fixed by $R$, $b$ and the labels before any boundary is
observed.
\end{enumerate}
\end{theorem}

\begin{proof}
By Lemma~\ref{lem:suppression} and Proposition~\ref{prop:controlled-score}, at each point of the
ranges the strict winner is the unique minimizer of $\Phi_{\tau,v}$ among the irreducible faces
inside the retained pair, when it is unique. Here $\tau_0\in[\frac1{4R},\tau_{\max}]$, since $b\le R$.

(a) The candidates are $\{i\}$, $\{1\}$ and, if $\{i,1\}$ is an edge, the pair. By Weyl's
inequality $\lambda_{\{i,1\}}(v)\ge\min(q_i-\xi,q_1)-\sqrt{w_{i1}}$, so at $\tau_0$ the score of
the pair exceeds the better singleton score by at least $1-\tau_0R=\frac12$, and the pair never
wins. The singleton $\{i\}$ has the smaller score exactly when $q_i-\xi<q_1$, that is when
$\xi>x_i$. Since $|x_i|\le R$, both sides of $x_i$ lie in $[-3R,3R]$. Hence $\{1\}$ wins strictly
on $[-3R,x_i)$, $\{i\}$ wins strictly on $(x_i,3R]$, and the boundary is $\xi=x_i$.

(b) Put $p=q_i-\bar\xi$ and $s=s_{ij}$. Since $x_i-x_j=q_i-q_j$ by (a), $s=q_j-p$. Then
$R\le s\le3R$, and $\{i\}$ beats $\{j\}$ by $\tau s>0$. By~\eqref{eq:lambda-sym} the
pair has $\lambda_{\{i,j\}}(v)=p-g$, the smaller eigenvalue of
$\bigl(\begin{smallmatrix}p&-\sqrt w\\-\sqrt w&q_j\end{smallmatrix}\bigr)$ with $w=w_{ij}$,
where
\[
 g=\frac{\sqrt{s^2+4w}-s}2=\frac{2w}{\sqrt{s^2+4w}+s}\in(0,\sqrt w\,] .
\]
So $\Phi(\{i,j\})-\Phi(\{i\})=1-\tau g$. Hence $\{i\}$ wins strictly below $t_{ij}=1/g$, the pair
wins strictly above it, and $1/g$ is the displayed formula. Also $g\le\sqrt w\le R$ gives
$t_{ij}\ge1/R>\frac1{4R}$. The formula increases in $s$ and decreases in $w$, and $s\le3R$,
$w\ge b^2$, $b\le R$, so $t_{ij}\le(\sqrt{13}+3)R/(2b^2)<\tau_{\max}/2$. At $\tau=t_{ij}$ the pair
eigenvalue $p-1/t_{ij}$ is a root of $(p-y)(q_j-y)=w$, which gives $t_{ij}^{-1}(s+t_{ij}^{-1})=w$,
the inversion formula.

(c) $B=S-q_1I$, since $S$ has diagonal $q_k$, entries $-\sqrt{w_{kl}}$ on edges and zero off the
graph. The positive unit vector $\sqrt\pi$ satisfies $S\sqrt\pi=D_{\sqrt\pi}(-Q)\ones=0$. As
$S\ge0$, the eigenvalue $0$ is $\lambda_{\min}(S)$, and it is simple with eigenvector $\sqrt\pi$
by the theorem of Perron and Frobenius for the irreducible nonnegative matrix $RI-S$. So
$\lambda=-q_1$, $S=B-\lambda I$, $\phi=\sqrt\pi$ and $\pi_i=\phi_i^2$, and
$Q=-D_\phi^{-1}SD_\phi$ is the definition of $S$. The coordinates used are $d-1$ values $x_i$
and $m$ values $t_{ij}$, and the fields and time scales of the design depend only on $R$, $b$
and the labels.
\end{proof}

The reconstruction of a reversible generator from its exit rates and edge products is
classical. Kim and Landim~\cite[Lemma~5.1 and Corollary~5.3]{kl2024} recover the same
quantities from the values of $J_V$, and with the stationary law from $J_V^{-1}(0)$ (their
Lemma~5.2) deduce that $J_V$ determines a reversible generator. Here the coordinates are read
off the boundaries of strict regions of the weighted composition law, with a design fixed in
advance.

\begin{proposition}[the dimension count]\label{thm:minimality}
The reversible generators on the fixed connected graph $E$ form a real-analytic manifold of
dimension $m+d-1$, with global coordinates the edge products $(w_{ij})$ and the exit differences
$(q_i-q_1)$. Let $W$ be a nonempty open subset of it, for example the set where $\|S\|_2<R$ and
$\sqrt{w_{ij}}>b$ when that set is nonempty. Let $k<m+d-1$, let $f\colon W\to\R^k$ be $C^1$ and
let $\ell\colon W\to\Lambda$ be locally constant with $\Lambda$ finite. Then the observation
$(f,\ell)$ is not injective on any nonempty open subset of $W$. So no family of fewer than
$m+d-1$ locally $C^1$ scalar observations with locally constant finite labels is locally
identifying.
\end{proposition}

\begin{proof}
Let $\mathcal Q_E$ be the set of reversible generators with $q_{ij}>0$ exactly on the edges in
$E$, and $\kappa(Q)=((w_{ij})_E,(q_i-q_1)_{i\ne1})$. Given $w\in(0,\infty)^m$ and
$x\in\R^{d-1}$ (with $x_1=0$), form $B$ as in Theorem~\ref{thm:reversible-reconstruction}(c),
let $\lambda$ and $\phi$ be its least eigenvalue and positive unit eigenvector, which is simple
since $-B$ plus a large multiple of $I$ is irreducible and nonnegative, and put
$\Psi(w,x)=-D_\phi^{-1}(B-\lambda I)D_\phi$. This is a reversible generator on $E$, with respect
to $\phi^2$, with products $w_{ij}$ and exit rates $x_i-\lambda$, so $\kappa\circ\Psi=\mathrm{id}$,
and Theorem~\ref{thm:reversible-reconstruction}(c), whose proof uses only reversibility and
irreducibility, gives $\Psi\circ\kappa=\mathrm{id}$. The map $\Psi$ is real-analytic by the
implicit function theorem at a simple eigenvalue, and $\kappa$ is polynomial. So $\kappa$ is a
global real-analytic chart of $\mathcal Q_E$, of dimension $m+d-1$.

For the lower bound let $W'\subseteq W$ be nonempty and open. Shrinking it, we may assume that
$\ell$ is constant on $W'$ and that $W'$ is a ball in the chart $\kappa$, so $f$ is a $C^1$ map on
a ball of $\R^p$, $p=m+d-1$. Let $\varrho_f$ be the largest rank of $df$ on $W'$, attained at
$y_0$. The rank is lower semicontinuous, so it equals $\varrho_f$ near $y_0$, and by the constant
rank theorem $f$ is constant on a $C^1$ submanifold through $y_0$ of dimension
$p-\varrho_f\ge p-k\ge1$. So $(f,\ell)$ is not injective on $W'$.
\end{proof}

The chart is the continuous-time version of that of Wolfer and
Watanabe~\cite[Theorems~3, 5 and~6]{wolferwatanabe2021} for reversible kernels, and the lower
bound is the classical rank argument. Regularity is needed, since interleaving the digits of
the $m+d-1$ coordinates gives one injective scalar.

\begin{remark}[partial data]\label{rem:partial-data}
Suppose that the products are known on $J\subseteq E$ and the exit differences along the edges
of a comparison graph $(V,E')$. In the chart $\kappa$ the generators with these data have the
given $w_e$ on $J$ and any $w_e>0$ on $E\setminus J$. Their $x$ is the given one plus a vector
that is constant on each component of $(V,E')$ and zero on that of the state $1$. So they
are parametrized by $(0,\infty)^{m-|J|}\times\R^{\operatorname{comp}(E')-1}$, with
$\operatorname{comp}(E')$ the number of components of $(V,E')$, and the data identify $Q$ exactly
when $J=E$ and $(V,E')$ is connected. For instance, for
$\ell>\sqrt2$ and $z=\ell-2/\ell$ the generators
\[
 Q_\ell=\begin{pmatrix}-\ell&2/\ell&z\\ \ell/2&-\ell&\ell/2\\ z&2/\ell&-\ell\end{pmatrix}
\]
satisfy detailed balance with $\nu=(1,4/\ell^2,1)$, have all exit rates equal to $\ell$ and
$w_{12}=w_{23}=1$, and have $w_{13}=z^2$. So the data of a spanning tree miss a parameter.
\end{remark}

\begin{theorem}[stability]\label{thm:stability}
In the setting of Theorem~\ref{thm:reversible-reconstruction}, let the observed dimensionless
coordinates $x_i/R$ and $Rt_{ij}$ have sup-norm errors at most $\varepsilon$, and let
$\widehat Q$ be the reconstruction from the observed values. There are $\varepsilon_0$ and $C$
depending only on $d$ and $b/R$ such that for $\varepsilon\le\varepsilon_0$ the reconstruction
$\widehat Q$ is a reversible generator on $E$ and
\[
 \|\widehat Q-Q\|_2\le C\,R\,\varepsilon .
\]
Among the constants are the bound $\phi_{\min}\ge \phi_*=(b/R)^{d-1}/\sqrt d$ for the ground
vector $\phi=\sqrt\pi$ and the bound $\lambda_2(S)\ge g_*=2b\phi_*^2/(d-1)$ for the gap of $S$.
\end{theorem}

\begin{proof}
The reconstruction commutes with the scaling $Q\mapsto Q/R$, so we take $R=1$, $\bar\xi=2$ and
$b\le1$. Then $\sigma_{ij}=t_{ij}\ge1$ and $s_{ij}=2-x_i+x_j\in[1,3]$, by
Theorem~\ref{thm:reversible-reconstruction}. Hats denote reconstructed values, and
$\widehat x_1=x_1=0$.

\emph{Products.} Let $\varepsilon\le b^2/20$. Then $\widehat\sigma\ge\frac34$ and
$|\widehat\sigma^{-1}-\sigma^{-1}|\le\frac43\varepsilon$, and $w=s/\sigma+\sigma^{-2}$ gives
$|\widehat w-w|\le(\frac{20}3+\frac{28}9)\varepsilon<10\varepsilon$. So $\widehat w\ge b^2/2$ and
$|\sqrt{\widehat w}-\sqrt w|\le10\varepsilon/b$, and the largest absolute row sum gives
$\|\widehat B-B\|_2\le\beta:=10d\varepsilon/b$. By Weyl's inequality
$|\widehat\lambda-\lambda|\le\beta$, so $\|\widehat S-S\|_2\le2\beta$. The matrix $\widehat B$
has negative entries exactly on the edges in $E$, so the chart $\Psi$ of the proof of
Proposition~\ref{thm:minimality} makes $\widehat Q$ a reversible generator with the graph $E$.

\emph{The constants.} Row $l$ of $S\phi=0$ reads $q_l\phi_l=\sum_{k\sim l}\sqrt{w_{lk}}\phi_k$, so
$\phi_l\ge b\phi_k$ on every edge. Some $\phi_{k_0}$ is at least $1/\sqrt d$, and every state is
joined to $k_0$ by at most $d-1$ edges, so $\phi_{\min}\ge\phi_*$. For $y\perp\phi$ write
$y=\phi\circ f$. Then $\langle y,Sy\rangle=\sum_{\{k,l\}\in E}\chi_{kl}(f_k-f_l)^2$ with
$\chi_{kl}=\sqrt{w_{kl}}\phi_k\phi_l\ge b\phi_*^2$, while $\|y\|^2=\sum_k\pi_kf_k^2$ and
$\sum_k\pi_kf_k=0$. In a spanning tree $\mathcal T$ of $E$ the path from $k$ to $l$ has at most
$d-1$ edges, so the inequality of Cauchy and Schwarz gives
\[
 \|y\|^2=\tfrac12\sum_{k,l}\pi_k\pi_l(f_k-f_l)^2\le\tfrac{d-1}2\sum_{e\in\mathcal T}(\nabla_ef)^2
 \le\frac{d-1}{2b\phi_*^2}\langle y,Sy\rangle,
\]
and $\lambda_2(S)\ge g_*$.

\emph{The generator.} Write $\widehat\phi=a\phi+z$ with $z\perp\phi$ and $a>0$. Since
$Sz=(S-\widehat S)\widehat\phi$, we get $\|z\|\le2\beta/g_*$, and since $1-a\le1-a^2=\|z\|^2$,
also $\gamma=\|\widehat\phi-\phi\|\le\sqrt2\|z\|\le2\sqrt2\beta/g_*$. This is the bound of Davis
and Kahan in its simplest form. Let $\varepsilon$ be so small that $\gamma\le\phi_*/2$. Then
$\widehat\phi_{\min}\ge\phi_*/2$,
$|\widehat\phi_l/\widehat\phi_k-\phi_l/\phi_k|\le4\gamma/\phi_*^2$ and
$\widehat\phi_l/\widehat\phi_k\le2/\phi_*$, so on an edge
\[
 |\widehat q_{kl}-q_{kl}|
 =\bigl|\sqrt{\widehat w_{kl}}\,\widehat\phi_l/\widehat\phi_k-\sqrt{w_{kl}}\,\phi_l/\phi_k\bigr|
 \le\frac{20\varepsilon}{b\phi_*}+\frac{4\gamma}{\phi_*^2}.
\]
The diagonal satisfies $|\widehat q_k-q_k|=|\widehat S_{kk}-S_{kk}|\le2\beta$, and off the graph
both matrices vanish. So every entry of $\widehat Q-Q$ is at most $C_1(d,b)\varepsilon$, and
$\|\widehat Q-Q\|_2\le dC_1\varepsilon$. The conditions on $\varepsilon$ depend only on $d$ and
$b$.
\end{proof}

The perturbation bounds are standard. The point is that the constants come from $R$ and $b$
alone, with no separate assumption on the spectral gap.

\subsection{Finite horizons}\label{sec:finite-horizon}

Theorem~\ref{thm:reversible-reconstruction} reads the coordinates off the limiting boundaries.
At a finite horizon the observed boundaries are shifted, and the shift has an explicit form,
uniformly over a class of chains and starts. Throughout, the connected graph $E$ on $[d]$ and
the numbers $0<b\le R$ and $0<\mu_*\le1/d$ are fixed. Let $\mathcal C$ be the set of reversible
irreducible generators on $E$ with $\|S\|_2\le R$ and $\sqrt{w_{ij}}\ge b$ on every edge, and let
$\mathcal M$ be the set of starts with $\min_i\mu_i\ge\mu_*$. For a pair $(i,j)$ and $\xi\in\R$
let $v^{ij}(\xi)$ be the design field of Section~\ref{sec:reversible-reconstruction}, and let
$\mathcal V$ be the compact set of these fields with $|\xi|\le3R$.

The fixed-chain face maxima are Theorem~\ref{thm:facemax}. The finite horizon also needs
constants that do not depend on the chain and the start in the class. This is the last result
we import from~\cite{paperC1}. There it holds for an arbitrary compact set of fields, and we apply
it with $\mathcal V$. Our edge bound $b$ is the bound $w_*$ of~\cite{paperC1}, and our $S$ is its
$S_V$.

\begin{proposition}[uniform face maxima on the class, {\cite[Proposition~5.12]{paperC1}}]\label{prop:uniform-facemax}
Let $E$ induce a connected graph on the face $F$, with $k=|F|\ge2$. There are $N_0$ and $C$,
depending only on $E$, $R$, $b$ and $\mu_*$, such that for all $Q\in\mathcal C$, $\mu\in\mathcal M$
and $v\in\mathcal V$ the following hold.
\begin{enumerate}[(a)]
\item For all $N\ge2$ and $T>0$ with $TR\le N/2$,
$M_F(T,v)\le C(T/N)^{k-1}e^{-T\lambda_F(v)+CT^2/N}$.
\item For $N\ge N_0$ and $1\le T\le N^{1/3}$, we have $C^{-1}\le K_F(v)\le C$ and
\[
 \biggl|\frac{M_F(T,v)}{K_F(v)\,N^{-(k-1)}T^{(k-1)/2}e^{-T\lambda_F(v)}}-1\biggr|
 \le C\Bigl(\frac1T+\frac{T^2}N\Bigr).
\]
\end{enumerate}
For a singleton, $M_{\{i\}}(T,v)$ is given exactly by~\eqref{eq:vertexmax}.
\end{proposition}

The proof in~\cite{paperC1} is a compactness argument. On an edge $q_{ij}\ge b^2/R$, so the
class, the starts and the fields form a compact set on which every connected face stays
irreducible, and every constant in the fixed-chain proofs of Lemma~\ref{lem:upper} and
Theorem~\ref{thm:facemax} is continuous there.

\begin{theorem}[finite-horizon recovery on the class]\label{thm:finite-horizon}
Run the design of Theorem~\ref{thm:reversible-reconstruction} at horizon $N$, with $L=\log N$ and
$T=\tau L$, and let $\mathcal W_N(\tau,v)$ be the set of supports of the maximizers of
$n\mapsto P_N(n)e^{Tv\cdot n/N}$. For an ordered pair $(i,j)$ with $\{i,j\}\in E$ put
$v=v^{ij}(\bar\xi)$,
\[
 \tau^-_N=\sup\Bigl\{\tau\in[\tfrac1{4R},\tau_{\max}]:\mathcal W_N(\tau',v)=\{\{i\}\}\text{ for all }\tau'\in[\tfrac1{4R},\tau]\Bigr\},
\]
\[
 \tau^+_N=\inf\Bigl\{\tau\in[\tfrac1{4R},\tau_{\max}]:\mathcal W_N(\tau',v)=\{\{i,j\}\}\text{ for all }\tau'\in[\tau,\tau_{\max}]\Bigr\},
\]
let $t_{ij,N}$ be any number in $[\tau^-_N,\tau^+_N]$, and let $K_{ij}=K_{\{i,j\}}(v)$. Put
$\varepsilon_N=\log L/(2L)$. There are $N_0$ and $C$, depending only on $E$, $R$, $b$ and
$\mu_*$, such that for all $Q\in\mathcal C$, $\mu\in\mathcal M$ and $N\ge N_0$ the following
hold.
\begin{enumerate}[(a)]
\item (Singleton switches.) Let $i\ne1$. For every $\xi\in[-3R,3R]$ other than
\[
 x_{i,N}=\frac1{\tau_0L}\Bigl[\log\frac{\mu_1}{\mu_i}+(N-1)\log\frac{N-\tau_0Lq_1}{N-\tau_0Lq_i}\Bigr],
\]
$\mathcal W_N(\tau_0,v^{i1}(\xi))$ is $\{\{1\}\}$ for $\xi<x_{i,N}$ and $\{\{i\}\}$ for
$\xi>x_{i,N}$. Moreover $|x_{i,N}-x_i-2R\log(\mu_1/\mu_i)/L|\le CL/N$, so
$|x_{i,N}-x_i|\le2R\log(1/\mu_*)/L+CL/N$.
\item (Pair entries.) We have $\tau^+_N-\tau^-_N\le C/L^2$ and
\[
 \Bigl|t_{ij,N}-t_{ij}(1-\varepsilon_N)-\frac{c_{ij}}L\Bigr|\le C\,\frac{\log L}{L^2},\qquad
 c_{ij}=-t_{ij}\Bigl(\frac12\log t_{ij}+\log\frac{K_{ij}}{\mu_i}\Bigr),\qquad|c_{ij}|\le C .
\]
\item (Reconstruction.) Let $\widehat Q_N$ be the reconstruction of
Theorem~\ref{thm:reversible-reconstruction}(c) from the values $x_{i,N}$ and $t_{ij,N}$, and
$\widetilde Q_N$ the one from $x_{i,N}$ and $t_{ij,N}/(1-\varepsilon_N)$. Both are reversible
generators on $E$, and
\[
 \|\widehat Q_N-Q\|_2\le C\,\frac{\log L}L,\qquad\|\widetilde Q_N-Q\|_2\le\frac CL .
\]
\end{enumerate}
\end{theorem}

The threshold $N_0$ is large. The design runs up to $\tau=\tau_{\max}=8R/b^2$, and $I+(T/N)Q$
is stochastic on the whole range exactly when $N\ge \tau_{\max}q_{\max}\log N$. So the least such $N$ has
$Nb^2/\log N=8Rq_{\max}$ up to rounding, and it grows like $b^{-2}\log(1/b)$ when $q_{\max}$ stays
of order $R$. Section~\ref{sec:verify} computes it for 2 families. The proof compares the exact
singleton atoms with the uniform face maximum of the pair and solves the resulting crossing
equation for $T$.

\begin{proof}
All constants below depend only on $E$, $R$, $b$ and $\mu_*$. On the ranges of the design
$T\in[L/(4R),\tau_{\max}L]$, so for $N\ge N_1$ we have $TR\le N/2$, $T\le N^{1/3}$ and
$1/T+T^2/N\le5R/L$, and Proposition~\ref{prop:uniform-facemax} applies.

\emph{Suppression.} Fix a pair and a point of the ranges, and let $m^*$ be the least score among
the irreducible faces inside the pair. The proof of Lemma~\ref{lem:suppression} uses only $R$,
$b$ and $\bar v$, so each connected face $G'$ outside the pair has
$\Phi_{\tau,v}(G')\ge m^*+\frac34$, and Proposition~\ref{prop:uniform-facemax}(a) or the singleton
atom gives $M_{G'}\le CT^{|G'|-1}N^{-\Phi_{\tau,v}(G')}$. The best face inside the pair has face
maximum at least $C^{-1}N^{-m^*}$, by the singleton atom with $\mu_i\ge\mu_*$ or by
Proposition~\ref{prop:uniform-facemax}(b). Faces that are not connected carry no mass. So for
$N\ge N_0$ every maximizer has its support inside the pair.

\emph{(a)} At $\tau_0$ the pair $\{i,1\}$, if it is an edge, has score at least the better
singleton score plus $\frac12$, by the proof of Theorem~\ref{thm:reversible-reconstruction}(a), so
as above its face maximum is smaller than $\max(M_{\{1\}},M_{\{i\}})$ by a factor $CLN^{-1/2}$.
The two singleton atoms are exact, and
\[
 \log M_{\{i\}}-\log M_{\{1\}}=\log\frac{\mu_i}{\mu_1}+(N-1)\log\frac{N-\tau_0Lq_i}{N-\tau_0Lq_1}+\tau_0L\,\xi ,
\]
which is affine and increasing in $\xi$ with zero $x_{i,N}$. This gives the claim on
$\mathcal W_N$. Since $|\log(1-y)+y|\le y^2$ for $0\le y\le\frac12$, the middle term is
$-\tau_0L(q_i-q_1)+O(L/N)+O(R^2L^2/N)$, and dividing by $\tau_0L=L/(2R)$ gives the expansion. For
large $N$ it gives $x_{i,N}\in(-3R,3R)$, since $|x_i|\le R$.

\emph{(b)} Let $\{i,j\}$ be an edge and $v=v^{ij}(\bar\xi)$. The exact atoms give
$\log M_{\{j\}}-\log M_{\{i\}}\le\log(1/\mu_*)-TR+1$ with $TR\ge L/4$, so $\{j\}$ loses to
$\{i\}$. By the proof of Theorem~\ref{thm:reversible-reconstruction}(b),
$\lambda_{\{i\}}(v)-\lambda_{\{i,j\}}(v)=1/t_{ij}$, and the exact atom gives
$\log M_{\{i\}}=\log\mu_i-T(q_i-\bar\xi)+O(1/L)$. So Proposition~\ref{prop:uniform-facemax}(b)
with $k=2$ gives, on the whole range,
\[
 \Lambda_N(T)=\log M_{\{i,j\}}-\log M_{\{i\}},\qquad\bigl|\Lambda_N(T)+L-\Xi(T)-\kappa\bigr|\le\frac{C_5}L,
\]
with $\Xi(T)=T/t_{ij}+\frac12\log T$ and $\kappa=\log(K_{ij}/\mu_i)$. Here $|\kappa|\le C$, since
$K_{ij}$ is bounded above and below and $\mu_i\ge\mu_*$. The function $\Xi$ increases with
$\Xi'\ge1/t_{ij}\ge2/\tau_{\max}$. Let $T^\circ$ solve $\Xi(T^\circ)=L-\kappa$. Then $\Lambda_N(T)<0$ for
$T\le T^\circ-C_5\tau_{\max}/L$ and $\Lambda_N(T)>0$ for $T\ge T^\circ+C_5\tau_{\max}/L$. As $t_{ij}\in[1/R,\tau_{\max}/2)$, the
window $|T-T^\circ|<C_5\tau_{\max}/L$ lies inside the range for large $N$. With the suppression step, every
maximizer has support $\{i\}$ below the window and $\{i,j\}$ above it. Hence
$T^\circ/L-C_5\tau_{\max}/L^2\le\tau^-_N\le\tau^+_N\le T^\circ/L+C_5\tau_{\max}/L^2$, which is the first claim, and
$t_{ij,N}=T^\circ/L+O(1/L^2)$. To solve for $T^\circ$ put $\tau^\circ=T^\circ/L$. Then
$\tau^\circ=t_{ij}\bigl(1-(\frac12\log T^\circ+\kappa)/L\bigr)$, so
$\tau^\circ/t_{ij}=1+O(\log L/L)$ and $\log T^\circ=\log L+\log t_{ij}+O(\log L/L)$. Putting this
back gives $\tau^\circ=t_{ij}(1-\varepsilon_N)+c_{ij}/L+O(\log L/L^2)$, uniformly, since
$t_{ij}\le \tau_{\max}/2$ and $|\kappa|\le C$.

\emph{(c)} By (a) and (b) the dimensionless coordinates $x_{i,N}/R$ and $Rt_{ij,N}$ differ from
$x_i/R$ and $Rt_{ij}$ by at most $C\log L/L$, and after the correction by at most $C/L$, since
$t_{ij,N}/(1-\varepsilon_N)-t_{ij}=(c_{ij}/L+O(\log L/L^2))/(1-\varepsilon_N)$.
Theorem~\ref{thm:stability} gives its admissible error size $\varepsilon_0$ and its constant as
functions of $d$ and $b/R$ alone. So for $N\ge N_0$ these errors are below $\varepsilon_0$, and
that theorem gives both bounds and the generator property.
\end{proof}

\begin{remark}[a start that defeats every horizon]\label{rem:finite-start}
The bound $\min_i\mu_i\ge\mu_*$ cannot be dropped, since the start enters through the singleton
switches as $2R\log(\mu_1/\mu_i)/L$. Let $d=2$, $q_{12}=q_{21}=1$, $\tau=\frac14$ and $v=(0,1)$.
The scores are $\Phi(\{2\})=0$, $\Phi(\{1\})=\frac14$ and
$\Phi(\{1,2\})=1+\frac14\cdot\frac{1-\sqrt5}2>\frac34$, so for a fixed start $\{2\}$ wins for large
$N$. Both singleton atoms carry the factor $(1-T/N)^{N-1}$, so
$M_{\{1\}}/M_{\{2\}}=(\mu_1/\mu_2)e^{-T}=(\mu_1/\mu_2)N^{-1/4}$ exactly. For the start
$\mu_2=1/(1+2N^{1/4})$ this ratio is $2$ at every $N$. The law is linear in the start, so
Lemma~\ref{lem:upper} for the two starts at a single state gives
$M_{\{1,2\}}\le CLN^{-1/2}M_{\{1\}}$ with $C$ free of $\mu$. So $\{1\}$
wins at every large $N$, against the limiting prediction. The switch between $\{1\}$ and $\{2\}$ in
the variable $v_2$ sits at $\log(\mu_1/\mu_2)/(\tau L)=1+4\log2/L$ instead of at $q_2-q_1=0$. On
$\mathcal M$ the same shift is at most $\log(1/\mu_*)/(\tau L)$, uniformly. For the pair
entries the start enters through the term $\log(K_{ij}/\mu_i)$ of $c_{ij}$ in
Theorem~\ref{thm:finite-horizon}(b), which is unbounded as $\mu_i\to0$.
\end{remark}

\begin{remark}[the size of the correction]\label{rem:finite-size}
On the two-state class with $R=3$ and $b=1$, the chain $q_1=q_2=1$ with the uniform start has
$t_{12}=6.16$ and $c_{12}=-16.5$, and at $N=10^7$ the exact law gives $t_{12,N}=4.65$, against
$5.63$ for $t_{12}(1-\varepsilon_N)$ and $4.60$ for the prediction of (b). So the term
$c_{ij}/L$, which no correction fixed in advance removes, dominates $t_{ij}\varepsilon_N$ at
accessible $N$.
\end{remark}


\section*{Data and code availability}

The code and data that support this paper are in the folder
\nolinkurl{kagey-problems/problem131/paperC2} of
\url{https://github.com/apemm/Kagey-Problems}. The command
\texttt{py -I verify\_all.py}, run there with Python~3, \texttt{numpy}, \texttt{scipy},
\texttt{sympy} and \texttt{mpmath}, reruns the checks of Section~\ref{sec:verify}. Each
output goes to \nolinkurl{paperC2/data}, headed by the sha256 of its script.

\section*{Funding}

This research did not receive any specific grant from funding agencies in the
public, commercial, or not-for-profit sectors.

\section*{Declaration of competing interest}

The author declares that he has no known competing financial interests or
personal relationships that could have appeared to influence the work reported
in this paper.

\section*{Acknowledgments}

This work grew out of Problem~131 of Peter Kagey's Open Problem
Collection~\cite{kagey131}. We thank him for posing it.

\begin{thebibliography}{99}

\bibitem{alahmadieh2023}
A.~Al Ahmadieh, The fiber of the principal minor map, preprint, arXiv:2309.00806v2, 2023.
\url{https://arxiv.org/abs/2309.00806v2}.

\bibitem{bermanplemmons1994}
A.~Berman, R.~J.~Plemmons, Nonnegative Matrices in the Mathematical Sciences, SIAM,
Philadelphia, 1994. \url{https://doi.org/10.1137/1.9781611971262}.

\bibitem{bfg2014}
L.~Bertini, A.~Faggionato, D.~Gabrielli, From level 2.5 to level 2 large deviations for
continuous time Markov chains, Markov Process. Related Fields 20 (2014) 545--562.
\url{https://arxiv.org/abs/1212.6908}.

\bibitem{bfg2015}
L.~Bertini, A.~Faggionato, D.~Gabrielli, Large deviations of the empirical flow for continuous
time Markov chains, Ann. Inst. Henri Poincar\'e Probab. Stat. 51 (2015) 867--900.
\url{https://doi.org/10.1214/14-AIHP601}.

\bibitem{bgl2024}
L.~Bertini, D.~Gabrielli, C.~Landim, Metastable $\Gamma$-expansion of finite state Markov chains
level two large deviations rate functions, Ann. Appl. Probab. 34 (2024) 3820--3869.
\url{https://doi.org/10.1214/24-AAP2051}.

\bibitem{catoni1999}
O.~Catoni, Simulated annealing algorithms and Markov chains with rare transitions, in:
S\'eminaire de Probabilit\'es XXXIII, Lecture Notes in Math. 1709, Springer, Berlin, 1999,
pp.~69--119. \url{https://www.numdam.org/item/SPS_1999__33__69_0/}.

\bibitem{cggr2024}
A.~Chatterjee, S.~Ghosh, R.~Gurjar, R.~Raj, Characterizing and testing principal minor
equivalence of matrices, preprint, arXiv:2410.01961v1, 2024.
\url{https://arxiv.org/abs/2410.01961v1}.

\bibitem{cirillonardisohier2015}
E.~N.~M.~Cirillo, F.~R.~Nardi, J.~Sohier, Metastability for general dynamics with rare
transitions: escape time and critical configurations, J. Stat. Phys. 161 (2015) 365--403.
\url{https://doi.org/10.1007/s10955-015-1334-6}.

\bibitem{cohen1981}
J.~E.~Cohen, Convexity of the dominant eigenvalue of an essentially nonnegative matrix, Proc.
Amer. Math. Soc. 81 (1981) 657--658. \url{https://doi.org/10.1090/S0002-9939-1981-0601750-2}.

\bibitem{ddz2024}
J.-F.~Delmas, D.~Dronnier, P.-A.~Zitt, Transformations preserving the effective spectral radius
of a matrix, preprint, arXiv:2406.12857v1, 2024. \url{https://arxiv.org/abs/2406.12857v1}.

\bibitem{dv1975}
M.~D.~Donsker, S.~R.~S.~Varadhan, On a variational formula for the principal eigenvalue for
operators with maximum principle, Proc. Natl. Acad. Sci. USA 72 (1975) 780--783.
\url{https://doi.org/10.1073/pnas.72.3.780}.

\bibitem{engelschneider1980}
G.~M.~Engel, H.~Schneider, Matrices diagonally similar to a symmetric matrix, Linear Algebra
Appl. 29 (1980) 131--138. \url{https://doi.org/10.1016/0024-3795(80)90234-7}.

\bibitem{freidlinwentzell2012}
M.~I.~Freidlin, A.~D.~Wentzell, Random Perturbations of Dynamical Systems, third ed.,
Grundlehren Math. Wiss. 260, Springer, Heidelberg, 2012.
\url{https://doi.org/10.1007/978-3-642-25847-3}.

\bibitem{friedland1981}
S.~Friedland, Convex spectral functions, Linear Multilinear Algebra 9 (1981) 299--316.
\url{https://doi.org/10.1080/03081088108817381}.

\bibitem{good1961}
I.~J.~Good, The frequency count of a Markov chain and the transition to continuous time, Ann.
Math. Statist. 32 (1961) 41--48. \url{https://doi.org/10.1214/aoms/1177705138}.

\bibitem{hartfielloewy1984}
D.~J.~Hartfiel, R.~Loewy, On matrices having equal corresponding principal minors, Linear
Algebra Appl. 58 (1984) 147--167. \url{https://doi.org/10.1016/0024-3795(84)90209-X}.

\bibitem{kagey131}
P.~Kagey, Problem 131, Open Problem Collection. \url{https://peterkagey.com/problems/131/},
2026 (accessed October 2026).

\bibitem{kl2024}
S.~Kim, C.~Landim, $\Gamma$-expansion of the measure-current large deviations rate functional of
non-reversible finite-state Markov chains, Stochastic Process. Appl. 189 (2025) 104733.
\url{https://doi.org/10.1016/j.spa.2025.104733}. Section~5 is cited with the numbering of
arXiv:2412.13515v2.

\bibitem{landim2023}
C.~Landim, Metastability from the large deviations point of view: a $\Gamma$-expansion of the
level two large deviations rate functional of non-reversible finite-state Markov chains,
Stochastic Process. Appl. 165 (2023) 275--315. \url{https://doi.org/10.1016/j.spa.2023.09.001}.

\bibitem{loewy1986}
R.~Loewy, Principal minors and diagonal similarity of matrices, Linear Algebra Appl. 78 (1986)
23--64. \url{https://doi.org/10.1016/0024-3795(86)90015-7}.

\bibitem{olivierivares2005}
E.~Olivieri, M.~E.~Vares, Large Deviations and Metastability, Encyclopedia Math. Appl. 100,
Cambridge University Press, Cambridge, 2005. \url{https://doi.org/10.1017/CBO9780511543272}.

\bibitem{paperC1}
A.~Pemmasani, Occupation counts of rarely switching Markov chains: stratified local laws,
crossings and the boundary layer, preprint, arXiv:2610.09258, 2026.
\url{https://arxiv.org/abs/2610.09258}.

\bibitem{trouve1996}
A.~Trouv\'e, Cycle decompositions and simulated annealing, SIAM J. Control Optim. 34 (1996)
966--986. \url{https://doi.org/10.1137/S0363012993258586}.

\bibitem{wentzell1972}
A.~D.~Wentzell [Ventcel'], On the asymptotics of eigenvalues of matrices with elements of order
$\exp\{-V_{ij}/(2\varepsilon^2)\}$, Dokl. Akad. Nauk SSSR 202 (1972) 263--265 (in Russian).

\bibitem{wolferwatanabe2021}
G.~Wolfer, S.~Watanabe, Information geometry of reversible Markov chains, Inf. Geom. 4 (2021)
393--433. \url{https://doi.org/10.1007/s41884-021-00061-6}.

\end{thebibliography}


\end{document}
